\documentclass[12pt,letterpaper]{article}

\usepackage[utf8]{inputenc}
\usepackage[T1]{fontenc}
\usepackage[spanish,mexico]{babel}
\usepackage[letterpaper,margin=2.5cm,top=3cm]{geometry}
\usepackage{setspace}
\onehalfspacing
\usepackage{amsmath,amssymb,amsfonts}
\usepackage{graphicx}
\usepackage{textcomp}
\usepackage{booktabs}
\usepackage{array}
\usepackage{siunitx}
\usepackage{listings}
\usepackage{float}
\usepackage{caption}
\usepackage{subcaption}
\usepackage{longtable}
\usepackage{enumitem}
\usepackage{xcolor}
\usepackage[colorlinks=true,linkcolor=black,citecolor=blue,urlcolor=blue]{hyperref}

% Ajuste del ancho reservado para numeración de subsecciones en el TOC.
\makeatletter
\renewcommand{\l@subsection}{\@dottedtocline{2}{1.5em}{3.5em}}
\makeatother


\captionsetup{justification=centering}
\graphicspath{{figuras/}{figures/}}

\lstset{
  basicstyle=\ttfamily\small,
  breaklines=true,
  frame=single,
  columns=flexible,
  keepspaces=true,
  showstringspaces=false
}

% Resultados generados automáticamente por Python.
% Este archivo no debe editarse manualmente.
\IfFileExists{datos/resultados_pmsm.tex}{\input{datos/resultados_pmsm}}{}


% Metadatos

\newcommand{\DocumentTitle}{Modelado, simulación y control FOC de un PMSM mediante una arquitectura híbrida}

\newcommand{\DocumentType}{Engineering Technical Report}

\newcommand{\DocumentID}{PMSM-HC-001}

\newcommand{\DocumentVersion}{0.2}

\newcommand{\DocumentStatus}{Technical Development and Architecture Definition}

\newcommand{\DocumentAuthor}{Ulises Ricardo Jasso Castillo}

\newcommand{\DocumentDate}{Septiembre 2026}

\begin{document}

\begin{titlepage}

\thispagestyle{empty}

\centering

\vspace*{1.0cm}

\includegraphics[width=2.5cm]{portada/ipn-eps-converted-to.pdf}

\vspace{0.5cm}

{\Large\textbf{INSTITUTO POLITÉCNICO NACIONAL}\par}

\vspace{0.5cm}

\includegraphics[width=3.5cm]{portada/Esime.png}

\vspace{0.5cm}

{\large\textbf{ESCUELA SUPERIOR DE INGENIERÍA MECÁNICA Y ELÉCTRICA}\par}

\vspace{1.5cm}

{\Large\textbf{\DocumentTitle}\par}

\vspace{1.5cm}

\begin{tabular}{>{\bfseries}p{0.30\textwidth}p{0.55\textwidth}}

Tipo de documento: & \DocumentType \\

Identificador: & \DocumentID \\

Versión: & \DocumentVersion \\

Estado: & \DocumentStatus \\

Autor: & \DocumentAuthor \\

Fecha: & \DocumentDate \\

\end{tabular}

\vfill

{\large Ciudad de México\par}

{\large \DocumentDate\par}

\end{titlepage}

\section{Introducción}

El motor síncrono de imanes permanentes (PMSM) constituye una planta electromecánica en la que las dinámicas eléctricas, electromagnéticas y mecánicas interactúan a diferentes escalas de tiempo. Esta característica hace necesario separar con claridad el modelo de la planta, las transformaciones de variables, el controlador y las capas de supervisión que posteriormente puedan aportar capacidad de adaptación o predicción.

El objetivo del proyecto es desarrollar una arquitectura reproducible de control para un PMSM en la que el FOC constituya el controlador principal y una capa inteligente permita estudiar y tratar las no linealidades del sistema. Las líneas de desarrollo consideradas son ANFIS, MPC y estrategias Tiny/TinyML, cuya integración todavía se encuentra en desarrollo.

El modelo histórico \texttt{skeleton.slx} se utiliza como referencia temporal de arquitectura y trazabilidad. Su análisis sirve para identificar módulos, señales y relaciones funcionales, pero no se adopta como la especificación permanente del sistema final. La especificación del sistema se consolida progresivamente en el desarrollo Python y en sus contratos de frontera.

\section{Objetivos y alcance}

\subsection{Objetivo general}
Desarrollar y documentar una arquitectura reproducible de control para un PMSM basada en FOC como controlador principal y preparada para incorporar ANFIS, MPC y Tiny/TinyML como capa inteligente de predicción, adaptación y optimización.

\subsection{Objetivos específicos}
\begin{enumerate}[label=\arabic*.]
  \item Formalizar el modelo de la planta PMSM y separar sus parámetros de los parámetros del controlador.
  \item Documentar las transformaciones Clarke y Park y las convenciones de ángulo empleadas.
  \item Consolidar la reconstrucción funcional de la cadena M3--M4--M6 en Python.
  \item Documentar detalladamente el algoritmo FOC y su papel como controlador principal.
  \item Establecer contratos de frontera entre las capas de planta, control e inteligencia.
  \item Desarrollar y validar primero en Python las estrategias inteligentes.
  \item Transferir posteriormente a Simulink/Simscape las soluciones ya verificadas para su implementación e integración.
\end{enumerate}


% Capítulo 02 — Máquina de referencia: Siemens 1FT7087
\section{Máquina de referencia: Siemens 1FT7087}

\subsection{Origen físico de la planta}

La planta PMSM desarrollada en este trabajo tiene como referencia física una máquina síncrona de imanes permanentes de la familia Siemens SIMOTICS S-1FT7. La referencia utilizada como base del proyecto es el motor identificado como \textit{SIMOTICS S-1FT7087-7WH70-4CH0}. Esta identificación constituye el punto de partida de la trazabilidad física de la planta y permite separar desde el comienzo los datos propios de la máquina de las condiciones de accionamiento, los parámetros del controlador y las decisiones posteriores de modelado.

La máquina física no se incorpora al proyecto únicamente como un conjunto de números para alimentar una simulación. Su documentación constituye la fuente primaria para establecer magnitudes eléctricas, mecánicas, térmicas y de operación que posteriormente se transforman en una parametrización concentrada y reproducible. A partir de esta parametrización se formula el modelo matemático de la PMSM y, posteriormente, se plantea su evolución hacia representaciones electromagnéticas de mayor fidelidad.

La genealogía adoptada en el proyecto puede resumirse como:

\begin{equation}
\begin{aligned}
\text{máquina física}
&\rightarrow \text{documentación}
\rightarrow \text{caracterización}
\rightarrow \text{parametrización}
\\
&\rightarrow \text{modelo matemático}
\rightarrow \text{planta digital}
\end{aligned}
\end{equation}

Esta secuencia evita interpretar directamente una implementación computacional como si fuera la máquina física. Cada nivel conserva su propia procedencia y sus propias hipótesis.

\subsection{Identificación de la referencia Siemens}

La referencia documental principal utilizada para la caracterización es el manual de configuración de motores síncronos S-1FT7 de Siemens, \textit{S-1FT7 synchronous motors --- Configuration Manual}, documento A5E45099423B AB, edición 03/2024. Este documento proporciona tablas de características, condiciones de operación, parámetros nominales y curvas de funcionamiento de la familia de motores.

La identificación del proyecto se mantiene como \textit{1FT7087-7WH70-4CH0}. En la documentación consultada aparecen, además, denominaciones de familia que no son idénticas en todos los elementos gráficos y tabulares. En particular, algunas curvas de desempeño están rotuladas como \textit{1FT7087-7WH71}, mientras que una tabla numérica utiliza la denominación \textit{1FT7087-7WH7}. Por esta razón, el proyecto conserva la identificación 7WH70-4CH0 como referencia principal y registra estas diferencias documentales sin asumir de manera automática que todas las variantes son idénticas.

Esta precaución es relevante para la trazabilidad. Un valor publicado para una variante de la familia puede servir como referencia documental, pero antes de utilizarlo como evidencia específica de la unidad 7WH70-4CH0 debe verificarse su correspondencia.

\subsection{Datos nominales de la máquina}

La documentación Siemens proporciona las principales magnitudes nominales y límites de operación de la máquina de referencia. Entre ellas se encuentran una velocidad nominal de 4500 rpm, un par nominal de 46 N m y una corriente nominal de 53 A. También se documentan un par estático de 61 N m para la condición S1(100K), una corriente de rotor bloqueado de 67 A, una velocidad mecánica máxima de 8000 rpm, un par máximo de 150 N m y una corriente máxima de 225 A.

Los datos anteriores pertenecen a la caracterización de la máquina y no deben confundirse con los límites utilizados posteriormente por el controlador reconstruido. Por ejemplo, el límite de corriente utilizado en una reconstrucción histórica del FOC constituye una condición de control y no sustituye al valor de corriente máxima publicado para el motor.

Los principales datos documentales utilizados como referencia son:

\begin{itemize}
    \item velocidad nominal: $4500\,\mathrm{rpm}$;
    \item par nominal: $46\,\mathrm{N\,m}$;
    \item corriente nominal: $53\,\mathrm{A}$;
    \item par estático S1(100K): $61\,\mathrm{N\,m}$;
    \item corriente de rotor bloqueado: $67\,\mathrm{A}$;
    \item velocidad mecánica máxima: $8000\,\mathrm{rpm}$;
    \item par máximo: $150\,\mathrm{N\,m}$;
    \item corriente máxima documentada: $225\,\mathrm{A}$;
    \item número de polos: $2p=8$.
\end{itemize}

La documentación también proporciona condiciones de convertidor asociadas a los datos de la máquina. Estas condiciones se mantienen separadas de los parámetros intrínsecos del motor. En particular, la documentación relaciona los datos nominales con un enlace de continua de referencia de $600\,\mathrm{V}$.

\subsection{Número de polos y pares de polos}

La máquina está documentada con ocho polos físicos. Para la formulación eléctrica en coordenadas $d$--$q$, este dato se expresa mediante el número de pares de polos:

\begin{equation}
2p=8
\end{equation}

por lo que:

\begin{equation}
p=4.
\end{equation}

Este valor de $p=4$ es el utilizado en la relación entre magnitudes mecánicas y eléctricas:

\begin{equation}
\omega_e=p\omega_m
\end{equation}

\begin{equation}
\theta_e=p\theta_m.
\end{equation}

\subsection{Parámetros eléctricos documentales}

La documentación Siemens proporciona una resistencia de devanado de $0.046\,\Omega$ a $20\,^{\circ}\mathrm{C}$ y una inductancia identificada como $L_D=1.2\,\mathrm{mH}$. También proporciona una constante de par y una constante de tensión que permiten construir referencias independientes para la parametrización electromagnética.

Los valores documentales principales son:

\begin{itemize}
    \item resistencia de devanado a $20\,^{\circ}\mathrm{C}$: $R=0.046\,\Omega$;
    \item inductancia documentada: $L_D=1.2\,\mathrm{mH}$;
    \item constante de par: $k_T=0.91\,\mathrm{N\,m/A}$;
    \item constante de tensión: $k_E=58\,\mathrm{V/(1000\,rpm)}$.
\end{itemize}

La presencia de un único valor de inductancia documental no permite afirmar directamente que la máquina física tenga $L_d=L_q$ bajo todas las condiciones de operación. La igualdad $L_d=L_q=1.2\,\mathrm{mH}$ adoptada en el nivel M4 corresponde a una hipótesis del modelo concentrado actual. La eventual anisotropía, saturación y dependencia angular pertenecen a niveles de fidelidad posteriores.

\subsection{Constante de par y flujo magnético equivalente}

La constante de par publicada permite construir una estimación del enlace de flujo equivalente empleado en el modelo $d$--$q$:

\begin{equation}
k_T=\frac{3}{2}\sqrt{2}\,p\lambda_m.
\end{equation}

Despejando:

\begin{equation}
\lambda_m=
\frac{k_T}{\left(\frac{3}{2}\right)\sqrt{2}\,p}.
\end{equation}

Con $k_T=0.91\,\mathrm{N\,m/A}$ y $p=4$, se obtiene aproximadamente:

\begin{equation}
\lambda_m\approx0.107245\,\mathrm{Wb}.
\end{equation}

Este valor es el utilizado como referencia numérica en la parametrización M4 actual. Es un valor derivado a partir de una constante publicada y de la convención matemática adoptada, no un valor publicado directamente por Siemens con esa misma magnitud y convención.

\subsection{Constante de tensión como referencia independiente}

La documentación también proporciona $k_E=58\,\mathrm{V/(1000\,rpm)}$. Esta magnitud constituye una segunda referencia para evaluar el orden de magnitud del enlace de flujo, pero debe conservar sus condiciones de definición.

La constante de par está asociada a una condición térmica indicada por el fabricante, mientras que la constante de tensión se especifica a $20\,^{\circ}\mathrm{C}$. Por ello, una diferencia entre una estimación basada en $k_T$ y otra basada en $k_E$ no se interpreta automáticamente como un error de la planta.

\subsection{Parámetros mecánicos}

La documentación proporciona la inercia del rotor según la configuración mecánica. Para la máquina sin freno se documenta:

\begin{equation}
J_{\mathrm{sin\ freno}}=0.00274\,\mathrm{kg\,m^2},
\end{equation}

mientras que para la configuración con freno:

\begin{equation}
J_{\mathrm{con\ freno}}=0.00416\,\mathrm{kg\,m^2}.
\end{equation}

La parametrización M4 utiliza actualmente:

\begin{equation}
J=0.00416\,\mathrm{kg\,m^2}.
\end{equation}

Esta elección se conserva como una decisión trazable de configuración. La dinámica mecánica M4 no incorpora todavía elasticidad torsional del eje ni un modelo detallado de transmisión.

\subsection{Límites de operación y condiciones del convertidor}

La máquina tiene una velocidad mecánica máxima documentada de $8000\,\mathrm{rpm}$, un par máximo de $150\,\mathrm{N\,m}$ y una corriente máxima documentada de $225\,\mathrm{A}$.

Estos valores describen características de la máquina según las condiciones especificadas por el fabricante. No constituyen automáticamente los límites internos de una simulación de control.

En particular, el proyecto utiliza una condición de enlace de continua de $540\,\mathrm{V}$ en el experimento B0. Este valor pertenece al dominio eléctrico del accionamiento utilizado en el experimento y no es un parámetro intrínseco del motor.

\begin{equation}
\text{parámetros de planta}
\neq
\text{condiciones del accionamiento}
\neq
\text{parámetros del controlador}.
\end{equation}

\subsection{Condiciones térmicas y de servicio}

La documentación distingue condiciones de operación continua S1 asociadas a incrementos de temperatura de $100\,\mathrm{K}$ y $60\,\mathrm{K}$, además de condiciones de servicio intermitente.

La constante $k_E$ está documentada a $20\,^{\circ}\mathrm{C}$, mientras que $k_T$ se especifica bajo una condición térmica diferente. La resistencia también se proporciona a $20\,^{\circ}\mathrm{C}$. Estas condiciones deben conservarse al comparar magnitudes.

La planta M4 almacena actualmente algunos coeficientes térmicos como información de trazabilidad, pero no incorpora todavía una dinámica térmica completa.

\subsection{Datos documentales frente a parámetros M4}

La caracterización distingue:

\begin{itemize}
    \item datos publicados directamente por el fabricante;
    \item magnitudes derivadas mediante ecuaciones y convenciones;
    \item decisiones de modelado adoptadas para M4;
    \item valores históricos o recuperados cuya procedencia aún debe verificarse.
\end{itemize}

En el modelo M4 actual se utilizan:

\begin{itemize}
    \item $R_s=0.05\,\Omega$;
    \item $L_d=L_q=1.2\,\mathrm{mH}$;
    \item $\lambda_m=0.107245\,\mathrm{Wb}$;
    \item $p=4$;
    \item $J=0.00416\,\mathrm{kg\,m^2}$.
\end{itemize}

La resistencia M4 de $0.05\,\Omega$ no sustituye al valor documental de $0.046\,\Omega$ a $20\,^{\circ}\mathrm{C}$. La diferencia queda registrada.

De igual manera, $L_d=L_q$ constituye una hipótesis del nivel M4 y no una demostración de que la máquina física carezca de anisotropía.

\subsection{Parámetros pendientes}

Los parámetros de fricción viscosa $B$ y par de fricción $T_{\mathrm{fric}}$ no se incorporan como valores M4 confirmados. Los valores históricos asociados a reconstrucciones anteriores se mantienen separados y pendientes de verificación.

También permanece pendiente la caracterización detallada de saturación, dependencia angular, armónicos, pérdidas y dinámica térmica.

\subsection{Separación entre planta, accionamiento y control}

La máquina de referencia define la planta física. El convertidor, el enlace de continua, la modulación y los límites del controlador forman capas posteriores.

Esta separación es necesaria para que el mismo modelo de planta pueda utilizarse posteriormente con FOC, MPC, ANFIS y TinyML sin modificar silenciosamente sus parámetros físicos.

\begin{equation}
\text{máquina PMSM}
\rightarrow
\text{planta}
\rightarrow
\text{interfaz de accionamiento}
\rightarrow
\text{controlador}.
\end{equation}

Las interfaces entre estos niveles deberán conservar contratos explícitos.

\subsection{Curvas de funcionamiento como evidencia futura}

La documentación Siemens incluye curvas par--velocidad de la familia 1FT7087. Estas curvas constituyen una fuente importante de evidencia para etapas posteriores de validación porque permiten comparar regiones de funcionamiento, límites de par, velocidad y condiciones del convertidor.

Las curvas distinguen condiciones de alimentación y enlace de continua. Por ello, una futura digitalización deberá conservar conjuntamente los ejes, la condición eléctrica, la configuración del convertidor y la identificación de la curva.

La digitalización deberá registrar como mínimo:

\begin{itemize}
    \item fuente documental;
    \item página de origen;
    \item condición de alimentación;
    \item condición del enlace de continua;
    \item variante indicada por la figura;
    \item sistema de unidades;
    \item método de digitalización;
    \item versión del conjunto de datos.
\end{itemize}

No se presentan aquí puntos digitalizados como resultados validados del modelo M4.

\subsection{Discrepancias de variante documental}

La trazabilidad conserva una diferencia importante entre la referencia física del proyecto y algunas denominaciones presentes en el manual. El proyecto utiliza como referencia principal:

\begin{equation}
\text{SIMOTICS S-1FT7087-7WH70-4CH0}.
\end{equation}

Sin embargo, determinadas curvas aparecen identificadas como 1FT7087-7WH71 y algunas tablas como 1FT7087-7WH7.

Esta diferencia no se resuelve suponiendo equivalencia. Cualquier utilización futura de esas curvas como evidencia cuantitativa deberá verificar primero la correspondencia de variante.

\subsection{Evolución prevista de la representación física}

La referencia Siemens constituye el origen físico de una cadena de fidelidad:

\begin{equation}
\mathrm{M4}
\rightarrow
\mathrm{mapas\ electromagn\acute{e}ticos}
\rightarrow
\mathrm{modelo\ no\ lineal}
\rightarrow
\mathrm{FEM/ROM}
\rightarrow
\mathrm{Digital\ Twin}.
\end{equation}

El nivel M4 proporciona actualmente una planta concentrada, ejecutable y reproducible. Los niveles posteriores incorporarán información electromagnética y fenómenos físicos adicionales únicamente cuando exista información trazable para hacerlo.

El aumento de fidelidad no consiste simplemente en agregar complejidad al código, sino en incorporar información física cuya procedencia y comportamiento puedan verificarse.

\subsection{Criterio para cerrar la caracterización}

La caracterización resulta suficientemente estable para el nivel M4 cuando los parámetros activos tienen procedencia identificable, las magnitudes derivadas conservan sus ecuaciones de cálculo, las decisiones de modelado están declaradas y los datos pendientes no se presentan como hechos confirmados.

A partir de esta base puede formalizarse la planta matemática sin confundir el nivel actual de conocimiento con una caracterización completa de todos los fenómenos físicos.

\subsection{Cierre del capítulo}

La máquina Siemens 1FT7087 constituye la referencia física sobre la cual se construye la planta PMSM del proyecto. La documentación del fabricante proporciona las magnitudes necesarias para establecer una primera parametrización y, al mismo tiempo, delimita las condiciones bajo las cuales dichas magnitudes deben interpretarse.

El valor de esta etapa no reside solamente en disponer de parámetros numéricos, sino en establecer una genealogía reproducible entre la máquina física, la documentación, las magnitudes derivadas y el modelo computacional. La parametrización M4 representa actualmente este primer nivel de fidelidad.

A partir de esta base, el capítulo siguiente formaliza la caracterización y parametrización de la máquina, distinguiendo los datos confirmados, las magnitudes derivadas, las hipótesis del modelo y los elementos que permanecen pendientes de verificación.


\section{Caracterización y parametrización de la máquina}

\subsection{Criterio de caracterización}

La caracterización de la máquina tiene como finalidad transformar la información física y documental identificada en el capítulo anterior en un conjunto de parámetros utilizables por la planta digital. Esta transformación no consiste únicamente en copiar valores desde la documentación del fabricante. Cada magnitud debe conservar su origen, condición de operación, unidad, procedimiento de obtención y estado de verificación.

Para mantener la trazabilidad durante la construcción del modelo se adopta una clasificación de cuatro estados:

\begin{itemize}
    \item \textbf{CONFIRMED}: valor respaldado directamente por documentación de referencia o por una derivación considerada cerrada a partir de datos documentales.
    \item \textbf{PROPOSED}: valor o representación adoptada para una etapa del modelo, cuya utilización es explícitamente una decisión de modelado.
    \item \textbf{PENDING\_VERIFICATION}: información recuperada o derivada cuya procedencia, interpretación o equivalencia física todavía requiere comprobación.
    \item \textbf{HISTORICAL}: valor conservado únicamente como evidencia de una implementación anterior y que no constituye la fuente de verdad de la parametrización actual.
\end{itemize}

Esta clasificación evita que un valor procedente de una reconstrucción anterior adquiera automáticamente el carácter de parámetro físico confirmado.

\subsection{Parámetros electromecánicos de referencia}

La siguiente tabla concentra los principales datos de la máquina utilizados como punto de partida de la planta digital.

\begin{table}[H]
\centering
\caption{Datos físicos y electromecánicos de referencia.}
\label{tab:parametros_documentales}
\begin{tabular}{llll}
\toprule
Parámetro & Valor & Unidad & Estado \\
\midrule
Velocidad nominal & 4500 & rpm & CONFIRMED \\
Par nominal & 46 & Nm & CONFIRMED \\
Corriente nominal & 53 & A & CONFIRMED \\
Par estático S1(100K) & 61 & Nm & CONFIRMED \\
Corriente de rotor bloqueado & 67 & A & CONFIRMED \\
Velocidad máxima & 8000 & rpm & CONFIRMED \\
Par máximo & 150 & Nm & CONFIRMED \\
Corriente máxima & 225 & A & CONFIRMED \\
Número de polos $2p$ & 8 & -- & CONFIRMED \\
Pares de polos $p$ & 4 & -- & CONFIRMED* \\
Constante de par $k_T$ & 0.91 & Nm/A & CONFIRMED \\
Constante de tensión $k_E$ & 58 & V/(1000 rpm) & CONFIRMED \\
Resistencia a 20~$^\circ$C & 0.046 & $\Omega$ & CONFIRMED \\
Inductancia de campo giratorio $L_D$ & 1.2 & mH & CONFIRMED \\
Constante de tiempo térmica & 11 & min & CONFIRMED \\
Inercia con freno & 0.00416 & kg\,m$^2$ & CONFIRMED \\
\bottomrule
\end{tabular}

\smallskip
\footnotesize
*Valor derivado directamente de $2p=8$ mediante $p=4$.
\end{table}

Los valores de esta tabla constituyen la referencia documental de la máquina. No todos ellos se incorporan directamente como parámetros de estado del modelo M4. Algunos corresponden a límites de operación, condiciones térmicas o características del accionamiento y serán utilizados posteriormente para validación y comparación.

\subsection{Separación entre datos de máquina y condiciones del accionamiento}

Una consideración fundamental consiste en distinguir los parámetros propios de la máquina de las condiciones impuestas por el accionamiento.

La tensión del enlace de corriente continua, por ejemplo, determina la capacidad eléctrica disponible para generar las tensiones de fase y limita la región de operación del accionamiento. Sin embargo, no constituye una propiedad intrínseca del rotor, estator o circuito magnético de la PMSM.

Por esta razón, el valor de \SI{540}{V} utilizado en el experimento B0 se mantiene como una condición de configuración del accionamiento:

\begin{equation}
V_{DC,\mathrm{B0}}=\SI{540}{V}.
\end{equation}

No se incorpora este valor como parámetro permanente dentro de la definición física de la máquina.

La misma separación se aplica a los límites de corriente y tensión utilizados por el controlador. Un límite impuesto por el FOC no sustituye automáticamente a la corriente máxima documentada para la máquina. Los dos valores pueden ser diferentes porque pertenecen a niveles funcionales distintos.

\subsection{Resistencia de los devanados}

La documentación de Siemens proporciona una resistencia de devanado de:

\begin{equation}
R_{20^\circ C}=\SI{0.046}{\ohm}.
\end{equation}

La implementación M4 actual utiliza:

\begin{equation}
R_s=\SI{0.05}{\ohm}.
\end{equation}

Estos valores no deben confundirse ni sobrescribirse entre sí. El primero corresponde al dato documental a \SI{20}{\celsius}; el segundo corresponde a la parametrización actualmente adoptada por el modelo M4.

Por tanto, la diferencia se conserva explícitamente como parte de la trazabilidad:

\begin{equation}
R_{\mathrm{Siemens},20^\circ C}
\neq
R_{s,\mathrm{M4}}.
\end{equation}

La utilización de \SI{0.05}{\ohm} en M4 constituye una decisión de parametrización que deberá contrastarse posteriormente con las condiciones térmicas y con la respuesta eléctrica de la planta.

\subsection{Inductancias}

El fabricante proporciona una inductancia de campo giratorio de:

\begin{equation}
L_D=\SI{1.2}{mH}.
\end{equation}

La parametrización M4 utiliza:

\begin{equation}
L_d=L_q=\SI{1.2}{mH}.
\end{equation}

Esta igualdad representa una hipótesis de modelado para el nivel concentrado actual. No debe interpretarse como una caracterización exhaustiva de la máquina real.

En particular, una máquina PMSM puede presentar dependencia de la inductancia con la corriente, saturación magnética, saliencia y dependencia angular. Dichos efectos requieren información adicional o modelos de mayor fidelidad. Por ello, la representación actual se considera un primer nivel de aproximación concentrada.

La evolución prevista es:

\begin{equation}
L_d,L_q
\longrightarrow
L_d(i_d,i_q,\theta_e),
L_q(i_d,i_q,\theta_e),
\end{equation}

cuando existan mapas electromagnéticos suficientemente sustentados.

\subsection{Número de pares de polos}

La documentación proporciona el número de polos mediante:

\begin{equation}
2p=8.
\end{equation}

De este dato se obtiene directamente:

\begin{equation}
p=4.
\end{equation}

El valor $p=4$ se utiliza posteriormente para relacionar las magnitudes eléctricas y mecánicas:

\begin{equation}
\omega_e=p\omega_m,
\end{equation}

y

\begin{equation}
\theta_e=p\theta_m.
\end{equation}

Esta derivación constituye un ejemplo de parámetro que no necesita ser tratado como una hipótesis independiente: se obtiene directamente de un dato documental.

\subsection{Flujo permanente del imán}

El flujo permanente utilizado actualmente por la planta es:

\begin{equation}
\lambda_m=\SI{0.107245}{Wb}.
\end{equation}

Este valor no aparece como tal en la tabla de parámetros de Siemens utilizada como fuente principal. Se obtiene a partir de la constante de par documentada:

\begin{equation}
k_T=\SI{0.91}{Nm/A},
\end{equation}

y del número de pares de polos $p=4$.

Utilizando la convención de corriente RMS adoptada para la derivación:

\begin{equation}
k_T=
\frac{3}{2}\sqrt{2}\,p\,\lambda_m,
\end{equation}

por lo que:

\begin{equation}
\lambda_m=
\frac{k_T}
{\left(\frac{3}{2}\right)\sqrt{2}\,p}.
\end{equation}

Sustituyendo:

\begin{equation}
\lambda_m=
\frac{0.91}
{\left(\frac{3}{2}\right)\sqrt{2}(4)}
\approx
\SI{0.107245}{Wb}.
\end{equation}

Por tanto, el valor utilizado por M4 es un parámetro \textbf{derivado} a partir de información del fabricante y no un valor copiado directamente de la documentación.

Su valor numérico se encuentra actualmente marcado como confirmado para la parametrización M4, mientras que su procedencia documental específica permanece bajo trazabilidad pendiente de verificación. Esta distinción se conserva para evitar confundir la confirmación de un valor numérico dentro del modelo con la demostración completa de su procedencia experimental.

\subsection{Comparación con la constante de tensión}

La documentación proporciona también:

\begin{equation}
k_E=\SI{58}{V/(1000\,rpm)}
\end{equation}

a \SI{20}{\celsius}.

La conversión de esta magnitud a una constante equivalente en función de la velocidad angular permite obtener una segunda referencia para el flujo permanente. Sin embargo, las condiciones de temperatura y las convenciones de línea/fase y RMS/pico deben mantenerse explícitas durante la conversión.

La existencia de una diferencia entre la estimación derivada desde $k_T$ y la estimación derivada desde $k_E$ no se interpreta automáticamente como un error de uno de los datos. Ambas constantes están asociadas a condiciones de operación diferentes y además poseen tolerancias de fabricación especificadas por el fabricante.

Por ello, la estimación obtenida desde $k_T$ se mantiene como la utilizada actualmente por M4, mientras que la estimación asociada a $k_E$ se conserva como referencia independiente para futuras verificaciones.

\subsection{Inercia mecánica}

La documentación proporciona una inercia de:

\begin{equation}
J=\SI{27.4e-4}{kg\,m^2}
\end{equation}

sin freno y:

\begin{equation}
J=\SI{41.6e-4}{kg\,m^2}
=\SI{0.00416}{kg\,m^2}
\end{equation}

con freno.

La parametrización M4 actual utiliza:

\begin{equation}
J=\SI{0.00416}{kg\,m^2}.
\end{equation}

Por tanto, el modelo actual corresponde a la condición de inercia asociada a la configuración con freno indicada en la documentación.

La elección debe conservarse como parte de la configuración de la planta, ya que la inercia modifica directamente la dinámica mecánica:

\begin{equation}
J\frac{d\omega_m}{dt}
=
T_e-T_L.
\end{equation}

\subsection{Parámetros térmicos conservados}

La parametrización M4 conserva actualmente algunos coeficientes térmicos:

\begin{equation}
\alpha=3.93\times10^{-3}\;K^{-1},
\end{equation}

\begin{equation}
\alpha_{\Phi}=-1.0\times10^{-3}\;K^{-1},
\end{equation}

y una temperatura de referencia:

\begin{equation}
T_{\mathrm{meas}}=\SI{298.15}{K}.
\end{equation}

Estos valores se mantienen almacenados para trazabilidad y desarrollos posteriores. No constituyen todavía una dinámica térmica completa de la máquina.

La documentación Siemens proporciona una constante de tiempo térmica de aproximadamente \SI{11}{min}, pero no se dispone todavía de un modelo térmico completo con capacidad térmica, resistencias térmicas y estados suficientemente establecidos para incorporarlo como dinámica activa de M4.

\subsection{Parámetros mecánicos pendientes}

Los parámetros asociados a fricción viscosa y par de fricción no forman parte de la parametrización M4 actualmente adoptada.

En particular:

\begin{equation}
B
\end{equation}

y

\begin{equation}
T_{\mathrm{fric}}
\end{equation}

permanecen pendientes de verificación.

Valores encontrados en implementaciones históricas, como $B=\SI{0.001}{N\,m\,s/rad}$ y $T_{\mathrm{fric}}=\SI{0.5}{Nm}$, no se incorporan como parámetros confirmados de la máquina.

La dinámica mecánica actual utiliza por tanto el par de carga externo como entrada:

\begin{equation}
J\frac{d\omega_m}{dt}=T_e-T_L.
\end{equation}

Esta decisión permite mantener separado el modelo confirmado de la máquina de una representación de pérdidas mecánicas cuya identificación todavía no se ha cerrado.

\subsection{Parámetros de la planta frente a parámetros del FOC}

La caracterización de la máquina debe mantenerse separada de los parámetros utilizados posteriormente por el controlador.

Por ejemplo, el modelo M4 utiliza:

\begin{equation}
\lambda_m=\SI{0.107245}{Wb},
\end{equation}

mientras que la implementación reconstruida del FOC conserva actualmente un valor histórico de:

\begin{equation}
\lambda_{m,\mathrm{FOC}}=\SI{0.1100}{Wb}.
\end{equation}

De manera similar, la documentación de la máquina proporciona una corriente máxima de \SI{225}{A}, mientras que el FOC reconstruido utiliza actualmente un límite de corriente $i_q$ de aproximadamente \SI{112.43}{A}.

Estas diferencias no se corrigen silenciosamente en esta etapa. Se documentan como diferencias entre capas y configuraciones históricas para evitar perder la trazabilidad de la reconstrucción.

El objetivo de la arquitectura es que la planta mantenga una parametrización propia y que los controladores reciban mediante contratos las magnitudes necesarias para operar sobre ella.

\subsection{Tabla maestra de parametrización M4}

La parametrización actualmente utilizada por la planta concentrada puede resumirse como:

\begin{table}[H]
\centering
\caption{Parámetros principales de la planta PMSM M4.}
\label{tab:parametros_m4}
\begin{tabular}{llll}
\toprule
Parámetro & Valor & Unidad & Estado \\
\midrule
$R_s$ & 0.05 & $\Omega$ & CONFIRMED \\
$M_s$ & $2.0\times10^{-5}$ & H & CONFIRMED \\
$L_0$ & $1.6\times10^{-4}$ & H & CONFIRMED \\
$L_d$ & $1.2\times10^{-3}$ & H & CONFIRMED \\
$L_q$ & $1.2\times10^{-3}$ & H & CONFIRMED \\
$\lambda_m$ & 0.107245 & Wb & CONFIRMED \\
$p$ & 4 & -- & CONFIRMED \\
$J$ & 0.00416 & kg\,m$^2$ & CONFIRMED \\
$B$ & -- & N\,m\,s/rad & PENDING\_VERIFICATION \\
$T_{\mathrm{fric}}$ & -- & Nm & PENDING\_VERIFICATION \\
\bottomrule
\end{tabular}
\end{table}

Los estados indicados en esta tabla describen la situación de la parametrización dentro del proyecto y no deben interpretarse como una certificación experimental de cada magnitud.

\subsection{Niveles de fidelidad de la parametrización}

La parametrización M4 representa el primer nivel de una evolución de fidelidad. El objetivo no es mantener indefinidamente todos los parámetros constantes, sino utilizar esta representación como base reproducible sobre la cual incorporar información adicional.

La evolución prevista puede expresarse como:

\begin{equation}
\mathrm{M4}_{\mathrm{concentrated}}
\rightarrow
\mathrm{M4}_{\mathrm{maps}}
\rightarrow
\mathrm{M4}_{\mathrm{nonlinear}}
\rightarrow
\mathrm{FEM/ROM}
\rightarrow
\mathrm{Digital\ Twin}.
\end{equation}

En el segundo nivel, los parámetros electromagnéticos podrán convertirse en funciones del punto de operación. Por ejemplo:

\begin{equation}
\psi_d=\psi_d(i_d,i_q,\theta_e),
\end{equation}

\begin{equation}
\psi_q=\psi_q(i_d,i_q,\theta_e),
\end{equation}

y:

\begin{equation}
T_e=T_e(i_d,i_q,\theta_e).
\end{equation}

Estas funciones no se consideran todavía parámetros confirmados de la máquina. Su construcción requiere datos electromagnéticos suficientemente trazables, procedentes de análisis, simulación, FEM, identificación experimental o una combinación documentada de estas fuentes.

\subsection{Criterio de cierre de la parametrización}

La parametrización de una etapa se considerará cerrada únicamente cuando cada magnitud empleada por el modelo pueda clasificarse según su procedencia y función.

En consecuencia, el criterio de cierre no consiste solamente en disponer de un archivo de parámetros que permita ejecutar una simulación. Se requiere además conocer:

\begin{enumerate}
    \item de dónde procede cada parámetro;
    \item bajo qué condición fue definido;
    \item si fue publicado, derivado o propuesto;
    \item en qué nivel del modelo se utiliza;
    \item qué incertidumbre o limitación conserva;
    \item y qué evidencia permitiría modificar su estado de verificación.
\end{enumerate}

Con este criterio, la parametrización deja de ser una colección de constantes aisladas y se convierte en una capa explícita de trazabilidad de la planta.

El siguiente capítulo utilizará esta parametrización para describir la implementación de la planta PMSM M4, sus estados, interfaces, entradas, salidas y límites de fidelidad.


\section{Planta PMSM}

\subsection{Función de la planta dentro del Digital Twin}

La planta PMSM constituye el núcleo físico-computacional sobre el cual se desarrolla progresivamente el Digital Twin de la máquina de referencia Siemens 1FT7087. Su función es proporcionar una representación matemática ejecutable de las principales dinámicas eléctricas, electromagnéticas y mecánicas de la máquina, manteniendo trazabilidad respecto de los datos físicos que la parametrizan.

La planta no debe entenderse como un bloque aislado cuya única finalidad sea producir una respuesta para un controlador. Dentro de la arquitectura propuesta, constituye una representación computacional de la máquina física que puede evolucionar en nivel de fidelidad y utilizarse como plataforma común para diferentes estrategias de control, análisis y validación.

El nivel implementado actualmente corresponde a la realización M4, denominada \texttt{M4\_BASE}. Esta realización constituye la primera implementación computacional reproducible de la planta a partir de un conjunto de parámetros físicamente trazables. Sobre esta base se podrá incorporar progresivamente información electromagnética adicional, incluyendo mapas de flujo y torque, dependencias no lineales, información procedente de modelos FEM y, cuando resulte necesario, modelos reducidos.

La linealidad de determinadas relaciones utilizadas en M4 constituye una característica del nivel de representación adoptado y no la definición final de la planta. La arquitectura se ha diseñado desde el inicio para permitir el enriquecimiento progresivo de la representación física.

La separación entre planta y control constituye una decisión fundamental de la arquitectura. La planta representa el comportamiento de la máquina, mientras que los algoritmos de control determinan las acciones que se aplican sobre ella. Por esta razón, el FOC, el MPC y las futuras capas ANFIS o TinyML no forman parte de la dinámica interna de \texttt{PMSMBase}.

Esta separación permite que una misma representación física pueda utilizarse posteriormente como plataforma común para comparar diferentes estrategias de control sin modificar la definición fundamental de la máquina.

El flujo conceptual de la planta dentro de la arquitectura es:

\begin{equation}
\begin{aligned}
\text{Máquina física}
&\rightarrow \text{Parametrización trazable}
\rightarrow \text{Planta PMSM}
\\
&\rightarrow \text{Accionamiento y control}
\rightarrow \text{Validación}.
\end{aligned}
\end{equation}

De esta manera, la planta constituye un elemento estable de la arquitectura, mientras que su representación interna puede aumentar progresivamente en fidelidad.

\subsection{Arquitectura de implementación}

La implementación actual en Python está organizada en una cadena explícita:

\begin{equation}
\texttt{parameters.py}
\rightarrow
\texttt{PMSMParameters}
\rightarrow
\texttt{PMSMBase}
\rightarrow
\texttt{PlantInput/PlantOutput}.
\end{equation}

Cada elemento tiene una responsabilidad específica.

El archivo \texttt{parameters.py} contiene la parametrización utilizada por el modelo M4. No contiene parámetros propios del FOC ni valores históricos procedentes de CloudE que hayan sido descartados de la fuente M4.

La clase \texttt{PMSMParameters} agrupa los valores utilizados directamente por la planta y proporciona una validación básica de coherencia.

La clase \texttt{PMSMBase}, ubicada en \texttt{pmsm\_base.py}, implementa los estados, las ecuaciones dinámicas, las transformaciones necesarias y el paso de integración de la planta.

Finalmente, las estructuras \texttt{PlantInput} y \texttt{PlantOutput} definen la interfaz actual de intercambio de información entre la planta y el resto del sistema.

Esta separación permite distinguir claramente:

\begin{enumerate}
\item la información física utilizada para parametrizar la máquina;
\item el modelo dinámico que representa su comportamiento;
\item las señales mediante las cuales la planta se conecta con otros bloques.
\end{enumerate}

La organización también facilita la evolución posterior de la representación. La incorporación de mapas o de relaciones electromagnéticas de mayor fidelidad puede realizarse dentro de la frontera de la planta, manteniendo, cuando sea posible, contratos compatibles con las capas externas.

\subsection{Contenedor y trazabilidad de parámetros}

La implementación utiliza la estructura \texttt{TraceableValue} para asociar a una magnitud numérica su estado de trazabilidad, fuente, unidad y, cuando corresponde, una nota adicional.

Los estados definidos en el código son:

\begin{itemize}
\item \texttt{CONFIRMED}: valor considerado confirmado dentro de la fuente de datos M4;
\item \texttt{PROPOSED}: decisión interna de diseño o modelado;
\item \texttt{PENDING\_VERIFICATION}: información cuya procedencia o interpretación todavía requiere verificación.
\end{itemize}

La clase \texttt{PMSMParameters} reúne los parámetros empleados directamente por la planta:

\begin{equation}
\Theta =
\left\{
R_s,L_d,L_q,\lambda_m,p,J
\right\}.
\end{equation}

Además, conserva otros valores definidos por la estructura de datos M4, entre ellos $M_s$, $L_0$ y $N$, que deben distinguirse de los parámetros que participan directamente en las ecuaciones dinámicas activas.

Una distinción importante es que $N$ y $p$ no representan la misma magnitud. En la implementación:

\begin{equation}
N=6
\end{equation}

corresponde a un parámetro asociado a la configuración del bloque, mientras que:

\begin{equation}
p=4
\end{equation}

representa el número de pares de polos utilizado por la planta. Esta distinción evita reutilizar de forma incorrecta un parámetro de configuración como si fuera una propiedad electromecánica.

Los parámetros de fricción y amortiguamiento:

\begin{equation}
B
\end{equation}

y:

\begin{equation}
T_{\mathrm{fric}}
\end{equation}

se conservan como referencias pendientes de verificación, pero no participan actualmente en la dinámica M4.

Asimismo, el valor histórico:

\begin{equation}
\lambda_m=0.11,\mathrm{Wb}
\end{equation}

procedente de implementaciones anteriores no sustituye al valor utilizado actualmente por M4:

\begin{equation}
\lambda_m=0.107245,\mathrm{Wb}.
\end{equation}

Esta separación permite que el código actual permanezca ligado a la fuente M4 mientras los valores históricos se conservan únicamente como evidencia de las etapas anteriores de reconstrucción.

\subsection{Estados de la planta}

La planta M4 utiliza cuatro estados internos:

\begin{equation}
x=
\begin{bmatrix}
i_d &
i_q &
\omega_m &
\theta_m
\end{bmatrix}^{T}.
\end{equation}

Los estados eléctricos son las corrientes $i_d$ e $i_q$, expresadas en el sistema de coordenadas rotante. La corriente $i_d$ representa la componente asociada al eje $d$, mientras que $i_q$ representa la componente asociada al eje $q$, que en la parametrización actual participa directamente en la producción del par electromagnético.

Los otros dos estados representan la parte mecánica de la máquina:

\begin{itemize}
\item $\omega_m$: velocidad angular mecánica;
\item $\theta_m$: posición angular mecánica.
\end{itemize}

Los estados iniciales configurados por la implementación son:

\begin{equation}
i_d(0)=0,
\qquad
i_q(0)=0,
\end{equation}

\begin{equation}
\omega_m(0)=0,
\qquad
\theta_m(0)=0.
\end{equation}

A partir de la posición mecánica se obtiene la posición eléctrica mediante:

\begin{equation}
\theta_e=p\theta_m.
\end{equation}

De forma equivalente, la velocidad eléctrica se obtiene mediante:

\begin{equation}
\omega_e=p\omega_m.
\end{equation}

Estas relaciones conectan directamente la dinámica mecánica con la representación eléctrica utilizada por el modelo.

\subsection{Entradas de la planta}

La interfaz \texttt{PlantInput} define actualmente cinco magnitudes:

\begin{equation}
u_{\mathrm{plant}}=
\begin{bmatrix}
v_a &
v_b &
v_c &
T_{\mathrm{load}} &
\theta_m
\end{bmatrix}^{T}.
\end{equation}

Las tensiones $v_a$, $v_b$ y $v_c$ representan las tensiones aplicadas a las tres fases de la máquina.

El término $T_{\mathrm{load}}$ representa el par resistente externo utilizado por la simulación. Debido a que los términos $B$ y $T_{\mathrm{fric}}$ todavía no están verificados, el modelo actual utiliza el par de carga como representación del esfuerzo mecánico externo aplicado a la máquina.

La interfaz también recibe $\theta_m$. Esta magnitud se utiliza actualmente para determinar el ángulo eléctrico requerido por las transformaciones de coordenadas.

Esta estructura corresponde a la interfaz actual de la planta reconstruida. No reproduce directamente los nombres de los buses históricos de CloudE. La arquitectura actual define interfaces propias y trazables que pueden mantenerse independientemente de la organización de los modelos históricos.

\subsection{Salidas de la planta}

La estructura \texttt{PlantOutput} proporciona las principales variables necesarias para observar el comportamiento de la máquina:

\begin{equation}
y_{\mathrm{plant}}=
\begin{bmatrix}
i_a &
i_b &
i_c &
i_d &
i_q &
\omega_m &
\theta_m &
T_e
\end{bmatrix}^{T}.
\end{equation}

Las corrientes $i_a$, $i_b$ e $i_c$ permiten observar el comportamiento trifásico de la máquina.

Las corrientes $i_d$ e $i_q$ permiten observar el mismo comportamiento en el marco rotante utilizado por el modelo y posteriormente por los algoritmos de control.

Las variables $\omega_m$ y $\theta_m$ representan la respuesta mecánica, mientras que $T_e$ representa el par electromagnético calculado por la planta.

Esta combinación permite utilizar las mismas salidas para análisis de la máquina, validación del modelo y realimentación de los controladores.

\subsection{Representación de las variables eléctricas}

Aunque la frontera de la planta utiliza variables trifásicas, el núcleo dinámico se formula en coordenadas $dq$.

El flujo de representación utilizado es:

\begin{equation}
abc
\rightarrow
\alpha\beta
\rightarrow
dq.
\end{equation}

Las tensiones trifásicas recibidas mediante \texttt{PlantInput} son transformadas primero al plano estacionario $\alpha\beta$ y posteriormente al sistema rotante $dq$.

Una vez calculadas las corrientes internas $i_d$ e $i_q$, la transformación inversa permite obtener las corrientes trifásicas:

\begin{equation}
dq
\rightarrow
\alpha\beta
\rightarrow
abc.
\end{equation}

La transformación de coordenadas no representa un cambio físico en la máquina. Se trata de un cambio de representación matemática que permite expresar las variables eléctricas en un sistema de referencia conveniente para describir la interacción entre corriente, flujo, velocidad y par.

En la implementación actual, las transformaciones Clarke y Park se utilizan como parte del paso de ejecución de la planta. Por tanto, la planta puede recibir variables de fase y entregar nuevamente variables de fase, mientras que sus ecuaciones internas trabajan en el marco $dq$.

Esta característica será especialmente importante en los capítulos posteriores dedicados al FOC y a la formulación matemática.

\subsection{Dinámica eléctrica}

El núcleo eléctrico de la planta está definido en coordenadas $dq$.

La evolución de la corriente en el eje $d$ se implementa como:

\begin{equation}
\frac{di_d}{dt} =
\frac{v_d-R_si_d+\omega_eL_qi_q}{L_d}.
\end{equation}

La corriente en el eje $q$ evoluciona según:

\begin{equation}
\frac{di_q}{dt} =
\frac{v_q-R_si_q-\omega_e(L_di_d+\lambda_m)}{L_q}.
\end{equation}

Estas expresiones muestran que las corrientes no dependen únicamente de las tensiones aplicadas. También intervienen la resistencia de los devanados, las inductancias, el flujo permanente y la velocidad eléctrica.

Por tanto, la planta contiene una interacción entre los dominios eléctrico y mecánico: cuando aumenta $\omega_m$, aumenta también $\omega_e$, modificando los términos de acoplamiento presentes en las ecuaciones eléctricas.

La relación:

\begin{equation}
\omega_e=p\omega_m
\end{equation}

establece precisamente este vínculo entre ambos dominios.

La obtención física de estas ecuaciones se desarrolla posteriormente en el capítulo dedicado a la formulación matemática, donde se partirá de las relaciones de tensión, flujo y fuerza electromotriz para justificar cada término.

\subsection{Modelo electromagnético}

El modelo M4 utiliza como expresión de referencia para el par electromagnético:

\begin{equation}
T_e=
\frac{3}{2}p
\left[
\lambda_mi_q+
(L_d-L_q)i_di_q
\right].
\end{equation}

La estructura completa de la ecuación se conserva en la implementación aunque actualmente:

\begin{equation}
L_d=L_q.
\end{equation}

Por esta razón:

\begin{equation}
L_d-L_q=0,
\end{equation}

y el término asociado al par de reluctancia se anula dentro de la parametrización M4 actual.

El par queda entonces reducido a la contribución asociada al flujo permanente:

\begin{equation}
T_e=
\frac{3}{2}p\lambda_mi_q.
\end{equation}

Esta simplificación es consecuencia de la parametrización utilizada en este nivel del modelo y no debe interpretarse como una propiedad universal de todas las máquinas PMSM.

Una PMSM real puede presentar saliencia, saturación, dependencia de los enlaces de flujo con la corriente, dependencia angular, armónicos y otros fenómenos electromagnéticos. Estos efectos podrán incorporarse progresivamente cuando exista información suficiente para aumentar la fidelidad de la planta.

Además, el código identifica explícitamente la expresión del par como una ecuación de referencia o candidata. Por ello, no se afirma en esta etapa que constituya la ley interna exacta del modelo Simscape de referencia.

\subsection{Dinámica mecánica}

La parte mecánica de la planta se describe mediante el balance entre el par electromagnético generado y el par de carga aplicado:

\begin{equation}
J\frac{d\omega_m}{dt}
=
T_e-T_{\mathrm{load}}.
\end{equation}

La velocidad mecánica determina a su vez la evolución de la posición:

\begin{equation}
\frac{d\theta_m}{dt}
=
\omega_m.
\end{equation}

La dinámica mecánica proporciona así el vínculo entre la generación de par y el movimiento del rotor.

En el modelo actual no se introducen de forma independiente los términos de amortiguamiento viscoso ni fricción seca porque los parámetros $B$ y $T_{\mathrm{fric}}$ permanecen pendientes de verificación.

Esto no significa que la máquina física carezca de pérdidas mecánicas. Significa que esas pérdidas todavía no han sido identificadas con suficiente trazabilidad para formar parte de la parametrización M4.

\subsection{Integración numérica}

Una vez calculadas las derivadas de los estados, la implementación actual utiliza integración explícita de Euler.

Para un estado genérico $x$:

\begin{equation}
x_{k+1}
=
x_k+
\dot{x}_k\Delta t.
\end{equation}

El intervalo $\Delta t$ se proporciona a la función \texttt{step()} y debe ser positivo.

El procedimiento de un paso puede resumirse como:

\begin{enumerate}
\item recibir las variables de entrada;
\item calcular el ángulo eléctrico;
\item transformar las variables trifásicas al marco $dq$;
\item calcular las derivadas de las corrientes y variables mecánicas;
\item integrar los estados;
\item transformar las corrientes $dq$ nuevamente a $abc$;
\item construir la estructura \texttt{PlantOutput}.
\end{enumerate}

La elección de Euler explícito está clasificada en el código como una decisión \texttt{PROPOSED}. Su objetivo en esta fase es proporcionar una implementación transparente y reproducible.

El método numérico no forma parte de la física de la máquina. Por tanto, podrá modificarse posteriormente cuando los experimentos requieran diferentes compromisos entre precisión, estabilidad y costo computacional.

\subsection{Nivel de representación M4}

M4 no intenta representar simultáneamente todos los fenómenos presentes en una máquina PMSM real. Su función dentro de la arquitectura es establecer una primera realización computacional físicamente trazada, reproducible y suficientemente estructurada para permitir su posterior enriquecimiento.

Entre los fenómenos que todavía no forman parte de la dinámica activa se encuentran:

\begin{itemize}
\item saturación magnética;
\item dependencia de $L_d$ y $L_q$ con las corrientes;
\item dependencia angular de los enlaces de flujo;
\item armónicos espaciales;
\item pérdidas electromagnéticas detalladas;
\item dinámica térmica completa;
\item identificación de fricción y amortiguamiento;
\item resolución mediante FEM;
\item mapas electromagnéticos físicamente validados.
\end{itemize}

Estos elementos delimitan el nivel de representación actual y, al mismo tiempo, establecen los ejes sobre los cuales puede aumentar la fidelidad de la planta.

\subsection{Mapas y evolución del modelo}

El repositorio contiene estructuras destinadas a soportar una evolución de la representación electromagnética.

El módulo \texttt{flux\_maps.py} proporciona actualmente una representación analítica conservadora de los enlaces de flujo basada en los parámetros constantes de M4. Esta representación se mantiene como referencia propuesta y no como un mapa electromagnético experimental o FEM validado.

De forma equivalente, \texttt{torque\_map.py} conserva una representación analítica del par basada en la ecuación de referencia.

El repositorio también conserva datos recuperados de etapas anteriores. Estos datos no se activan automáticamente porque su semántica, correspondencia de ejes, unidades y procedencia deben verificarse antes de utilizarlos como fuente de un modelo de mayor fidelidad.

El archivo \texttt{pmsm\_nonlinear.py} constituye la estructura prevista para una futura extensión del modelo. Actualmente no implementa una dinámica no lineal activa y, por tanto, no se presenta como el modelo operativo de la máquina.

La evolución de la representación no se interpreta como una secuencia rígida, sino como un proceso de enriquecimiento de la información física:

\begin{equation}
\text{M4}
\rightarrow
\text{información electromagnética verificada}
\rightarrow
\text{modelo de mayor fidelidad}.
\end{equation}

La información FEM podrá constituir una fuente independiente de alta fidelidad para generar o contrastar mapas y modelos reducidos.

Cada transición deberá conservar la trazabilidad respecto de la representación anterior.

\subsection{Frontera entre planta y control}

La planta PMSM se mantiene conceptualmente separada de los algoritmos de control.

El intercambio puede representarse como:

\begin{equation}
\mathrm{Controlador}
\rightarrow
\mathrm{Contrato}
\rightarrow
\mathrm{Planta}
\rightarrow
\mathrm{Contrato}
\rightarrow
\mathrm{Controlador}.
\end{equation}

Los contratos y buses actúan como mecanismos de intercambio de información entre bloques. Esto permite evitar que cada bloque dependa directamente de la implementación interna de otro.

La planta, por tanto, no debe conocer si la acción de control procede de un FOC, de un MPC o de una futura capa supervisora. Su responsabilidad consiste en recibir las variables definidas por su interfaz, evolucionar sus estados y devolver las variables correspondientes.

Esta separación será especialmente importante cuando se incorporen diferentes estrategias de control sobre una misma representación física.

El FOC continúa siendo el controlador de referencia establecido en la arquitectura actual. Las estrategias MPC, ANFIS y TinyML se consideran desarrollos posteriores que deberán operar sobre la misma frontera física y respetar los contratos establecidos.

\subsection{Particularidad de la posición mecánica}

La implementación actual presenta una particularidad que debe conservarse explícitamente en la trazabilidad.

La estructura \texttt{PlantInput} contiene la posición mecánica $\theta_m$, pero $\theta_m$ también forma parte de los estados internos de \texttt{PMSMBase}.

Durante la ejecución de \texttt{step()}, la posición recibida se utiliza para calcular:

\begin{equation}
\theta_e=p\theta_m.
\end{equation}

Posteriormente, el valor recibido se asigna al estado interno antes de realizar la integración.

El estado interno vuelve entonces a evolucionar mediante:

\begin{equation}
\theta_{m,k+1}
=
\theta_{m,k}
+
\omega_{m,k}\Delta t.
\end{equation}

Esta característica representa la implementación actual de la planta reconstruida. No se modifica en esta etapa porque hacerlo implicaría cambiar la semántica de la interfaz y podría afectar los experimentos existentes.

Sin embargo, queda registrada como un punto de consolidación de la arquitectura futura. Una versión posterior deberá establecer de forma inequívoca si la posición del rotor constituye una entrada externa, un estado interno o una variable medida o estimada procedente de otro bloque.

\subsection{Independencia respecto del accionamiento}

La planta M4 tampoco incorpora como parámetro propio la tensión del enlace de corriente continua.

El valor $V_{DC}$ pertenece al dominio eléctrico del accionamiento y determina las tensiones que pueden generarse para alimentar la máquina. Por ello, no forma parte de \texttt{PMSMParameters}.

Esta separación permite estudiar la misma máquina bajo diferentes condiciones del accionamiento sin modificar su identidad física.

Por ejemplo, una simulación puede utilizar:

\begin{equation}
V_{DC}=540,\mathrm{V},
\end{equation}

como condición experimental del sistema de potencia, mientras que los parámetros propios de la máquina permanecen sin cambios.

Esta distinción resulta necesaria posteriormente para comparar regiones de operación, límites de tensión, debilitamiento de campo y curvas de par--velocidad.

\subsection{Nivel de fidelidad M4}

La implementación descrita corresponde a la primera realización computacional reproducible de la planta PMSM.

Su valor no reside en representar desde el inicio todos los fenómenos electromagnéticos de la máquina real, sino en establecer una base física, matemática y computacional trazable sobre la cual pueda aumentarse progresivamente la fidelidad.

El nivel actual puede resumirse como:

\begin{equation}
\boxed{
\text{PMSM concentrada}
+
\text{parámetros trazables}
+
\text{dinámica }dq
+
\text{dinámica mecánica}
}
\end{equation}

A partir de esta base podrán incorporarse progresivamente las dependencias que actualmente no están representadas.

El objetivo final no es simplemente disponer de un modelo matemático más complejo, sino construir una representación digital de la máquina que pueda utilizarse para experimentar, comparar y validar diferentes estrategias de control y diferentes niveles de representación física.

\subsection{Resumen de la planta implementada}

La planta M4 actual puede resumirse mediante la siguiente cadena:

\begin{equation}
\begin{aligned}
\text{Parámetros M4}
&\rightarrow \text{Estados}
\rightarrow \text{Transformaciones}
\\
&\rightarrow \text{Dinámica eléctrica}
\rightarrow \text{Torque}
\rightarrow \text{Dinámica mecánica}.
\end{aligned}
\end{equation}

En términos de software:

\begin{equation}
\texttt{parameters.py}
\rightarrow
\texttt{PMSMBase}
\rightarrow
\texttt{PlantInput}
\rightarrow
\texttt{step()}
\rightarrow
\texttt{PlantOutput}.
\end{equation}

Esta implementación constituye la primera representación computacional reproducible de la planta física de referencia.

A partir de ella, los capítulos siguientes desarrollarán primero el fundamento físico de las ecuaciones utilizadas, posteriormente su formulación en espacio de estados y finalmente su relación con las implementaciones de referencia MATLAB/Simscape, la reconstrucción Python y los procesos de validación.

La planta M4 queda así definida como una base de trabajo físicamente trazada, reproducible y evolutiva, sobre la cual puede aumentar progresivamente el nivel de fidelidad del Digital Twin.


\section{Formulación matemática de la planta PMSM}

\subsection{Propósito de la formulación matemática}

La formulación matemática constituye el vínculo entre la máquina física de referencia y la implementación computacional de la planta PMSM. Su función no es únicamente proporcionar un conjunto de ecuaciones para ejecutar una simulación, sino establecer de manera explícita cómo los fenómenos físicos de la máquina se traducen en variables, relaciones matemáticas, estados y entradas que pueden ser implementados en software.

La formulación utilizada en este trabajo parte de la descripción de una máquina eléctrica trifásica con imanes permanentes y conduce progresivamente hacia una representación en coordenadas rotantes $dq$. Esta representación permite separar de manera conveniente las componentes asociadas a los ejes magnéticos de la máquina y construir un modelo compacto de sus dinámicas eléctrica, electromagnética y mecánica.

El proceso puede resumirse conceptualmente como:

\begin{equation}
\begin{aligned}
\text{fenómeno físico}
&\rightarrow \text{ley física}
\rightarrow \text{variable matemática}
\\
&\rightarrow \text{ecuación}
\rightarrow \text{modelo de estados}
\rightarrow \text{implementación}.
\end{aligned}
\end{equation}

La formulación actual corresponde al nivel concentrado M4. Por tanto, sus parámetros electromagnéticos se consideran constantes dentro de este nivel y no se introducen todavía dependencias explícitas de saturación, posición angular, temperatura, armónicos espaciales o mapas obtenidos mediante FEM.

Esta decisión no pretende representar de manera definitiva todos los fenómenos de la máquina real. Su objetivo es establecer una base matemática coherente y trazable que pueda evolucionar posteriormente hacia representaciones de mayor fidelidad.

\subsection{Origen físico del modelo}

La planta se origina en una máquina PMSM trifásica de referencia Siemens 1FT7087. La información documental de la máquina proporciona magnitudes como resistencia de devanado, inductancia, constante de par, constante de tensión, velocidad, corriente, número de polos e inercia.

Sin embargo, los datos de catálogo no constituyen por sí mismos un modelo dinámico completo. Para obtener una planta ejecutable es necesario establecer las relaciones que gobiernan la evolución de las variables eléctricas y mecánicas.

La formulación se construye a partir de tres dominios físicos principales:

\begin{enumerate}
    \item dominio eléctrico;
    \item dominio electromagnético;
    \item dominio mecánico.
\end{enumerate}

Estos dominios están acoplados.

Las tensiones aplicadas producen corrientes; las corrientes producen enlaces de flujo y par electromagnético; el par produce aceleración mecánica; la velocidad mecánica modifica a su vez los términos eléctricos mediante la velocidad eléctrica.

Por ello, la planta no puede tratarse como tres modelos independientes. Su comportamiento surge de la interacción entre ellos.

El acoplamiento fundamental puede representarse como:

\begin{equation}
(v_a,v_b,v_c)
\rightarrow
(i_a,i_b,i_c)
\rightarrow
(\psi,T_e)
\rightarrow
\omega_m
\rightarrow
\theta_m.
\end{equation}

A su vez, la posición y velocidad del rotor regresan al dominio eléctrico mediante:

\begin{equation}
\theta_e=p\theta_m,
\qquad
\omega_e=p\omega_m.
\end{equation}

Esta realimentación física constituye una de las características fundamentales del modelo PMSM.

\subsection{Variables eléctricas trifásicas}

La máquina posee tres fases eléctricas, convencionalmente identificadas como $a$, $b$ y $c$.

Las variables de tensión y corriente pueden expresarse como:

\begin{equation}
\mathbf{v}_{abc}
=
\begin{bmatrix}
v_a\\
v_b\\
v_c
\end{bmatrix},
\qquad
\mathbf{i}_{abc}
=
\begin{bmatrix}
i_a\\
i_b\\
i_c
\end{bmatrix}.
\end{equation}

De forma equivalente, los enlaces de flujo pueden representarse mediante:

\begin{equation}
\boldsymbol{\psi}_{abc}
=
\begin{bmatrix}
\psi_a\\
\psi_b\\
\psi_c
\end{bmatrix}.
\end{equation}

La descripción directa en coordenadas $abc$ conserva la naturaleza trifásica de la máquina, pero genera ecuaciones con una dependencia angular que resulta menos conveniente para el análisis y control.

Por esta razón se realiza una transformación de coordenadas.

La representación utilizada en la planta es:

\begin{equation}
abc
\rightarrow
\alpha\beta
\rightarrow
dq.
\end{equation}

El paso por $\alpha\beta$ transforma las tres variables de fase en dos componentes ortogonales en un marco estacionario. Posteriormente, la transformación de Park lleva esas componentes a un marco que gira con el rotor.

El objetivo no es alterar la física de la máquina, sino cambiar la base matemática utilizada para describirla.

\subsection{Ley fundamental de tensión}

El punto de partida para la dinámica eléctrica es la relación entre tensión, resistencia y enlace de flujo.

Para una fase, la relación básica puede escribirse como:

\begin{equation}
v=Ri+\frac{d\psi}{dt}.
\end{equation}

Esta expresión contiene dos contribuciones principales.

La primera es la caída resistiva:

\begin{equation}
v_R=Ri.
\end{equation}

La segunda es la tensión asociada a la variación temporal del enlace de flujo:

\begin{equation}
v_{\psi}=\frac{d\psi}{dt}.
\end{equation}

La expresión completa puede interpretarse, por tanto, como el resultado de aplicar una tensión que debe vencer tanto la caída resistiva de los devanados como la fuerza electromotriz asociada a la variación de los enlaces de flujo.

Para una máquina trifásica, esta relación se aplica a las tres fases. Sin embargo, cuando el sistema se expresa en coordenadas rotantes, la derivada de los enlaces de flujo debe considerar también el movimiento del sistema de referencia.

Este hecho dará origen a los términos de acoplamiento asociados a $\omega_e$ en las ecuaciones $dq$.

\subsection{Enlaces de flujo de la PMSM}

En el nivel M4, los enlaces de flujo en el marco $dq$ se representan mediante:

\begin{equation}
\psi_d=L_di_d+\lambda_m
\end{equation}

y:

\begin{equation}
\psi_q=L_qi_q.
\end{equation}

La primera expresión contiene dos contribuciones.

La primera es el flujo asociado a la corriente del eje $d$:

\begin{equation}
\psi_{d,i}=L_di_d.
\end{equation}

La segunda corresponde al flujo producido por los imanes permanentes:

\begin{equation}
\psi_{d,PM}=\lambda_m.
\end{equation}

Por tanto:

\begin{equation}
\psi_d
=
\psi_{d,i}
+
\psi_{d,PM}.
\end{equation}

En el modelo M4, el flujo permanente se representa mediante un único parámetro constante $\lambda_m$.

La componente $q$ se representa como:

\begin{equation}
\psi_q=L_qi_q.
\end{equation}

Esta formulación constituye un modelo concentrado. En una representación de mayor fidelidad, los enlaces de flujo podrían expresarse como funciones de las corrientes y de la posición del rotor:

\begin{equation}
\psi_d
=
\psi_d(i_d,i_q,\theta_e,T,\ldots)
\end{equation}

y:

\begin{equation}
\psi_q
=
\psi_q(i_d,i_q,\theta_e,T,\ldots).
\end{equation}

Estas expresiones más generales constituyen precisamente el camino hacia los mapas electromagnéticos y los modelos no lineales.

\subsection{Transformación Clarke}

La transformación de Clarke convierte las variables trifásicas $abc$ en un sistema ortogonal estacionario $\alpha\beta$.

Para las corrientes, la implementación utiliza una representación equivalente en la que:

\begin{equation}
i_\alpha=i_a
\end{equation}

y:

\begin{equation}
i_\beta=
\frac{i_b-i_c}{\sqrt{3}}.
\end{equation}

La misma idea puede aplicarse a otras variables trifásicas.

La transformación reduce la descripción de tres componentes de fase a dos componentes ortogonales, bajo las condiciones de balance apropiadas para la representación utilizada.

El plano $\alpha\beta$ permanece fijo respecto al estator. Por ello, todavía no incorpora la rotación del rotor.

Conceptualmente:

\begin{equation}
abc
\rightarrow
\alpha\beta
\end{equation}

cambia la base de representación, pero no introduce todavía el movimiento del sistema de coordenadas asociado al rotor.

\subsection{Transformación Park}

La transformación de Park lleva las variables $\alpha\beta$ a un marco de referencia rotante $dq$.

De manera conceptual:

\begin{equation}
\alpha\beta
\rightarrow
dq.
\end{equation}

El ángulo utilizado para esta transformación es el ángulo eléctrico:

\begin{equation}
\theta_e=p\theta_m.
\end{equation}

La razón de utilizar el ángulo eléctrico, en lugar del ángulo mecánico directamente, es que el campo electromagnético se relaciona con la periodicidad eléctrica de la máquina.

La transformación puede representarse de manera general como:

\begin{equation}
\begin{bmatrix}
i_d\\
i_q
\end{bmatrix}
=
\mathbf{T}(\theta_e)
\begin{bmatrix}
i_\alpha\\
i_\beta
\end{bmatrix}.
\end{equation}

La matriz $\mathbf{T}(\theta_e)$ depende del convenio de transformación seleccionado.

Este punto es importante para la trazabilidad del modelo: diferentes convenciones de Clarke y Park pueden producir signos, factores de escala o definiciones de eje diferentes. Por ello, la transformación utilizada por el código debe considerarse parte de la definición del modelo y no una operación puramente intercambiable.

En la implementación actual, la planta utiliza las funciones de transformación disponibles en el paquete de control para realizar la conversión entre los marcos correspondientes.

\subsection{\texorpdfstring{Significado físico de los ejes $d$ y $q$}{Significado físico de los ejes d y q}}

Una vez realizado el cambio de coordenadas, las variables eléctricas se expresan respecto a un marco asociado al rotor.

El eje $d$ se alinea con la referencia magnética principal del rotor y, en el modelo utilizado, contiene la contribución del flujo de los imanes permanentes.

El eje $q$ es ortogonal al eje $d$ y participa directamente en la producción del par electromagnético.

Esta separación permite interpretar el comportamiento de la máquina mediante dos componentes de corriente:

\begin{equation}
i_d
\quad\text{y}\quad
i_q.
\end{equation}

En un modelo PMSM simplificado con $L_d=L_q$, la relación entre el par y la corriente $q$ se vuelve particularmente directa.

Sin embargo, esta interpretación no debe llevar a concluir que $i_d$ carezca de importancia en una máquina real. En presencia de saliencia, saturación, límites de tensión, debilitamiento de campo y otros fenómenos, la componente $d$ puede desempeñar un papel fundamental.

\subsection{Velocidad mecánica y velocidad eléctrica}

La máquina posee una velocidad mecánica $\omega_m$ y una velocidad eléctrica $\omega_e$.

La relación entre ambas está determinada por el número de pares de polos:

\begin{equation}
\omega_e=p\omega_m.
\end{equation}

Para la máquina de referencia se tiene:

\begin{equation}
2p=8,
\end{equation}

por lo que:

\begin{equation}
p=4.
\end{equation}

En consecuencia:

\begin{equation}
\omega_e=4\omega_m.
\end{equation}

De forma equivalente, la posición eléctrica es:

\begin{equation}
\theta_e=p\theta_m.
\end{equation}

Estas relaciones son fundamentales porque conectan la dinámica mecánica con las ecuaciones eléctricas.

Un incremento de velocidad mecánica produce un incremento de velocidad eléctrica y, por tanto, modifica los términos de fuerza electromotriz y acoplamiento presentes en el modelo $dq$.

\subsection{\texorpdfstring{Obtención de las ecuaciones eléctricas en $dq$}{Obtención de las ecuaciones eléctricas en dq}}

Partiendo de la relación general:

\begin{equation}
v=Ri+\frac{d\psi}{dt},
\end{equation}

la transformación al marco rotante introduce términos asociados a la velocidad angular eléctrica.

Las ecuaciones de tensión en coordenadas $dq$ pueden expresarse como:

\begin{equation}
v_d
=
R_si_d
+
\frac{d\psi_d}{dt}
-
\omega_e\psi_q
\end{equation}

y:

\begin{equation}
v_q
=
R_si_q
+
\frac{d\psi_q}{dt}
+
\omega_e\psi_d.
\end{equation}

Los términos:

\begin{equation}
-\omega_e\psi_q
\end{equation}

y:

\begin{equation}
+\omega_e\psi_d
\end{equation}

aparecen debido al hecho de que el sistema de coordenadas $dq$ está girando.

No constituyen pérdidas adicionales ni componentes externas. Son términos cinemáticos que aparecen al expresar las leyes eléctricas en un sistema de referencia rotante.

\subsection{Sustitución de los enlaces de flujo}

Utilizando:

\begin{equation}
\psi_d=L_di_d+\lambda_m
\end{equation}

y:

\begin{equation}
\psi_q=L_qi_q,
\end{equation}

la derivada del enlace de flujo del eje $d$ queda:

\begin{equation}
\frac{d\psi_d}{dt}
=
L_d\frac{di_d}{dt}
\end{equation}

cuando $L_d$ y $\lambda_m$ se consideran constantes.

De forma equivalente:

\begin{equation}
\frac{d\psi_q}{dt}
=
L_q\frac{di_q}{dt}.
\end{equation}

Sustituyendo estas relaciones en las ecuaciones de tensión se obtiene:

\begin{equation}
v_d
=
R_si_d
+
L_d\frac{di_d}{dt}
-
\omega_eL_qi_q
\end{equation}

y:

\begin{equation}
v_q
=
R_si_q
+
L_q\frac{di_q}{dt}
+
\omega_e
\left(
L_di_d+\lambda_m
\right).
\end{equation}

Estas son las ecuaciones eléctricas utilizadas como base de la implementación M4.

\subsection{Forma diferencial de las corrientes}

Despejando las derivadas de las corrientes:

\begin{equation}
\frac{di_d}{dt}
=
\frac{
v_d-R_si_d+\omega_eL_qi_q
}{L_d}
\end{equation}

y:

\begin{equation}
\frac{di_q}{dt}
=
\frac{
v_q-R_si_q-\omega_e(L_di_d+\lambda_m)
}{L_q}.
\end{equation}

Estas expresiones son directamente implementables en una simulación temporal.

Además, muestran explícitamente los principales acoplamientos de la máquina.

La ecuación del eje $d$ contiene el término:

\begin{equation}
+\omega_eL_qi_q,
\end{equation}

mientras que la ecuación del eje $q$ contiene:

\begin{equation}
-\omega_e(L_di_d+\lambda_m).
\end{equation}

Por esta razón, las dinámicas de $i_d$ e $i_q$ no son completamente independientes.

Esta interacción constituye una de las razones por las que el control de corriente de una PMSM requiere compensar los términos de acoplamiento cuando se busca un comportamiento desacoplado.

\subsection{Fuerza electromotriz y flujo permanente}

El término:

\begin{equation}
\omega_e\lambda_m
\end{equation}

que aparece en la ecuación del eje $q$ representa la contribución asociada al flujo de los imanes permanentes cuando el rotor gira.

Este término puede interpretarse como parte de la fuerza electromotriz asociada al movimiento del flujo permanente respecto a los devanados.

Por ello, el flujo permanente no es simplemente un parámetro estático sin efecto dinámico. Su presencia modifica directamente la tensión requerida para mantener determinadas corrientes cuando la máquina gira.

A medida que aumenta $\omega_e$, también aumenta la magnitud del término:

\begin{equation}
\omega_e\lambda_m.
\end{equation}

Esta relación es fundamental para comprender posteriormente los límites de tensión y la necesidad de estrategias como debilitamiento de campo en regiones de alta velocidad.

\subsection{Obtención del flujo permanente a partir de la constante de par}

El valor de $\lambda_m$ utilizado actualmente por M4 no procede directamente de una tabla que publique explícitamente el enlace de flujo del imán.

Se obtiene a partir de la constante de par documentada para la máquina.

La documentación de referencia proporciona:

\begin{equation}
k_T=0.91\;\mathrm{N\,m/A}
\end{equation}

y el número de pares de polos:

\begin{equation}
p=4.
\end{equation}

Para la convención utilizada en este modelo, la relación empleada es:

\begin{equation}
k_T
=
\frac{3}{2}\sqrt{2}\,p\,\lambda_m.
\end{equation}

Despejando:

\begin{equation}
\lambda_m
=
\frac{k_T}
{\frac{3}{2}\sqrt{2}\,p}.
\end{equation}

Sustituyendo los valores:

\begin{equation}
\lambda_m
=
\frac{0.91}
{\frac{3}{2}\sqrt{2}(4)}.
\end{equation}

Se obtiene aproximadamente:

\begin{equation}
\boxed{
\lambda_m\approx0.107245\;\mathrm{Wb}
}
\end{equation}

Este valor es el utilizado actualmente por la planta M4.

Es importante distinguir entre una magnitud documentada directamente y una magnitud derivada. La documentación proporciona $k_T$, mientras que $\lambda_m$ se obtiene mediante una relación matemática dependiente de la convención de corriente y del modelo adoptado.

Por esta razón, aunque el valor numérico esté incorporado en M4 como confirmado, su procedencia se mantiene identificada como derivada y su trazabilidad documental continúa marcada para verificación.

\subsection{Comprobación independiente mediante la constante de tensión}

La documentación de la máquina también proporciona una constante de tensión:

\begin{equation}
k_E=58\;\mathrm{V/(1000\,rpm)}
\end{equation}

a 20\,${}^{\circ}$C.

Esta magnitud ofrece una segunda referencia para estudiar el flujo equivalente de la máquina.

Sin embargo, la conversión entre $k_E$ y $\lambda_m$ depende de la convención empleada para distinguir tensión de fase, tensión de línea, valores RMS, valores pico y velocidad angular eléctrica.

Por ello, una comparación directa entre $k_T$ y $k_E$ sin considerar sus convenciones puede generar discrepancias aparentes.

La constante $k_E$ proporciona aproximadamente una segunda estimación del flujo equivalente, pero no se utiliza actualmente para sustituir el valor M4 derivado de $k_T$.

La diferencia observada entre ambas estimaciones se conserva como parte de la trazabilidad y deberá analizarse posteriormente considerando las condiciones de ensayo, temperatura, tolerancias documentadas y convenciones de magnitud.

Este tratamiento evita forzar una coincidencia numérica artificial entre dos magnitudes que pueden haber sido especificadas bajo condiciones diferentes.

\subsection{Obtención del par electromagnético}

El par electromagnético de una PMSM puede expresarse en el marco $dq$ mediante:

\begin{equation}
T_e
=
\frac{3}{2}p
\left(
\psi_di_q-\psi_qi_d
\right).
\end{equation}

Esta expresión relaciona directamente los enlaces de flujo con las corrientes y permite representar tanto la contribución asociada a los imanes permanentes como la contribución de reluctancia.

Sustituyendo:

\begin{equation}
\psi_d=L_di_d+\lambda_m
\end{equation}

y:

\begin{equation}
\psi_q=L_qi_q,
\end{equation}

se obtiene:

\begin{equation}
T_e
=
\frac{3}{2}p
\left[
(L_di_d+\lambda_m)i_q
-
L_qi_qi_d
\right].
\end{equation}

Reordenando:

\begin{equation}
T_e
=
\frac{3}{2}p
\left[
\lambda_mi_q
+
(L_d-L_q)i_di_q
\right].
\end{equation}

Esta forma permite identificar dos contribuciones:

\begin{equation}
T_{PM}
=
\frac{3}{2}p\lambda_mi_q
\end{equation}

y:

\begin{equation}
T_{rel}
=
\frac{3}{2}p(L_d-L_q)i_di_q.
\end{equation}

Por tanto:

\begin{equation}
T_e=T_{PM}+T_{rel}.
\end{equation}

La primera contribución está asociada al flujo de los imanes permanentes. La segunda corresponde al par de reluctancia producido por la diferencia entre las inductancias de los ejes $d$ y $q$.

\subsection{\texorpdfstring{Consecuencia de la parametrización $L_d=L_q$}{Consecuencia de la parametrización Ld = Lq}}

El modelo M4 utiliza actualmente:

\begin{equation}
L_d=L_q=1.2\;\mathrm{mH}.
\end{equation}

Por tanto:

\begin{equation}
L_d-L_q=0.
\end{equation}

La contribución de reluctancia queda anulada dentro de este modelo:

\begin{equation}
T_{rel}=0.
\end{equation}

En consecuencia:

\begin{equation}
\boxed{
T_e=
\frac{3}{2}p\lambda_mi_q
}
\end{equation}

para la parametrización M4 actual.

Esta simplificación debe interpretarse correctamente.

No significa que una PMSM real necesariamente tenga:

\begin{equation}
L_d=L_q.
\end{equation}

Tampoco significa que una máquina real carezca de saliencia o de efectos de saturación.

Significa únicamente que el primer nivel de la planta utiliza un modelo concentrado en el cual se han establecido ambas inductancias con el mismo valor.

Cuando se disponga de información suficiente para construir:

\begin{equation}
L_d=L_d(i_d,i_q,\theta_e)
\end{equation}

y:

\begin{equation}
L_q=L_q(i_d,i_q,\theta_e),
\end{equation}

la contribución de reluctancia podrá reaparecer de forma dependiente del estado y de la posición.

\subsection{Dinámica mecánica}

La parte mecánica se obtiene mediante el balance de torques.

El principio básico es:

\begin{equation}
J\frac{d\omega_m}{dt}
=
\sum T.
\end{equation}

En la representación M4 actual se consideran el par electromagnético y el par de carga:

\begin{equation}
J\frac{d\omega_m}{dt}
=
T_e-T_{\mathrm{load}}.
\end{equation}

Despejando:

\begin{equation}
\frac{d\omega_m}{dt}
=
\frac{T_e-T_{\mathrm{load}}}{J}.
\end{equation}

La posición mecánica se obtiene integrando la velocidad:

\begin{equation}
\frac{d\theta_m}{dt}
=
\omega_m.
\end{equation}

Estas dos ecuaciones cierran el acoplamiento entre los dominios electromagnético y mecánico.

El par generado por las corrientes modifica la velocidad y la velocidad modifica la frecuencia eléctrica utilizada por las ecuaciones de corriente.

\subsection{Modelo mecánico y parámetros no incorporados}

En una representación más completa, la ecuación mecánica podría incluir pérdidas dependientes de la velocidad y otros pares resistentes:

\begin{equation}
J\frac{d\omega_m}{dt}
=
T_e
-
T_{\mathrm{load}}
-
B\omega_m
-
T_{\mathrm{fric}}.
\end{equation}

Sin embargo, los parámetros $B$ y $T_{\mathrm{fric}}$ no se consideran confirmados en la parametrización M4.

Por esta razón, la implementación actual utiliza:

\begin{equation}
J\frac{d\omega_m}{dt}
=
T_e-T_{\mathrm{load}}.
\end{equation}

Esta decisión mantiene la trazabilidad del modelo y evita introducir valores históricos no verificados como si fueran propiedades confirmadas de la máquina.

\subsection{Modelo completo de ecuaciones diferenciales}

La planta M4 puede expresarse mediante el conjunto:

\begin{equation}
\frac{di_d}{dt}
=
\frac{
v_d-R_si_d+p\omega_mL_qi_q
}{L_d},
\end{equation}

\begin{equation}
\frac{di_q}{dt}
=
\frac{
v_q-R_si_q-p\omega_m(L_di_d+\lambda_m)
}{L_q},
\end{equation}

\begin{equation}
\frac{d\omega_m}{dt}
=
\frac{
T_e-T_{\mathrm{load}}
}{J},
\end{equation}

\begin{equation}
\frac{d\theta_m}{dt}
=
\omega_m,
\end{equation}

con:

\begin{equation}
T_e
=
\frac{3}{2}p
\left[
\lambda_mi_q+
(L_d-L_q)i_di_q
\right].
\end{equation}

Estas ecuaciones constituyen el núcleo matemático de la planta concentrada utilizada en M4.

\subsection{Formulación en espacio de estados}

Para expresar la planta en forma de espacio de estados se define el vector:

\begin{equation}
x=
\begin{bmatrix}
i_d\\
i_q\\
\omega_m\\
\theta_m
\end{bmatrix}.
\end{equation}

Las entradas dinámicas principales pueden definirse como:

\begin{equation}
u=
\begin{bmatrix}
v_d\\
v_q\\
T_{\mathrm{load}}
\end{bmatrix}.
\end{equation}

El conjunto de parámetros puede escribirse como:

\begin{equation}
\Theta=
\left\{
R_s,L_d,L_q,\lambda_m,p,J
\right\}.
\end{equation}

La planta puede entonces representarse de manera general como:

\begin{equation}
\dot{x}=f(x,u,\Theta).
\end{equation}

En forma desarrollada:

\begin{equation}
\dot{x}
=
\begin{bmatrix}
\dfrac{
v_d-R_si_d+p\omega_mL_qi_q
}{L_d}
\\[10pt]
\dfrac{
v_q-R_si_q-p\omega_m(L_di_d+\lambda_m)
}{L_q}
\\[10pt]
\dfrac{
T_e-T_{\mathrm{load}}
}{J}
\\[10pt]
\omega_m
\end{bmatrix}.
\end{equation}

El vector de salida puede definirse, por ejemplo, como:

\begin{equation}
y=
\begin{bmatrix}
i_d\\
i_q\\
\omega_m\\
\theta_m\\
T_e
\end{bmatrix}.
\end{equation}

Las corrientes trifásicas pueden calcularse adicionalmente mediante las transformaciones inversas:

\begin{equation}
(i_d,i_q,\theta_e)
\rightarrow
(i_\alpha,i_\beta)
\rightarrow
(i_a,i_b,i_c).
\end{equation}

Esta formulación permite identificar claramente qué variables pertenecen al estado interno de la planta y cuáles corresponden a entradas o salidas.

\subsection{Relación entre el modelo matemático y la implementación Python}

La formulación anterior se implementa en \texttt{PMSMBase} mediante cuatro estados internos:

\begin{equation}
(i_d,i_q,\omega_m,\theta_m).
\end{equation}

Las derivadas se calculan en el método interno de la planta y posteriormente se realiza la integración temporal.

La correspondencia conceptual es directa:

\begin{center}
\begin{tabular}{ll}
\textbf{Modelo matemático} & \textbf{Implementación} \\
\hline
$i_d$ & \texttt{\_id} \\
$i_q$ & \texttt{\_iq} \\
$\omega_m$ & \texttt{\_omega\_m} \\
$\theta_m$ & \texttt{\_theta\_m} \\
$T_e$ & \texttt{electromagnetic\_torque()} \\
$\dot{x}$ & \texttt{\_derivatives()} \\
$x_{k+1}$ & \texttt{step()} \\
\end{tabular}
\end{center}

Esta correspondencia es importante para la reproducibilidad porque permite verificar que cada elemento matemático tenga una representación concreta dentro del código.

La ecuación por sí sola no constituye todavía una planta ejecutable. La implementación requiere definir estados, entradas, condiciones iniciales, parámetros, transformaciones y método de integración.

\subsection{Condiciones iniciales}

La implementación M4 utiliza inicialmente:

\begin{equation}
i_d(0)=0,
\end{equation}

\begin{equation}
i_q(0)=0,
\end{equation}

\begin{equation}
\omega_m(0)=0,
\end{equation}

y:

\begin{equation}
\theta_m(0)=0.
\end{equation}

Estas condiciones corresponden al estado inicial empleado por la clase base de la planta.

No deben confundirse con condiciones nominales de operación de la máquina.

Los datos nominales del fabricante, como:

\begin{equation}
I_N=53\;\mathrm{A}
\end{equation}

y:

\begin{equation}
M_N=46\;\mathrm{N\,m},
\end{equation}

describen una condición de operación de referencia de la máquina y no constituyen automáticamente condiciones iniciales para la integración de la planta.

Esta distinción permite separar la definición física de la máquina de la configuración específica de cada experimento.

\subsection{Parámetros documentales, derivados y decisiones de modelado}

La formulación matemática depende de tres tipos de información que deben mantenerse separados.

\subsubsection{Parámetros documentales}

Son magnitudes obtenidas directamente de la documentación de referencia o de fuentes primarias del proyecto.

Entre ellas se encuentran:

\begin{equation}
R,
\quad
L_D,
\quad
k_T,
\quad
k_E,
\quad
J,
\quad
2p.
\end{equation}

Estos valores deben conservar la referencia documental correspondiente.

\subsubsection{Parámetros derivados}

Son magnitudes calculadas a partir de parámetros documentales y de una formulación explícita.

Un ejemplo central es:

\begin{equation}
\lambda_m
=
\frac{k_T}
{\frac{3}{2}\sqrt{2}p}.
\end{equation}

El valor resultante depende de la convención de modelado utilizada, por lo que su trazabilidad debe conservar tanto la fuente original como la ecuación empleada.

\subsubsection{Decisiones de modelado}

Son decisiones necesarias para construir una representación ejecutable cuando la documentación disponible no proporciona directamente toda la información requerida.

En M4, un ejemplo importante es:

\begin{equation}
L_d=L_q=1.2\;\mathrm{mH}.
\end{equation}

La documentación proporciona una inductancia $L_D$, pero no establece en el material utilizado una pareja independiente de valores $L_d$ y $L_q$. Por ello, la igualdad se utiliza como decisión de modelado para el nivel concentrado actual.

Esta clasificación evita presentar una hipótesis de modelado como si fuera un dato físico medido directamente.

\subsection{Supuestos del nivel M4}

La formulación actual se apoya en varios supuestos explícitos:

\begin{enumerate}
    \item $R_s$ se considera constante;
    \item $L_d$ y $L_q$ se consideran constantes;
    \item el flujo permanente $\lambda_m$ se considera constante;
    \item no se representa explícitamente la saturación magnética;
    \item no se representan armónicos espaciales;
    \item no se introduce dependencia angular de los enlaces de flujo;
    \item no se introducen pérdidas electromagnéticas detalladas;
    \item $B$ y $T_{\mathrm{fric}}$ no forman parte de la dinámica M4;
    \item la temperatura no modifica todavía los parámetros durante la simulación;
    \item la máquina se representa mediante un modelo concentrado.
\end{enumerate}

Estos supuestos delimitan exactamente qué significa el modelo M4.

Por tanto, cuando una simulación reproduce una determinada respuesta, dicha respuesta corresponde al conjunto de ecuaciones y supuestos definidos aquí, no necesariamente a todos los fenómenos presentes en la máquina física.

\subsection{Qué fenómenos aparecen al aumentar la fidelidad}

Una máquina real puede presentar dependencias más generales:

\begin{equation}
\psi_d
=
\psi_d(i_d,i_q,\theta_e,T),
\end{equation}

\begin{equation}
\psi_q
=
\psi_q(i_d,i_q,\theta_e,T),
\end{equation}

y:

\begin{equation}
T_e
=
T_e(i_d,i_q,\theta_e,T).
\end{equation}

En este caso, las inductancias dejan de ser necesariamente constantes:

\begin{equation}
L_d
=
L_d(i_d,i_q,\theta_e,T)
\end{equation}

y:

\begin{equation}
L_q
=
L_q(i_d,i_q,\theta_e,T).
\end{equation}

También pueden aparecer armónicos, pérdidas dependientes de la velocidad y la frecuencia, variaciones térmicas y efectos geométricos.

La transición hacia estas representaciones no consiste simplemente en añadir términos arbitrarios a las ecuaciones actuales. Requiere disponer de información física suficiente para identificar las nuevas dependencias.

Por ello, el siguiente nivel natural del proyecto es la construcción de mapas electromagnéticos.

\subsection{Relación con los mapas electromagnéticos}

En un modelo basado en mapas, los enlaces de flujo pueden representarse mediante superficies o tablas:

\begin{equation}
\psi_d=f_d(i_d,i_q,\theta_e)
\end{equation}

y:

\begin{equation}
\psi_q=f_q(i_d,i_q,\theta_e).
\end{equation}

El par puede obtenerse entonces de la relación electromagnética correspondiente o de un mapa:

\begin{equation}
T_e=f_T(i_d,i_q,\theta_e).
\end{equation}

Esta representación permite introducir no linealidades que el modelo concentrado no puede reproducir.

Sin embargo, una tabla recuperada de una fuente histórica no debe activarse automáticamente como si fuera un mapa validado. Es necesario determinar:

\begin{itemize}
    \item qué variable representa cada eje;
    \item qué unidades utiliza;
    \item qué sistema de referencia corresponde;
    \item qué condiciones de temperatura se aplicaron;
    \item qué rango de corriente representa;
    \item si existe dependencia angular;
    \item cuál es su procedencia;
    \item qué método se utilizó para obtenerla.
\end{itemize}

La construcción de mapas constituye por tanto una etapa de identificación y validación, no solamente una tarea de programación.

\subsection{Relación con FEM y ROM}

Los modelos FEM permiten obtener información espacialmente distribuida del campo electromagnético a partir de la geometría y materiales de la máquina.

Una eventual cadena de construcción puede representarse como:

\begin{equation}
\text{geometría}
\rightarrow
\text{materiales}
\rightarrow
\text{FEM}
\rightarrow
\text{soluciones electromagnéticas}
\rightarrow
\text{mapas}
\rightarrow
\text{ROM}.
\end{equation}

El modelo reducido, o ROM, puede utilizar posteriormente la información obtenida mediante FEM sin requerir resolver el problema electromagnético completo durante cada paso de una simulación de control.

En este contexto, el modelo M4 no compite con FEM. Constituye el nivel inicial sobre el cual se puede establecer la evolución de fidelidad.

\subsection{Separación entre formulación de planta y formulación del controlador}

Las ecuaciones obtenidas en este capítulo describen la planta.

El controlador utiliza estas relaciones para decidir las acciones que deben aplicarse, pero sus ecuaciones no forman parte de la dinámica física de la máquina.

Por ejemplo, el FOC puede generar referencias y tensiones:

\begin{equation}
(i_d^\ast,i_q^\ast)
\rightarrow
(v_d^\ast,v_q^\ast),
\end{equation}

mientras que la planta responde mediante:

\begin{equation}
(v_d,v_q)
\rightarrow
(i_d,i_q,\omega_m,\theta_m).
\end{equation}

Esta separación permite que posteriormente otro controlador utilice la misma planta:

\begin{equation}
\mathrm{FOC}
\quad\text{o}\quad
\mathrm{MPC}
\quad\text{o}\quad
\mathrm{supervisor\ ANFIS}
\quad\text{o}\quad
\mathrm{otra\ estrategia}.
\end{equation}

La planta permanece definida por su física y sus parámetros, mientras que el controlador se define por sus objetivos y leyes de control.

\subsection{Consistencia con la arquitectura por contratos}

La formulación matemática también establece qué variables deben atravesar la frontera entre la planta y los controladores.

En el dominio $dq$, las variables principales son:

\begin{equation}
v_d,\quad
v_q,\quad
i_d,\quad
i_q,\quad
\omega_m,\quad
\theta_m,\quad
T_e.
\end{equation}

En la frontera eléctrica trifásica aparecen:

\begin{equation}
v_a,v_b,v_c
\end{equation}

y:

\begin{equation}
i_a,i_b,i_c.
\end{equation}

Estas variables constituyen información que puede ser transportada mediante los contratos definidos por la arquitectura.

El modelo matemático, por tanto, no solamente describe la dinámica interna, sino que ayuda a establecer una frontera física coherente entre la planta y los sistemas que interactúan con ella.

\subsection{Relación entre modelo matemático y Digital Twin}

La formulación matemática constituye una de las capas fundamentales del Digital Twin.

La relación completa puede representarse como:

\begin{equation}
\text{Máquina física}
\rightarrow
\text{datos}
\rightarrow
\text{parámetros}
\rightarrow
\text{ecuaciones}
\rightarrow
\text{modelo computacional}.
\end{equation}

La etapa posterior incorpora validación y aumento de fidelidad:

\begin{equation}
\text{modelo computacional}
\rightarrow
\text{validación}
\rightarrow
\text{mapas}
\rightarrow
\text{modelo no lineal}
\rightarrow
\text{FEM/ROM}.
\end{equation}

Finalmente, la planta puede convertirse en el núcleo de una plataforma de experimentación para diferentes estrategias de control.

El Digital Twin no se define, por tanto, únicamente por la existencia de una ecuación diferencial. Se construye mediante la combinación de:

\begin{equation}
\text{trazabilidad}
+
\text{modelo físico}
+
\text{datos}
+
\text{validación}
+
\text{evolución de fidelidad}.
\end{equation}

\subsection{Resumen de la derivación}

La formulación desarrollada puede resumirse mediante la siguiente secuencia:

\begin{equation}
v=Ri+\frac{d\psi}{dt}
\end{equation}

\begin{equation}
\Downarrow
\end{equation}

\begin{equation}
\psi_d=L_di_d+\lambda_m,
\qquad
\psi_q=L_qi_q
\end{equation}

\begin{equation}
\Downarrow
\end{equation}

\begin{equation}
v_d=
R_si_d+
L_d\frac{di_d}{dt}
-
\omega_eL_qi_q
\end{equation}

\begin{equation}
v_q=
R_si_q+
L_q\frac{di_q}{dt}
+
\omega_e(L_di_d+\lambda_m)
\end{equation}

\begin{equation}
\Downarrow
\end{equation}

\begin{equation}
T_e=
\frac{3}{2}p
\left[
\lambda_mi_q+
(L_d-L_q)i_di_q
\right]
\end{equation}

\begin{equation}
\Downarrow
\end{equation}

\begin{equation}
J\frac{d\omega_m}{dt}
=
T_e-T_{\mathrm{load}}
\end{equation}

\begin{equation}
\Downarrow
\end{equation}

\begin{equation}
x=
\begin{bmatrix}
i_d&i_q&\omega_m&\theta_m
\end{bmatrix}^{T}.
\end{equation}

Esta cadena constituye la formulación matemática fundamental de la planta M4.

\subsection{Alcance de la formulación actual}

La formulación presentada describe de manera completa el nivel concentrado actualmente implementado, pero no pretende cerrar todavía la representación de alta fidelidad de la máquina.

El modelo matemático actual permite:

\begin{itemize}
    \item representar las dinámicas eléctricas principales;
    \item representar la interacción electromecánica;
    \item calcular el par electromagnético de referencia;
    \item representar la evolución mecánica;
    \item definir una representación de estados reproducible;
    \item implementar la planta en Python;
    \item conectar la planta con estrategias de control.
\end{itemize}

Todavía no permite representar explícitamente:

\begin{itemize}
    \item saturación magnética;
    \item variación de inductancias con corriente;
    \item dependencia angular detallada;
    \item armónicos espaciales;
    \item pérdidas electromagnéticas detalladas;
    \item evolución térmica acoplada;
    \item mapas FEM validados.
\end{itemize}

La incorporación de estos fenómenos constituye la evolución natural del modelo, pero deberá realizarse únicamente cuando exista información física suficiente para parametrizarlos y validarlos.

\subsection{Cierre del capítulo}

La formulación matemática establece la conexión formal entre la máquina PMSM de referencia, sus parámetros y la implementación computacional utilizada en M4.

El proceso parte de la ley fundamental de tensión, introduce los enlaces de flujo asociados a las corrientes y a los imanes permanentes, transforma las variables trifásicas a coordenadas $dq$, incorpora la velocidad eléctrica y obtiene las ecuaciones de corriente. A partir de los enlaces de flujo se deriva el par electromagnético y, finalmente, el balance mecánico completa el sistema dinámico.

La formulación también permite identificar con claridad qué parte del modelo corresponde a datos documentales, qué parte corresponde a magnitudes derivadas y qué parte constituye una decisión de modelado.

Esta distinción resulta esencial para la trazabilidad del Digital Twin.

El siguiente paso consiste en llevar estas ecuaciones a una representación explícita de espacio de estados y relacionarla con la implementación de referencia MATLAB/Simscape y con la reconstrucción reproducible en Python. Posteriormente, la validación experimental y documental permitirá determinar qué elementos del modelo concentrado deben conservarse y cuáles deberán sustituirse por representaciones de mayor fidelidad.

La evolución prevista permanece:

\begin{equation}
\mathrm{M4}_{\mathrm{concentrated}}
\rightarrow
\mathrm{maps}
\rightarrow
\mathrm{nonlinear}
\rightarrow
\mathrm{FEM/ROM}
\rightarrow
\mathrm{Digital\ Twin}.
\end{equation}


\section{Representación en espacio de estados}

\subsection{Propósito de la representación en espacio de estados}

El desarrollo de la formulación matemática realizado en el capítulo anterior permite establecer las ecuaciones diferenciales que describen la dinámica eléctrica, electromagnética y mecánica de la máquina PMSM. Sin embargo, para implementar dichas ecuaciones de forma sistemática en un entorno de simulación es necesario organizar el modelo alrededor de un conjunto explícito de variables de estado, entradas, parámetros y salidas.

La representación en espacio de estados proporciona precisamente esta organización. En lugar de considerar cada ecuación de manera aislada, el modelo se expresa como una relación entre un vector de estados, un vector de entradas y un conjunto de parámetros físicos. Esta representación permite definir de forma precisa qué variables contienen la memoria dinámica de la planta, qué magnitudes actúan sobre ella y qué variables pueden observarse o entregarse a otros bloques del sistema.

En este trabajo, el espacio de estados constituye además el puente entre tres niveles de representación:

\begin{enumerate}
    \item la formulación física de la máquina PMSM;
    \item la implementación computacional reproducible en Python;
    \item la futura integración e implementación del modelo en MATLAB/Simscape y dentro del Digital Twin.
\end{enumerate}

Por tanto, el espacio de estados no se considera solamente una forma matemática alternativa, sino una definición formal de la dinámica interna de la planta.

\subsection{Del modelo físico al modelo de estados}

El procedimiento seguido parte de las ecuaciones obtenidas a partir de la ley de tensión, los enlaces de flujo, la transformación a coordenadas $dq$, la expresión del par electromagnético y el balance mecánico.

El modelo concentrado utilizado actualmente contiene cuatro variables dinámicas principales:

\begin{equation}
x =
\begin{bmatrix}
i_d &
i_q &
\omega_m &
\theta_m
\end{bmatrix}^{T}.
\end{equation}

Estas cuatro variables representan, respectivamente, las componentes de corriente en los ejes $d$ y $q$, la velocidad mecánica del rotor y la posición mecánica del rotor.

La elección de estos estados responde directamente a las ecuaciones diferenciales de la planta. Las corrientes determinan la evolución electromagnética, mientras que la velocidad y la posición determinan la evolución mecánica y la orientación del sistema de coordenadas rotante.

De esta manera, cada componente del vector de estados tiene una interpretación física directa.

\subsection{Estados eléctricos}

Los primeros dos estados corresponden a las corrientes expresadas en el marco rotante:

\begin{equation}
x_1=i_d.
\end{equation}

\begin{equation}
x_2=i_q.
\end{equation}

La corriente $i_d$ representa la componente asociada al eje directo del rotor, mientras que $i_q$ representa la componente asociada al eje en cuadratura.

En el modelo PMSM utilizado, estas corrientes no son simplemente variables matemáticas obtenidas mediante una transformación. Constituyen las variables dinámicas que determinan, junto con los parámetros magnéticos, los enlaces de flujo y el par electromagnético.

A partir de las ecuaciones obtenidas en el capítulo anterior, sus derivadas son:

\begin{equation}
\frac{di_d}{dt}
=
\frac{
v_d-R_si_d+\omega_eL_qi_q
}{
L_d
}.
\end{equation}

\begin{equation}
\frac{di_q}{dt}
=
\frac{
v_q-R_si_q-\omega_e(L_di_d+\lambda_m)
}{
L_q
}.
\end{equation}

Estas dos ecuaciones constituyen la parte eléctrica principal del espacio de estados.

\subsection{Estado de velocidad mecánica}

El tercer estado corresponde a la velocidad angular mecánica:

\begin{equation}
x_3=\omega_m.
\end{equation}

Su evolución está determinada por el balance de par de la máquina:

\begin{equation}
J\frac{d\omega_m}{dt}
=
T_e-T_{\mathrm{load}}.
\end{equation}

Por tanto:

\begin{equation}
\frac{d\omega_m}{dt}
=
\frac{
T_e-T_{\mathrm{load}}
}{
J
}.
\end{equation}

La velocidad mecánica conecta directamente la parte electromagnética con la parte mecánica del modelo. El par generado por las corrientes modifica la velocidad, mientras que la velocidad modifica a su vez la frecuencia eléctrica y, por tanto, las ecuaciones eléctricas.

Esta relación constituye uno de los principales acoplamientos físicos de la máquina.

\subsection{Estado de posición mecánica}

El cuarto estado corresponde a la posición angular mecánica:

\begin{equation}
x_4=\theta_m.
\end{equation}

Su ecuación diferencial es:

\begin{equation}
\frac{d\theta_m}{dt}
=
\omega_m.
\end{equation}

La posición del rotor tiene una función especialmente importante porque determina la orientación relativa entre el rotor y las variables eléctricas trifásicas.

A partir de la posición mecánica se obtiene la posición eléctrica mediante:

\begin{equation}
\theta_e=p\theta_m.
\end{equation}

De forma equivalente, la velocidad eléctrica se obtiene mediante:

\begin{equation}
\omega_e=p\omega_m.
\end{equation}

En estas expresiones, $p$ representa el número de pares de polos. Para la máquina de referencia se establece:

\begin{equation}
p=4,
\end{equation}

a partir de los ocho polos físicos documentados para la máquina de referencia Siemens 1FT7087.

\subsection{Vector de estados completo}

Reuniendo las cuatro variables, el estado de la planta se define como:

\begin{equation}
x =
\begin{bmatrix}
i_d\\
i_q\\
\omega_m\\
\theta_m
\end{bmatrix}.
\end{equation}

La derivada temporal del estado es:

\begin{equation}
\dot{x}
=
\begin{bmatrix}
\dot{i}_d\\
\dot{i}_q\\
\dot{\omega}_m\\
\dot{\theta}_m
\end{bmatrix}.
\end{equation}

Sustituyendo las ecuaciones obtenidas previamente:

\begin{equation}
\dot{x}
=
\begin{bmatrix}
\dfrac{v_d-R_si_d+\omega_eL_qi_q}{L_d}
\\[6pt]
\dfrac{v_q-R_si_q-\omega_e(L_di_d+\lambda_m)}{L_q}
\\[6pt]
\dfrac{T_e-T_{\mathrm{load}}}{J}
\\[6pt]
\omega_m
\end{bmatrix}.
\end{equation}

Esta expresión constituye la definición dinámica central del modelo M4.

\subsection{Entradas de la planta}

Para expresar la dinámica de manera compacta se define el vector de entradas como:

\begin{equation}
u =
\begin{bmatrix}
v_d\\
v_q\\
T_{\mathrm{load}}
\end{bmatrix}.
\end{equation}

Las dos primeras entradas corresponden a las componentes de tensión utilizadas para excitar la dinámica eléctrica del modelo. La tercera entrada representa el par de carga aplicado al eje mecánico.

La definición anterior no implica que estas señales sean generadas internamente por la planta. Por el contrario, representan magnitudes que atraviesan la frontera de la planta y que pueden provenir de otras partes de la arquitectura de control o del escenario experimental.

En una implementación mediante FOC, por ejemplo, $v_d$ y $v_q$ proceden de la acción de control calculada por el controlador. La planta no calcula el FOC; recibe las acciones que le corresponden a través de su interfaz.

\subsection{Parámetros de la planta}

El comportamiento del espacio de estados depende de un conjunto de parámetros físicos y de modelado:

\begin{equation}
\Theta =
\left\{
R_s,
L_d,
L_q,
\lambda_m,
p,
J
\right\}.
\end{equation}

En el modelo M4 estos parámetros corresponden a:

\begin{itemize}
    \item $R_s$: resistencia estatórica;
    \item $L_d$: inductancia equivalente en el eje directo;
    \item $L_q$: inductancia equivalente en el eje en cuadratura;
    \item $\lambda_m$: enlace de flujo equivalente producido por los imanes permanentes;
    \item $p$: número de pares de polos;
    \item $J$: inercia mecánica utilizada por el modelo.
\end{itemize}

La presencia explícita de estos parámetros en la formulación permite mantener separadas las propiedades de la máquina de las condiciones de operación y de los parámetros pertenecientes al controlador.

Esta separación es importante para la trazabilidad del Digital Twin.

\subsection{Par electromagnético dentro del espacio de estados}

El par electromagnético actúa como variable intermedia entre la dinámica eléctrica y la dinámica mecánica.

Su forma general para el modelo PMSM utilizado es:

\begin{equation}
T_e
=
\frac{3}{2}p
\left(
\psi_di_q-\psi_qi_d
\right).
\end{equation}

Con:

\begin{equation}
\psi_d=L_di_d+\lambda_m
\end{equation}

y:

\begin{equation}
\psi_q=L_qi_q.
\end{equation}

Por tanto:

\begin{equation}
T_e
=
\frac{3}{2}p
\left[
\lambda_mi_q+
(L_d-L_q)i_di_q
\right].
\end{equation}

En el modelo M4 se utiliza:

\begin{equation}
L_d=L_q.
\end{equation}

En consecuencia, la contribución de reluctancia desaparece en esta formulación y el par queda reducido a:

\begin{equation}
T_e
=
\frac{3}{2}p\lambda_mi_q.
\end{equation}

Esta simplificación es una consecuencia del modelo concentrado adoptado. No debe interpretarse como una afirmación general de que cualquier PMSM tenga exactamente la misma relación entre par y corriente.

Cuando posteriormente se introduzcan inductancias dependientes del estado, mapas electromagnéticos o información FEM, la expresión utilizada para calcular el par podrá adquirir una dependencia adicional con las corrientes y con la posición del rotor.

\subsection{Dependencia entre estados eléctricos y mecánicos}

Una característica fundamental del espacio de estados PMSM es que sus ecuaciones están acopladas.

La velocidad mecánica determina la velocidad eléctrica:

\begin{equation}
\omega_e=p\omega_m.
\end{equation}

La velocidad eléctrica aparece en las ecuaciones eléctricas:

\begin{equation}
\frac{di_d}{dt}
=
\frac{
v_d-R_si_d+\omega_eL_qi_q
}{
L_d
},
\end{equation}

\begin{equation}
\frac{di_q}{dt}
=
\frac{
v_q-R_si_q-\omega_e(L_di_d+\lambda_m)
}{
L_q
}.
\end{equation}

Por otra parte, las corrientes determinan el par electromagnético, y este par determina la aceleración mecánica:

\begin{equation}
\frac{d\omega_m}{dt}
=
\frac{
T_e-T_{\mathrm{load}}
}{
J
}.
\end{equation}

Existe, por tanto, una cadena de dependencia:

\begin{equation}
(i_d,i_q)
\rightarrow
T_e
\rightarrow
\omega_m
\rightarrow
\omega_e
\rightarrow
(\dot{i}_d,\dot{i}_q).
\end{equation}

Esta relación explica por qué la planta eléctrica y la planta mecánica no pueden considerarse completamente independientes.

\subsection{Forma general del modelo}

El conjunto de ecuaciones puede expresarse de forma compacta como:

\begin{equation}
\dot{x}=f(x,u,\Theta).
\end{equation}

La salida de la planta puede representarse mediante:

\begin{equation}
y=h(x,u,\Theta).
\end{equation}

Esta notación es deliberadamente general. Permite conservar la estructura del modelo aunque posteriormente se sustituyan algunas relaciones analíticas por mapas experimentales, expresiones no lineales o modelos derivados de FEM.

En el nivel M4, la función $f$ contiene las ecuaciones concentradas descritas anteriormente. En niveles posteriores, la estructura de estados podrá conservarse mientras cambien las funciones que describen los fenómenos electromagnéticos.

\subsection{Naturaleza no lineal del modelo}

Aunque el modelo M4 utiliza parámetros constantes y presenta una estructura relativamente compacta, la formulación contiene acoplamientos entre variables de estado.

Por ejemplo, la velocidad eléctrica depende de la velocidad mecánica:

\begin{equation}
\omega_e=p\omega_m,
\end{equation}

y esta magnitud multiplica las corrientes en las ecuaciones eléctricas.

De forma general:

\begin{equation}
\dot{x}=f(x,u,\Theta)
\end{equation}

no tiene por qué reducirse a una relación lineal constante entre $x$ y $u$.

El término:

\begin{equation}
\omega_eL_qi_q
\end{equation}

constituye un ejemplo directo de acoplamiento entre estados.

Asimismo, la expresión general del par contiene:

\begin{equation}
(L_d-L_q)i_di_q,
\end{equation}

que representa la contribución de reluctancia cuando existe saliencia magnética.

En el caso particular $L_d=L_q$, esta última contribución se anula, pero el acoplamiento electromecánico asociado a la velocidad permanece.

\subsection{Interpretación física de un estado}

Un estado no debe entenderse simplemente como una variable almacenada por el programa de simulación. Desde el punto de vista físico, un estado representa una cantidad necesaria para determinar la evolución futura del sistema junto con las entradas y los parámetros.

En la planta PMSM:

\begin{itemize}
    \item $i_d$ y $i_q$ contienen la información dinámica asociada a las corrientes estatóricas en el marco rotante;
    \item $\omega_m$ determina la evolución de la velocidad y participa en la conversión entre frecuencia mecánica y eléctrica;
    \item $\theta_m$ determina la posición instantánea del rotor y la orientación utilizada por las transformaciones de coordenadas.
\end{itemize}

Esta interpretación permite distinguir los estados de otras variables calculadas.

Por ejemplo, $T_e$ es una magnitud fundamental del sistema, pero en el modelo actual se calcula a partir de los estados eléctricos y de los parámetros electromagnéticos. Por ello se considera una variable derivada o intermedia, y no un quinto estado independiente.

\subsection{Condiciones iniciales}

Para ejecutar una simulación es necesario definir las condiciones iniciales del estado:

\begin{equation}
x(0)=
\begin{bmatrix}
i_d(0)\\
i_q(0)\\
\omega_m(0)\\
\theta_m(0)
\end{bmatrix}.
\end{equation}

La implementación base actual inicia los estados internos en cero:

\begin{equation}
i_d(0)=0,
\end{equation}

\begin{equation}
i_q(0)=0,
\end{equation}

\begin{equation}
\omega_m(0)=0,
\end{equation}

\begin{equation}
\theta_m(0)=0.
\end{equation}

Estas condiciones corresponden a una condición inicial de reposo del modelo computacional.

No deben interpretarse como una afirmación acerca del estado inicial de una máquina física durante una prueba determinada. Una prueba experimental o una simulación diferente puede requerir condiciones iniciales distintas.

\subsection{Salidas de la planta}

Las variables entregadas por la planta dependen de la interfaz utilizada. Para la implementación Python actual se define una estructura de salida que contiene:

\begin{equation}
y =
\begin{bmatrix}
i_a &
i_b &
i_c &
i_d &
i_q &
\omega_m &
\theta_m &
T_e
\end{bmatrix}^{T}.
\end{equation}

Las corrientes $i_a$, $i_b$ e $i_c$ permiten interactuar con elementos definidos en coordenadas trifásicas.

Las corrientes $i_d$ e $i_q$ permiten observar directamente el estado electromagnético en el marco rotante.

La velocidad y posición permiten observar la evolución mecánica, mientras que el par electromagnético proporciona una variable fundamental para evaluar el intercambio de energía entre la parte eléctrica y la mecánica.

\subsection{Correspondencia con la implementación Python}

La representación matemática tiene una correspondencia directa con la arquitectura de software desarrollada para la migración Python.

La clase de parámetros:

\begin{equation}
\mathrm{PMSMParameters}
\end{equation}

contiene los parámetros utilizados por la planta.

La clase:

\begin{equation}
\mathrm{PMSMBase}
\end{equation}

implementa la dinámica del modelo M4.

La estructura de entrada:

\begin{equation}
\mathrm{PlantInput}
\end{equation}

transporta las señales que atraviesan la frontera de la planta.

La estructura:

\begin{equation}
\mathrm{PlantOutput}
\end{equation}

transporta las variables generadas por la planta.

La correspondencia conceptual puede resumirse como:

\begin{equation}
\{
x,u,\Theta
\}
\longrightarrow
\mathrm{PMSMBase}
\longrightarrow
y.
\end{equation}

Esta correspondencia permite conservar una relación trazable entre la formulación matemática y el código ejecutable.

\subsection{Transformación entre coordenadas e implementación}

El espacio de estados utiliza principalmente las variables $dq$ porque las ecuaciones dinámicas de la máquina se expresan de forma compacta en este marco.

Sin embargo, la interfaz de la planta puede utilizar variables trifásicas $abc$.

La implementación realiza por tanto la cadena:

\begin{equation}
abc
\rightarrow
\alpha\beta
\rightarrow
dq.
\end{equation}

Para recuperar las variables trifásicas:

\begin{equation}
dq
\rightarrow
\alpha\beta
\rightarrow
abc.
\end{equation}

Las transformaciones no constituyen una planta diferente. Son cambios de representación de las mismas magnitudes físicas.

Esta distinción es importante para evitar confundir el espacio de estados con el sistema de coordenadas utilizado para expresarlo.

\subsection{Integración temporal}

El espacio de estados define las derivadas, pero una implementación digital necesita además un procedimiento para obtener los estados en instantes discretos.

En la implementación base actual se utiliza integración explícita de Euler.

Para un paso temporal $\Delta t$:

\begin{equation}
x_{k+1}
=
x_k+\Delta t\,f(x_k,u_k,\Theta).
\end{equation}

Aplicado individualmente:

\begin{equation}
i_{d,k+1}
=
i_{d,k}
+
\Delta t
\frac{
v_{d,k}-R_si_{d,k}
+\omega_{e,k}L_qi_{q,k}
}{
L_d
}.
\end{equation}

De forma equivalente:

\begin{equation}
i_{q,k+1}
=
i_{q,k}
+
\Delta t
\frac{
v_{q,k}-R_si_{q,k}
-\omega_{e,k}
(L_di_{d,k}+\lambda_m)
}{
L_q
}.
\end{equation}

La velocidad se actualiza mediante:

\begin{equation}
\omega_{m,k+1}
=
\omega_{m,k}
+
\Delta t
\frac{
T_{e,k}-T_{\mathrm{load},k}
}{
J
}.
\end{equation}

Y la posición mediante:

\begin{equation}
\theta_{m,k+1}
=
\theta_{m,k}
+
\Delta t\,\omega_{m,k}.
\end{equation}

Esta formulación hace explícita la relación entre las ecuaciones continuas y la ejecución digital del modelo.

\subsection{Escalas temporales}

La planta PMSM contiene fenómenos con escalas temporales diferentes.

Las corrientes eléctricas pueden presentar variaciones considerablemente más rápidas que la dinámica mecánica de la velocidad. Por esta razón, el paso temporal utilizado para la simulación debe ser compatible con la dinámica eléctrica que se pretende resolver.

En la reconstrucción actual se utiliza un paso base de:

\begin{equation}
\Delta t=50\,\mu\mathrm{s}.
\end{equation}

Este valor también permite mantener una correspondencia con la escala temporal utilizada históricamente en la reconstrucción del lazo de control de corriente.

La selección del paso de integración pertenece a la implementación numérica y no debe confundirse con un parámetro físico intrínseco de la máquina.

\subsection{Relación entre espacio de estados y parámetros trazables}

Una ventaja de esta representación consiste en que cada término de la dinámica puede asociarse con un parámetro concreto.

Por ejemplo:

\begin{equation}
R_s
\rightarrow
\text{caída resistiva},
\end{equation}

\begin{equation}
L_d,L_q
\rightarrow
\text{dinámica de corriente},
\end{equation}

\begin{equation}
\lambda_m
\rightarrow
\text{flujo permanente y FEM},
\end{equation}

\begin{equation}
p
\rightarrow
\text{conversión mecánica-eléctrica},
\end{equation}

\begin{equation}
J
\rightarrow
\text{dinámica mecánica}.
\end{equation}

Esta relación permite conectar la formulación matemática con la matriz de trazabilidad de parámetros establecida para el proyecto.

Los parámetros documentales, derivados y las decisiones de modelado no se mezclan arbitrariamente. Cada uno participa en la formulación con un significado explícito.

\subsection{Nivel de fidelidad correspondiente al modelo M4}

La representación en espacio de estados descrita en este capítulo corresponde al nivel concentrado M4 utilizado actualmente como base reproducible.

En este nivel:

\begin{itemize}
    \item las inductancias se representan mediante valores concentrados;
    \item el enlace de flujo de los imanes se representa mediante $\lambda_m$;
    \item el modelo utiliza una representación analítica del par;
    \item no se incorporan mapas FEM validados;
    \item no se incorporan armónicos espaciales detallados;
    \item no se incorporan pérdidas electromagnéticas dependientes de frecuencia mediante mapas;
    \item no se incorpora todavía una dinámica térmica completa;
    \item los parámetros de fricción viscosa y par de fricción permanecen pendientes de verificación.
\end{itemize}

Por tanto, el espacio de estados M4 debe considerarse una representación física concentrada y trazable, no la representación final de máxima fidelidad de la máquina.

\subsection{Qué cambia cuando aumente la fidelidad}

La estructura de estados no necesariamente debe desaparecer cuando el modelo evolucione.

Una evolución natural consiste en mantener las variables principales y sustituir algunas relaciones internas por funciones de mayor fidelidad.

Por ejemplo, el enlace de flujo podría pasar de:

\begin{equation}
\psi_d=L_di_d+\lambda_m
\end{equation}

a una función:

\begin{equation}
\psi_d
=
\psi_d(i_d,i_q,\theta_m),
\end{equation}

obtenida mediante mapas electromagnéticos.

De manera equivalente:

\begin{equation}
\psi_q
=
\psi_q(i_d,i_q,\theta_m).
\end{equation}

El par podría entonces obtenerse mediante una relación dependiente del estado:

\begin{equation}
T_e
=
T_e(i_d,i_q,\theta_m).
\end{equation}

Esto permitiría incorporar saliencia, saturación magnética, dependencia angular y otros fenómenos cuando exista evidencia suficiente para parametrizarlos.

La transición prevista es:

\begin{equation}
\text{modelo concentrado}
\rightarrow
\text{mapas}
\rightarrow
\text{modelo no lineal}
\rightarrow
\text{FEM/ROM}.
\end{equation}

La estructura de estados constituye así una base sobre la cual pueden evolucionar las funciones físicas internas.

\subsection{Relación con mapas electromagnéticos}

Los mapas electromagnéticos no deben considerarse un sustituto arbitrario del espacio de estados. Su función es proporcionar relaciones físicas más detalladas para aquellas magnitudes que en el modelo concentrado se representan mediante constantes o expresiones simplificadas.

Por ejemplo, un mapa de flujo podría representar:

\begin{equation}
\psi_d=f(i_d,i_q,\theta_e),
\end{equation}

\begin{equation}
\psi_q=g(i_d,i_q,\theta_e).
\end{equation}

El espacio de estados continuaría describiendo la evolución temporal de las corrientes, velocidad y posición, mientras que las funciones electromagnéticas proporcionarían relaciones internas de mayor fidelidad.

Los mapas recuperados existentes en el proyecto se mantienen separados de la implementación M4 activa hasta que su semántica, unidades, ejes, condiciones y procedencia hayan sido verificadas.

\subsection{Relación con FEM y ROM}

Un modelo obtenido mediante FEM puede aportar información sobre la distribución espacial del campo magnético, saturación, dependencia angular, armónicos y otras características que no aparecen explícitamente en el modelo concentrado.

Sin embargo, la incorporación de resultados FEM al Digital Twin requiere una etapa de reducción y organización de la información.

Una ROM puede transformar información de alta fidelidad en relaciones utilizables por una simulación dinámica.

La relación conceptual es:

\begin{equation}
\mathrm{FEM}
\rightarrow
\mathrm{datos}
\rightarrow
\mathrm{ROM/mapas}
\rightarrow
f(x,u,\Theta)
\rightarrow
\text{planta ejecutable}.
\end{equation}

Por ello, el espacio de estados constituye también una frontera natural para integrar progresivamente información electromagnética de mayor fidelidad.

\subsection{Condición particular de la posición del rotor}

La implementación actual presenta una particularidad que debe mantenerse explícita para evitar una interpretación incorrecta del código.

La estructura de entrada de la planta contiene una señal denominada $\theta_m$, mientras que la planta también mantiene internamente un estado de posición mecánica.

Conceptualmente, ambas cantidades representan la misma magnitud física, pero su coexistencia en la implementación actual responde a la reconstrucción incremental del modelo.

Por esta razón, esta interfaz no debe reinterpretarse silenciosamente como una arquitectura definitiva.

La consolidación futura deberá decidir de forma explícita si la posición del rotor es:

\begin{enumerate}
    \item un estado interno integrado por la planta;
    \item una señal externa proporcionada por una fuente de posición;
    \item o una interfaz cuidadosamente definida que permita ambas posibilidades sin duplicar la dinámica.
\end{enumerate}

Esta observación pertenece a la evolución de la implementación y no modifica la formulación física presentada anteriormente.

\subsection{Separación entre planta y controlador}

El espacio de estados pertenece a la planta PMSM.

El controlador FOC utiliza información de la planta para calcular las acciones necesarias para seguir una referencia, pero sus ecuaciones no forman parte del vector de estados de la planta.

De manera conceptual:

\begin{equation}
\text{referencia}
\rightarrow
\text{controlador}
\rightarrow
\text{acción}
\rightarrow
\text{planta}
\rightarrow
\text{realimentación}.
\end{equation}

La planta se encarga de responder físicamente a las entradas.

El FOC, MPC u otra estrategia de control se encarga de calcular acciones sobre la planta.

Esta separación permite que una misma planta pueda ser utilizada posteriormente para comparar diferentes estrategias de control sin modificar arbitrariamente su definición física.

\subsection{Espacio de estados y contratos de arquitectura}

La separación entre planta y control se extiende mediante los contratos de interfaz del sistema.

Un bloque no debe depender directamente de la implementación interna de otro bloque. Las señales atraviesan interfaces definidas que actúan como contratos.

En términos conceptuales:

\begin{equation}
\text{Controlador}
\rightarrow
[\text{contrato}]
\rightarrow
\text{Planta}.
\end{equation}

Y en sentido de realimentación:

\begin{equation}
\text{Planta}
\rightarrow
[\text{contrato}]
\rightarrow
\text{Controlador}.
\end{equation}

El espacio de estados define la dinámica interna de la planta, mientras que el contrato define qué información puede entrar o salir de ella.

Esta distinción permite mantener independientes la física de la máquina y la arquitectura de control.

\subsection{Relación con MATLAB y Simscape}

La formulación de estados también proporciona una referencia para comparar implementaciones realizadas en diferentes plataformas.

MATLAB/Simscape puede representar la máquina mediante componentes físicos y resolver internamente sus ecuaciones. Python puede implementar explícitamente el modelo concentrado utilizado en la reconstrucción.

Aunque las dos implementaciones puedan utilizar mecanismos numéricos diferentes, ambas pueden compararse mediante variables físicas comunes:

\begin{itemize}
    \item corrientes;
    \item velocidad;
    \item posición;
    \item par;
    \item tensiones;
    \item respuesta temporal;
    \item condiciones de operación.
\end{itemize}

Esta comparación permite distinguir entre una diferencia debida al algoritmo numérico y una diferencia debida a la formulación física.

La trazabilidad no exige que ambos entornos tengan exactamente la misma implementación interna. Exige que las hipótesis, parámetros, entradas, salidas y condiciones de comparación estén suficientemente definidas.

\subsection{Relación con el Digital Twin}

El espacio de estados constituye una de las capas fundamentales del Digital Twin PMSM.

La relación conceptual puede expresarse como:

\begin{equation}
\begin{aligned}
\text{máquina física}
&\rightarrow \text{parámetros}
\rightarrow \text{modelo matemático}
\\
&\rightarrow \text{espacio de estados}
\rightarrow \text{modelo ejecutable}.
\end{aligned}
\end{equation}

A partir de esta base, el Digital Twin puede incorporar progresivamente datos experimentales, mapas electromagnéticos, modelos de pérdidas, información térmica y modelos derivados de FEM.

El espacio de estados proporciona entonces la estructura dinámica sobre la cual se integran esos elementos.

No se pretende que el modelo M4 represente desde el inicio todos los fenómenos de la máquina real. Su función es proporcionar una base explícita, reproducible y trazable que pueda evolucionar mediante evidencia adicional.

\subsection{Resumen de la formulación en espacio de estados}

La planta M4 puede resumirse mediante:

\begin{equation}
x=
\begin{bmatrix}
i_d\\
i_q\\
\omega_m\\
\theta_m
\end{bmatrix},
\qquad
u=
\begin{bmatrix}
v_d\\
v_q\\
T_{\mathrm{load}}
\end{bmatrix},
\qquad
\Theta=
\left\{
R_s,L_d,L_q,\lambda_m,p,J
\right\}.
\end{equation}

Su dinámica queda definida por:

\begin{equation}
\dot{x}=f(x,u,\Theta).
\end{equation}

Las principales relaciones físicas son:

\begin{equation}
\omega_e=p\omega_m,
\end{equation}

\begin{equation}
T_e
=
\frac{3}{2}p
\left[
\lambda_mi_q+
(L_d-L_q)i_di_q
\right],
\end{equation}

\begin{equation}
\frac{d\omega_m}{dt}
=
\frac{T_e-T_{\mathrm{load}}}{J},
\end{equation}

\begin{equation}
\frac{d\theta_m}{dt}
=
\omega_m.
\end{equation}

Esta representación concentra en una estructura única la información necesaria para ejecutar la dinámica principal de la planta.

\subsection{Alcance del modelo actual}

El espacio de estados desarrollado en este capítulo debe interpretarse como la representación matemática correspondiente al nivel M4 actualmente establecido.

Su valor principal es proporcionar una base:

\begin{itemize}
    \item físicamente interpretable;
    \item parametrizada;
    \item trazable;
    \item reproducible;
    \item independiente del controlador;
    \item compatible con la reconstrucción Python;
    \item comparable con implementaciones de referencia;
    \item preparada para evolucionar hacia modelos de mayor fidelidad.
\end{itemize}

Los elementos que todavía requieren información adicional no deben introducirse como si fueran parámetros confirmados. Entre ellos se encuentran la fricción detallada, mapas electromagnéticos verificados, pérdidas dependientes de frecuencia, comportamiento térmico completo y relaciones obtenidas de FEM/ROM.

\subsection{Cierre del capítulo}

La representación en espacio de estados transforma la formulación matemática de la PMSM en una estructura directamente utilizable por una implementación computacional.

Los estados $i_d$, $i_q$, $\omega_m$ y $\theta_m$ representan las variables dinámicas principales del modelo M4. Las entradas representan las acciones eléctricas y mecánicas que atraviesan la frontera de la planta, mientras que los parámetros contienen las propiedades utilizadas para determinar su evolución.

Esta estructura permite mantener una separación clara entre la dinámica física de la máquina y las estrategias de control que posteriormente actuarán sobre ella.

La representación también establece una base para aumentar progresivamente la fidelidad del modelo. Las relaciones analíticas actuales pueden ser sustituidas o complementadas por mapas electromagnéticos, modelos no lineales, resultados FEM y ROM sin perder necesariamente la estructura general de la planta.

De esta manera, el espacio de estados constituye un punto de unión entre la formulación matemática, la implementación Python, la referencia MATLAB/Simscape y la evolución futura hacia el Digital Twin PMSM.

El siguiente capítulo documentará la implementación de referencia MATLAB/Simscape, estableciendo su función dentro de la trazabilidad del modelo y diferenciándola de la reconstrucción reproducible desarrollada posteriormente en Python.


\section{Implementación de referencia MATLAB/Simscape}

\subsection{Propósito de la implementación de referencia}

La formulación matemática desarrollada en los capítulos anteriores define la estructura física de la planta PMSM. Sin embargo, para establecer una referencia de implementación resulta necesario disponer de una representación ejecutable que permita observar el comportamiento dinámico de la máquina bajo condiciones de simulación.

En este trabajo, MATLAB y Simscape constituyen una referencia de implementación física y una fuente de evidencia para la reconstrucción progresiva de la planta. Su función no consiste en sustituir la formulación matemática ni en convertirse automáticamente en la especificación permanente del sistema. Su función es proporcionar una realización computacional con la cual pueden contrastarse variables, condiciones de operación, estructuras de conexión y comportamientos dinámicos.

La relación conceptual utilizada es:

\begin{equation}
\begin{aligned}
\text{máquina física}
&\rightarrow \text{documentación}
\rightarrow \text{formulación matemática}
\\
&\rightarrow \text{implementación MATLAB/Simscape}
\rightarrow \text{evidencia}
\rightarrow \text{reconstrucción Python}.
\end{aligned}
\end{equation}

Esta secuencia permite separar tres conceptos que podrían confundirse si se considera únicamente el archivo de simulación:

\begin{itemize}
    \item el modelo físico que se pretende representar;
    \item la implementación concreta utilizada para simularlo;
    \item la evidencia obtenida a partir de dicha implementación.
\end{itemize}

Esta separación es fundamental para conservar la trazabilidad del Digital Twin.

\subsection{Origen de la implementación}

La implementación de referencia se encuentra asociada al trabajo previo de modelado y validación de la máquina PMSM de referencia Siemens 1FT7087.

La existencia de una implementación MATLAB/Simscape permitió disponer de una representación ejecutable antes de iniciar la reconstrucción completa en Python. Esto resulta especialmente importante en un proceso de migración o recuperación, porque proporciona un punto de comparación independiente del nuevo código.

La implementación no debe interpretarse como una fuente única de todos los parámetros. Los parámetros de la máquina deben conservar su propia procedencia documental, derivada o de modelado, tal como se estableció en la caracterización de la máquina.

Por tanto:

\begin{equation}
\text{Simscape}
\neq
\text{fuente automática de parámetros}.
\end{equation}

En cambio:

\begin{equation}
\text{Simscape}
\rightarrow
\text{implementación y evidencia de comportamiento}.
\end{equation}

Esta distinción permite evitar que un valor recuperado de una implementación histórica sea considerado automáticamente como una propiedad confirmada de la máquina.

\subsection{Relación con la máquina Siemens 1FT7087}

La implementación de referencia se interpreta dentro de la genealogía de la máquina Siemens 1FT7087 documentada en el capítulo 02.

La máquina física constituye el origen del problema de modelado. La documentación del fabricante proporciona los datos nominales y las características disponibles. A partir de ellos se establece la parametrización y posteriormente se construye una representación matemática.

MATLAB/Simscape proporciona una plataforma para realizar esa representación mediante componentes de simulación y bloques que permiten conectar dominios eléctricos y mecánicos.

La cadena de trazabilidad puede representarse como:

\begin{equation}
\begin{aligned}
\text{Siemens 1FT7087}
&\rightarrow \text{datos documentales}
\rightarrow \text{parámetros}
\\
&\rightarrow \text{modelo matemático}
\rightarrow \text{Simscape}.
\end{aligned}
\end{equation}

El sentido de esta cadena es importante: el modelo debe estar sustentado en la máquina y en la evidencia disponible, y no al contrario.

\subsection{Representación física de la PMSM}

Una implementación mediante Simscape permite representar la máquina desde una perspectiva física, conectando diferentes dominios y conservando relaciones entre variables eléctricas, electromagnéticas y mecánicas.

Desde el punto de vista conceptual, la planta puede dividirse en:

\begin{enumerate}
    \item dominio eléctrico;
    \item conversión electromecánica;
    \item dominio mecánico;
    \item medición y realimentación;
    \item interfaz con el sistema de control.
\end{enumerate}

La máquina recibe excitación eléctrica y produce una respuesta electromecánica. El par desarrollado modifica la dinámica mecánica, mientras que la posición y velocidad del rotor participan en la evolución de las variables eléctricas.

Esta estructura es coherente con el modelo de estados desarrollado anteriormente.

\subsection{Correspondencia con el modelo matemático}

Aunque Simscape puede resolver internamente las ecuaciones del componente físico, el modelo conceptual sigue siendo el mismo.

Las principales magnitudes consideradas son:

\begin{equation}
i_d,\quad i_q,\quad
\omega_m,\quad
\theta_m,
\end{equation}

junto con las magnitudes eléctricas y mecánicas asociadas.

El espacio de estados descrito anteriormente puede escribirse como:

\begin{equation}
\dot{x}=f(x,u,\Theta).
\end{equation}

Una implementación física de Simscape representa estas relaciones mediante componentes y conexiones. El solver se encarga posteriormente de resolver el sistema resultante.

Por ello, no es necesario que el código Python y Simscape tengan la misma estructura interna para que exista una relación física entre ambos. Lo necesario es que las variables comparadas representen las mismas magnitudes y que las condiciones de simulación estén definidas.

\subsection{Parámetros de la máquina y parámetros de simulación}

Una distinción esencial consiste en separar los parámetros propios de la máquina de las condiciones utilizadas para realizar una simulación.

Entre los parámetros de la máquina se encuentran:

\begin{equation}
R_s,\quad L_d,\quad L_q,\quad
\lambda_m,\quad p,\quad J.
\end{equation}

En cambio, una simulación puede incorporar condiciones tales como:

\begin{equation}
V_{DC},
\quad
T_{\mathrm{load}},
\quad
\omega_{\mathrm{ref}},
\quad
\Delta t,
\quad
t_{\mathrm{stop}}.
\end{equation}

Estas magnitudes no deben mezclarse automáticamente.

En particular, el voltaje del bus DC pertenece a las condiciones eléctricas del accionamiento y de la excitación de la máquina. No constituye una propiedad intrínseca de la PMSM equivalente a su resistencia, inductancia, flujo o inercia.

Por ello, el valor de $V_{DC}$ debe documentarse dentro de la configuración eléctrica del experimento o del accionamiento correspondiente.

\subsection{Condición del bus DC}

La documentación del fabricante relaciona los datos nominales de la máquina con determinadas condiciones de alimentación y accionamiento.

En la documentación de referencia se encuentra una condición nominal asociada a un enlace DC de 600 V. Durante la reconstrucción del proyecto también se identificó una condición de 540 V asociada a la configuración utilizada como base experimental.

Ambos valores deben conservarse con su contexto.

Por tanto, no se debe trasladar un valor de $V_{DC}$ a la parametrización intrínseca de la planta como si fuera una propiedad permanente de la máquina.

La separación conceptual es:

\begin{equation}
\text{PMSM}
+
\text{accionamiento}
+
V_{DC}
\rightarrow
\text{condición de operación}.
\end{equation}

Esta separación resulta especialmente importante al comparar diferentes simulaciones.

\subsection{Dominio mecánico}

La implementación de referencia también permite observar la interacción entre el par electromagnético y la dinámica mecánica.

La relación fundamental utilizada como referencia es:

\begin{equation}
J\frac{d\omega_m}{dt}
=
T_e-T_{\mathrm{load}}.
\end{equation}

La posición se obtiene mediante:

\begin{equation}
\frac{d\theta_m}{dt}
=
\omega_m.
\end{equation}

De esta forma, la respuesta mecánica no constituye un bloque independiente de la dinámica eléctrica. El par electromagnético generado por la máquina depende de las variables eléctricas y, a su vez, la velocidad mecánica modifica la velocidad eléctrica.

La implementación de referencia permite observar esta interacción como parte del comportamiento global del sistema.

\subsection{Transformaciones y variables de referencia}

El modelo utiliza variables trifásicas y variables transformadas para representar el comportamiento eléctrico.

La relación conceptual es:

\begin{equation}
abc
\rightarrow
\alpha\beta
\rightarrow
dq.
\end{equation}

La transformación a $dq$ permite expresar las ecuaciones eléctricas en un marco asociado al rotor.

La posición eléctrica se relaciona con la posición mecánica mediante:

\begin{equation}
\theta_e=p\theta_m.
\end{equation}

De forma equivalente:

\begin{equation}
\omega_e=p\omega_m.
\end{equation}

La implementación de referencia constituye una evidencia de cómo estas variables se integran en el modelo completo, mientras que su formulación matemática permanece definida de forma independiente en los capítulos anteriores.

\subsection{Interacción entre planta y accionamiento}

La PMSM no opera de manera aislada. Durante una simulación, las tensiones aplicadas a sus terminales dependen del sistema de accionamiento utilizado.

Por ello es útil separar conceptualmente:

\begin{equation}
\text{fuente DC}
\rightarrow
\text{convertidor}
\rightarrow
\text{máquina PMSM}.
\end{equation}

El controlador actúa sobre el sistema de accionamiento, mientras que el accionamiento genera las condiciones eléctricas aplicadas a la máquina.

Esta arquitectura permite distinguir entre:

\begin{itemize}
    \item variables de la máquina;
    \item variables del convertidor;
    \item variables del controlador;
    \item condiciones externas de operación.
\end{itemize}

La separación resulta esencial para evitar atribuir a la planta parámetros que en realidad pertenecen al accionamiento o al experimento.

\subsection{Qué aporta Simscape a la trazabilidad}

La implementación de referencia aporta diferentes tipos de evidencia.

En primer lugar, permite observar si las conexiones conceptuales entre los dominios eléctrico y mecánico son coherentes.

En segundo lugar, permite establecer condiciones de simulación y observar respuestas temporales.

En tercer lugar, proporciona una referencia independiente frente a la implementación Python.

Finalmente, permite conservar información sobre la arquitectura histórica del sistema, incluso cuando determinadas partes de esa arquitectura posteriormente sean sustituidas.

La evidencia obtenida debe clasificarse según su naturaleza.

Un resultado de simulación puede constituir evidencia de comportamiento bajo determinadas condiciones, pero no necesariamente evidencia documental de un parámetro físico.

Esta distinción se mantiene durante toda la reconstrucción.

\subsection{El papel histórico de \texttt{skeleton.slx}}

Dentro del material recuperado existe un modelo denominado \texttt{skeleton.slx}. Su importancia para este proyecto es principalmente histórica y arquitectónica.

El archivo permite estudiar cómo se organizaban determinados bloques del sistema, qué conexiones existían entre ellos y qué señales se utilizaban en la arquitectura anterior.

Por esta razón se conserva como referencia de trazabilidad.

Sin embargo, no se considera la especificación permanente de la planta.

La interpretación adoptada es:

\begin{equation}
\texttt{skeleton.slx}
=
\text{referencia histórica/temporal}.
\end{equation}

No:

\begin{equation}
\texttt{skeleton.slx}
=
\text{especificación definitiva}.
\end{equation}

Esta distinción permite aprovechar la información recuperada sin quedar condicionados por decisiones arquitectónicas que pertenecían a una etapa anterior del proyecto.

\subsection{Autopsia de la arquitectura histórica}

El análisis del \texttt{skeleton.slx} permitió identificar una arquitectura organizada alrededor de diferentes módulos y escalas de control.

Entre los elementos recuperados se encuentran referencias a:

\begin{itemize}
    \item una etapa eléctrica;
    \item una etapa de máquina;
    \item control de corriente;
    \item control de velocidad;
    \item transformaciones de coordenadas;
    \item señales de adaptación;
    \item bloques asociados a capas posteriores de la arquitectura.
\end{itemize}

La información recuperada es útil para reconstruir la genealogía del sistema, pero no todos los elementos poseen el mismo nivel de madurez o evidencia.

Por ello, el contenido del modelo histórico debe clasificarse antes de incorporarlo a la implementación actual.

\subsection{Escalas temporales recuperadas}

La autopsia histórica identificó una configuración basada en un paso de integración de:

\begin{equation}
5\times10^{-5}\,\mathrm{s},
\end{equation}

equivalente a una frecuencia de 20 kHz.

También se identificaron escalas diferenciadas para el control de corriente, el control de velocidad y otras capas de la arquitectura.

En particular, la reconstrucción histórica utiliza:

\begin{equation}
T_{s,\mathrm{current}}=50\,\mu\mathrm{s},
\end{equation}

y:

\begin{equation}
T_{s,\mathrm{speed}}=200\,\mu\mathrm{s}.
\end{equation}

Estos valores constituyen evidencia de la configuración histórica del sistema de control. No deben confundirse con constantes físicas de la PMSM.

La reconstrucción actual utiliza $50\,\mu\mathrm{s}$ como paso base del experimento B0, manteniendo esta condición como parte de la configuración experimental y no como parámetro de la planta.

\subsection{Parámetros de control recuperados}

La autopsia también permitió recuperar valores utilizados históricamente en el control de corriente y velocidad.

Entre ellos:

\begin{equation}
K_{p,id}=K_{p,iq}=15.0796,
\end{equation}

\begin{equation}
K_{i,id}=K_{i,iq}=628.3185,
\end{equation}

y para el lazo de velocidad:

\begin{equation}
K_{p,\omega}=5.2804,
\end{equation}

\begin{equation}
K_{i,\omega}=1.2693.
\end{equation}

Estos valores pertenecen al controlador reconstruido y no deben incorporarse al conjunto de parámetros de la planta.

Esta distinción permite mantener separados:

\begin{equation}
\Theta_{\mathrm{plant}}
\neq
\Theta_{\mathrm{FOC}}.
\end{equation}

El primer conjunto describe la máquina, mientras que el segundo describe una implementación particular del controlador.

\subsection{Diferencias de parametrización recuperadas}

La recuperación histórica también puso de manifiesto que no todos los valores utilizados anteriormente coinciden con la parametrización M4 actualmente establecida.

Un ejemplo importante es el enlace de flujo.

La planta M4 utiliza actualmente:

\begin{equation}
\lambda_m=0.107245\,\mathrm{Wb},
\end{equation}

mientras que determinadas implementaciones históricas del FOC utilizaban:

\begin{equation}
\lambda_m=0.1100\,\mathrm{Wb}.
\end{equation}

Estos valores no deben fusionarse silenciosamente.

El primero pertenece a la parametrización M4 actualmente adoptada para la planta, mientras que el segundo corresponde a una configuración histórica del control reconstruido.

La existencia de esta diferencia constituye precisamente un ejemplo de por qué la trazabilidad debe conservar el origen y el contexto de cada parámetro.

\subsection{Límites de corriente recuperados}

Otro ejemplo aparece en el límite de corriente utilizado por el control histórico.

La documentación de la máquina contiene un valor de corriente máxima de 225 A asociado a las condiciones correspondientes del fabricante.

La reconstrucción del controlador, en cambio, contiene un límite de $i_q$ de aproximadamente:

\begin{equation}
|i_q|\leq112.43\,\mathrm{A}.
\end{equation}

Este segundo valor pertenece a la estrategia de control reconstruida y no debe interpretarse como la corriente máxima física de la máquina.

De esta manera:

\begin{equation}
I_{\mathrm{machine,max}}
\neq
I_{\mathrm{FOC,limit}}.
\end{equation}

La diferencia puede ser necesaria para definir una estrategia de operación, pero su significado debe permanecer documentado.

\subsection{Límite de tensión del accionamiento}

La arquitectura histórica también utiliza una limitación de tensión asociada al bus DC:

\begin{equation}
V_{\max}
=
0.95
\frac{V_{DC}}{\sqrt{3}}.
\end{equation}

Esta relación pertenece a la cadena de accionamiento/control y no a la ecuación constitutiva de la máquina.

Por tanto, la planta recibe finalmente las tensiones que resultan del sistema de accionamiento, mientras que el controlador determina sus referencias de tensión respetando los límites disponibles.

Esta separación permite reutilizar la misma planta bajo diferentes condiciones de bus o estrategias de modulación.

\subsection{Elementos que se conservaron como evidencia}

El material MATLAB/Simscape recuperado puede utilizarse para conservar evidencia sobre:

\begin{itemize}
    \item estructura de la arquitectura anterior;
    \item señales y conexiones históricas;
    \item escalas temporales;
    \item parámetros de control;
    \item condiciones eléctricas;
    \item relación entre control y máquina;
    \item resultados de simulaciones anteriores;
    \item decisiones de modelado realizadas durante etapas anteriores.
\end{itemize}

Estos elementos se almacenan conceptualmente como información de trazabilidad.

Su presencia en el material histórico no implica que todos ellos deban permanecer activos en la arquitectura futura.

\subsection{Elementos que no deben convertirse automáticamente en especificación}

La recuperación también exige reconocer qué información no puede elevarse automáticamente a la categoría de requisito.

Entre estos elementos se encuentran:

\begin{itemize}
    \item valores cuya procedencia no está confirmada;
    \item parámetros de simulaciones históricas utilizados fuera de su contexto;
    \item valores del controlador confundidos con parámetros de la planta;
    \item bloques placeholder;
    \item señales cuyo significado no haya sido establecido;
    \item estructuras del \texttt{skeleton.slx} que pertenecían a una etapa temporal;
    \item valores de modelos anteriores que contradicen la fuente M4 actualmente establecida.
\end{itemize}

Esta clasificación permite aprovechar la recuperación sin introducir errores de trazabilidad en la nueva implementación.

\subsection{Relación entre Simscape y la reconstrucción Python}

La migración a Python no pretende reproducir literalmente la estructura interna de Simscape.

El objetivo es reconstruir de forma explícita y reproducible el comportamiento del modelo que puede justificarse mediante la evidencia disponible.

La relación conceptual es:

\begin{equation}
\text{Simscape}
\rightarrow
\text{evidencia}
\rightarrow
\text{formulación}
\rightarrow
\text{Python}.
\end{equation}

Python permite hacer explícitas las ecuaciones, los parámetros y las decisiones numéricas que en un entorno de simulación física pueden estar encapsuladas en bloques.

Esto facilita:

\begin{itemize}
    \item control de versiones;
    \item experimentos reproducibles;
    \item generación automática de datos;
    \item construcción de mapas;
    \item análisis numérico;
    \item comparación de modelos;
    \item preparación de futuras capas inteligentes.
\end{itemize}

La migración no significa eliminar el valor de MATLAB/Simscape. Significa separar su papel de referencia del papel de desarrollo reproducible que desempeña Python.

\subsection{Condiciones de comparación entre implementaciones}

Para comparar una simulación MATLAB/Simscape con una simulación Python no basta con observar si las curvas tienen una forma similar.

La comparación debe especificar:

\begin{itemize}
    \item parámetros de la máquina;
    \item condición de $V_{DC}$;
    \item carga aplicada;
    \item condiciones iniciales;
    \item referencia de velocidad;
    \item paso temporal;
    \item duración;
    \item estrategia de control;
    \item señales utilizadas para la comparación.
\end{itemize}

De esta forma, una comparación puede distinguir entre diferencias debidas a la planta y diferencias debidas al experimento.

Una condición de comparación bien definida puede expresarse conceptualmente como:

\begin{equation}
C=
\{
\Theta,
V_{DC},
T_{\mathrm{load}},
x_0,
\Delta t,
t_{\mathrm{stop}},
\text{control}
\}.
\end{equation}

El conjunto $C$ forma parte de la trazabilidad del experimento.

\subsection{Papel de MATLAB/Simscape en la validación}

La implementación de referencia permite realizar una validación por capas.

En una primera capa se verifica la coherencia física general.

En una segunda capa se comparan variables dinámicas entre implementaciones.

En una tercera capa se contrastan condiciones nominales y documentación de fabricante.

En una cuarta capa, conforme aumente la fidelidad, se podrán comparar resultados frente a mapas electromagnéticos y posteriormente frente a modelos FEM o ROM.

La validación debe ser incremental.

No se considera suficiente que una simulación reproduzca una respuesta visualmente plausible. Cada nivel de fidelidad debe estar asociado a evidencia adecuada.

\subsection{Relación con mapas electromagnéticos}

La implementación de referencia constituye una base para identificar qué relaciones físicas deberán conservarse cuando el modelo evolucione hacia mapas.

En el modelo concentrado:

\begin{equation}
\psi_d=L_di_d+\lambda_m,
\end{equation}

\begin{equation}
\psi_q=L_qi_q.
\end{equation}

En un modelo basado en mapas, estas relaciones pueden transformarse en funciones dependientes de corriente y posición:

\begin{equation}
\psi_d=
f(i_d,i_q,\theta_e),
\end{equation}

\begin{equation}
\psi_q=
g(i_d,i_q,\theta_e).
\end{equation}

La implementación MATLAB/Simscape puede servir como referencia de comportamiento para comprobar que la sustitución de estas relaciones no introduce inconsistencias en la dinámica global.

Sin embargo, los mapas deberán estar respaldados por datos propios o por evidencia cuya procedencia haya sido verificada.

\subsection{Relación con FEM y ROM}

La evolución hacia FEM y ROM representa un cambio de fidelidad, no un cambio de identidad de la planta.

La misma máquina de referencia puede representarse mediante diferentes niveles:

\begin{equation}
\mathrm{M4}
\rightarrow
\mathrm{maps}
\rightarrow
\mathrm{nonlinear}
\rightarrow
\mathrm{FEM}
\rightarrow
\mathrm{ROM}.
\end{equation}

MATLAB/Simscape puede permanecer como plataforma de comparación durante esta evolución.

Python puede utilizar los resultados para construir conjuntos de datos, mapas y modelos reducidos.

Posteriormente, los modelos evolucionados pueden volver a integrarse en Simulink/Simscape para realizar comparaciones cruzadas.

El flujo no es una sustitución unilateral de una plataforma por otra, sino un proceso de transferencia de conocimiento y evidencia.

\subsection{Separación entre referencia y especificación}

Una implementación de referencia puede contener decisiones que fueron adecuadas para una determinada etapa del proyecto, pero que posteriormente deban ser modificadas.

Por esta razón, el presente trabajo distingue entre:

\begin{equation}
\text{referencia histórica}
\end{equation}

y:

\begin{equation}
\text{especificación actual}.
\end{equation}

La especificación actual se construye a partir de la evidencia disponible, la formulación matemática, la trazabilidad de parámetros y las decisiones explícitamente adoptadas para la nueva arquitectura.

El material histórico se conserva para poder explicar cómo se llegó a dicha arquitectura.

Esta separación es particularmente importante en un proceso de recuperación de un proyecto previo.

\subsection{Relación con la arquitectura de control}

La implementación de referencia también permite observar la frontera entre la planta y el control.

Conceptualmente:

\begin{equation}
\text{referencia de velocidad}
\rightarrow
\text{FOC}
\rightarrow
\text{accionamiento}
\rightarrow
\text{PMSM}.
\end{equation}

La planta genera:

\begin{equation}
\text{PMSM}
\rightarrow
\{i_a,i_b,i_c,i_d,i_q,\omega_m,\theta_m,T_e\}.
\end{equation}

Estas variables pueden regresar mediante las interfaces definidas hacia el controlador.

La existencia de una implementación de control dentro del modelo histórico no convierte al controlador en parte de la definición física de la PMSM.

Esta separación será necesaria posteriormente para evaluar FOC, MPC, ANFIS y otras capas inteligentes utilizando una planta común.

\subsection{Preparación para la arquitectura híbrida}

El papel de MATLAB/Simscape adquiere especial relevancia cuando se considera la arquitectura futura de control híbrido.

El objetivo no consiste en construir una planta diferente para cada estrategia de control.

La planta debe constituir una referencia común:

\begin{equation}
\text{PMSM}
\rightarrow
\left\{
\begin{array}{c}
\text{FOC}\\
\text{MPC}\\
\text{ANFIS}\\
\text{TinyML}\\
\text{arquitectura híbrida}
\end{array}
\right.
\end{equation}

Las diferentes estrategias pueden compartir la misma representación física siempre que respeten los contratos de interfaz.

Esta condición es fundamental para que una comparación entre estrategias de control tenga significado.

\subsection{Contratos de frontera}

La arquitectura actual establece que los bloques deben comunicarse mediante contratos o buses de interfaz.

La relación conceptual es:

\begin{equation}
\text{controlador}
\rightarrow
[\text{contrato}]
\rightarrow
\text{accionamiento/planta}.
\end{equation}

Y en la realimentación:

\begin{equation}
\text{planta}
\rightarrow
[\text{contrato}]
\rightarrow
\text{controlador}.
\end{equation}

Los buses actúan como mecanismos de transporte y traducción de esos contratos.

La implementación histórica puede mostrar determinadas conexiones físicas entre bloques, pero la arquitectura futura se define a partir de contratos explícitos y no únicamente a partir de la topología de un archivo antiguo.

\subsection{Valor de la referencia MATLAB/Simscape para el Digital Twin}

La implementación de referencia forma parte de la genealogía del Digital Twin porque permite conectar la máquina física con una representación computacional ejecutable.

Sin embargo, el Digital Twin no se define por el archivo de Simscape.

Su construcción progresiva puede expresarse como:

\begin{equation}
\begin{aligned}
\text{máquina física}
&\rightarrow \text{documentación}
\rightarrow \text{parámetros}
\\
&\rightarrow \text{modelo matemático}
\rightarrow \text{referencia Simscape}
\\
&\rightarrow \text{datos}
\rightarrow \text{modelos de mayor fidelidad}
\rightarrow \text{Digital Twin}.
\end{aligned}
\end{equation}

En esta perspectiva, MATLAB/Simscape constituye una etapa importante de implementación y validación, pero no representa el final de la evolución del modelo.

\subsection{Limitaciones de la implementación de referencia}

La implementación de referencia no debe utilizarse como evidencia suficiente para afirmar que todos los fenómenos de la máquina real están representados.

Entre las limitaciones que deben mantenerse visibles se encuentran:

\begin{itemize}
    \item la parametrización puede corresponder a un nivel concentrado;
    \item algunos parámetros recuperados pertenecen a configuraciones históricas;
    \item no toda la información del modelo histórico tiene procedencia documental confirmada;
    \item el comportamiento electromagnético detallado requiere mapas o modelos de mayor fidelidad;
    \item la dependencia térmica completa no está representada en el modelo M4;
    \item la fricción y otras pérdidas requieren información adicional;
    \item la equivalencia exacta entre determinadas variantes documentales debe verificarse.
\end{itemize}

Estas limitaciones no invalidan el uso de la referencia. Determinan qué afirmaciones pueden realizarse a partir de ella.

\subsection{Función dentro de la estrategia de migración}

La implementación MATLAB/Simscape cumple cuatro funciones principales dentro del proceso de migración:

\begin{enumerate}
    \item conservar una referencia computacional del comportamiento de la planta;
    \item aportar evidencia para reconstruir relaciones y condiciones;
    \item permitir comparaciones independientes con la implementación Python;
    \item conservar la genealogía de la arquitectura anterior.
\end{enumerate}

Python asume posteriormente una función diferente: hacer explícitas las ecuaciones, organizar los datos, ejecutar experimentos reproducibles y preparar la evolución hacia modelos de mayor fidelidad.

Por tanto:

\begin{equation}
\text{MATLAB/Simscape}
\rightarrow
\text{referencia y evidencia},
\end{equation}

mientras:

\begin{equation}
\text{Python}
\rightarrow
\text{reconstrucción y evolución reproducible}.
\end{equation}

Ambas funciones son complementarias.

\subsection{Cierre del capítulo}

La implementación MATLAB/Simscape constituye una etapa fundamental en la genealogía computacional de la planta PMSM. Su valor principal reside en proporcionar una referencia ejecutable, una fuente de evidencia y un punto de comparación para la reconstrucción posterior.

La implementación debe interpretarse siempre dentro de la trazabilidad de la máquina Siemens 1FT7087 y de la formulación matemática establecida previamente. Los parámetros de la planta deben conservar su propia procedencia, mientras que las condiciones del accionamiento, los parámetros del controlador y las condiciones experimentales deben mantenerse separados.

El archivo \texttt{skeleton.slx} se conserva como referencia histórica y temporal de arquitectura y trazabilidad. Su información resulta útil para comprender la evolución del proyecto, pero no se adopta como especificación permanente.

Esta separación permite pasar de una implementación histórica a una reconstrucción reproducible sin perder la evidencia acumulada.

El siguiente paso consiste en documentar la reconstrucción Python, donde las ecuaciones y decisiones de modelado dejan de depender de la estructura interna de los bloques de Simscape y pasan a constituir una implementación explícita, versionable y reproducible de la planta.


\section{Reconstrucción reproducible en Python}

\subsection{Propósito de la reconstrucción}

La reconstrucción en Python constituye la siguiente etapa de la genealogía computacional de la planta PMSM. Después de establecer la máquina de referencia, su parametrización, la formulación matemática y la implementación MATLAB/Simscape, se desarrolla una representación explícita, versionable y reproducible mediante código Python.

El objetivo no es reproducir literalmente la estructura interna de MATLAB/Simscape. La finalidad es expresar de manera explícita las relaciones físicas que pueden justificarse mediante la documentación, la formulación matemática y la evidencia recuperada del proyecto.

La reconstrucción puede representarse como:

\begin{equation}
\text{referencia MATLAB/Simscape}
\rightarrow
\text{formulación física}
\rightarrow
\text{Python explícito}.
\end{equation}

Esta estrategia permite separar la evidencia histórica de la implementación actualmente mantenida.

La implementación Python constituye además la semilla computacional sobre la cual se desarrollarán posteriormente los mapas electromagnéticos, los modelos de mayor fidelidad y las diferentes estrategias de control que utilizarán la planta como plataforma común.

\subsection{Motivación de la migración a Python}

La utilización de Python proporciona características especialmente útiles para una reconstrucción reproducible:

\begin{itemize}
    \item código fuente versionable;
    \item parametrización explícita;
    \item experimentos automatizables;
    \item generación reproducible de datos;
    \item integración con herramientas de análisis numérico;
    \item construcción sistemática de mapas;
    \item generación automática de figuras y tablas;
    \item trazabilidad entre código, parámetros y resultados.
\end{itemize}

La migración permite, además, desacoplar la definición de la planta de una determinada herramienta gráfica.

La planta puede ejecutarse desde scripts y módulos Python, facilitando que diferentes experimentos utilicen exactamente la misma implementación base.

La decisión de utilizar Python no implica eliminar MATLAB/Simscape. Ambas plataformas cumplen funciones diferentes dentro de la arquitectura del proyecto.

\subsection{Relación entre referencia y reconstrucción}

La relación entre ambas implementaciones debe entenderse en términos de evidencia y formulación, no de copia estructural.

La referencia MATLAB/Simscape aporta información sobre:

\begin{itemize}
    \item representación física;
    \item conexiones entre dominios;
    \item condiciones de operación;
    \item variables observables;
    \item arquitectura histórica;
    \item resultados de simulación.
\end{itemize}

La reconstrucción Python convierte la información físicamente justificable en una implementación explícita.

Por tanto:

\begin{equation}
\text{Simscape}
\neq
\text{código fuente de Python}.
\end{equation}

La relación correcta es:

\begin{equation}
\text{evidencia}
+
\text{formulación}
+
\text{parametrización}
\rightarrow
\text{implementación Python}.
\end{equation}

Esta distinción evita atribuir a Python comportamientos que todavía no hayan sido demostrados o trazados.

\subsection{Arquitectura del código}

La implementación activa se organiza alrededor del paquete:

\begin{equation}
\texttt{pmsm\_hybrid\_control}.
\end{equation}

Dentro de esta estructura, la planta se separa de las capas de control.

La cadena principal de la planta es:

\begin{equation}
\texttt{parameters.py}
\rightarrow
\texttt{PMSMParameters}
\rightarrow
\texttt{PMSMBase}.
\end{equation}

Las interfaces de entrada y salida se representan mediante:

\begin{equation}
\texttt{PlantInput}
\rightarrow
\texttt{PMSMBase}
\rightarrow
\texttt{PlantOutput}.
\end{equation}

Esta separación permite que la planta pueda ser utilizada por diferentes experimentos sin incorporar directamente la lógica del controlador.

\subsection{Fuente de parámetros}

El módulo \texttt{parameters.py} establece la fuente de verdad para la parametrización M4 utilizada por la planta actual.

El propio módulo documenta que la fuente principal es:

\begin{equation}
\texttt{01\_M4\_HARD\_DATA.md}.
\end{equation}

Los parámetros se acompañan de información de trazabilidad.

La estructura utilizada permite asociar a un valor:

\begin{equation}
\{
\text{valor},
\text{estado},
\text{fuente},
\text{unidad},
\text{nota}
\}.
\end{equation}

Esta representación evita que un número aislado pierda su procedencia al ser utilizado en diferentes partes del proyecto.

\subsection{Estados de trazabilidad}

La implementación utiliza estados explícitos para distinguir el nivel de evidencia de los parámetros.

Entre ellos se encuentran:

\begin{itemize}
    \item \texttt{CONFIRMED};
    \item \texttt{PROPOSED};
    \item \texttt{PENDING\_VERIFICATION}.
\end{itemize}

Estos estados no representan niveles de calidad numérica, sino estados de conocimiento y procedencia.

Un valor confirmado puede utilizarse como parte de la parametrización establecida.

Un valor propuesto representa una decisión de modelado que todavía puede evolucionar.

Un valor pendiente no debe incorporarse silenciosamente como si fuera un parámetro físicamente confirmado.

\subsection{Parámetros M4 de la planta}

La parametrización concentrada actualmente utilizada contiene:

\begin{equation}
R_s=0.05\,\Omega,
\end{equation}

\begin{equation}
L_d=1.2\times10^{-3}\,\mathrm{H},
\end{equation}

\begin{equation}
L_q=1.2\times10^{-3}\,\mathrm{H},
\end{equation}

\begin{equation}
\lambda_m=0.107245\,\mathrm{Wb},
\end{equation}

\begin{equation}
p=4,
\end{equation}

y:

\begin{equation}
J=0.00416\,\mathrm{kg\,m^2}.
\end{equation}

Estos valores constituyen la parametrización M4 actualmente adoptada.

La igualdad:

\begin{equation}
L_d=L_q
\end{equation}

corresponde a la representación concentrada utilizada actualmente. No debe interpretarse como una afirmación universal de que la máquina real carezca de cualquier anisotropía.

\subsection{Flujo magnético permanente}

El valor utilizado actualmente para el enlace de flujo permanente es:

\begin{equation}
\lambda_m=0.107245\,\mathrm{Wb}.
\end{equation}

Este valor se obtuvo a partir de la constante de par documentada y del número de pares de polos:

\begin{equation}
\lambda_m=
\frac{k_T}
{(3/2)\sqrt{2}\,p}.
\end{equation}

Utilizando:

\begin{equation}
k_T=0.91\,\mathrm{N\,m/A}
\end{equation}

y:

\begin{equation}
p=4,
\end{equation}

se obtiene aproximadamente el valor utilizado por M4.

La procedencia numérica y la procedencia documental deben mantenerse diferenciadas. El valor se encuentra incorporado a la parametrización M4, mientras que su trazabilidad documental detallada permanece identificada para futura verificación.

\subsection{Separación de parámetros históricos}

Durante la reconstrucción se encontraron valores pertenecientes a implementaciones anteriores.

Un caso relevante es:

\begin{equation}
\lambda_m=0.1100\,\mathrm{Wb}.
\end{equation}

Este valor aparece en configuraciones históricas del FOC reconstruido.

No se utiliza como sustituto del valor M4 de la planta.

La separación se expresa como:

\begin{equation}
\lambda_{m,\mathrm{plant}}
\neq
\lambda_{m,\mathrm{FOC,historical}}.
\end{equation}

La diferencia no se corrige silenciosamente porque forma parte de la trazabilidad del proceso de recuperación.

La misma regla se aplica a otros parámetros históricos.

\subsection{Parámetros mecánicos}

La inercia actualmente utilizada es:

\begin{equation}
J=0.00416\,\mathrm{kg\,m^2}.
\end{equation}

Este valor corresponde a la configuración con freno considerada en la parametrización M4.

La documentación del fabricante también proporciona una inercia de aproximadamente:

\begin{equation}
J=0.00274\,\mathrm{kg\,m^2}
\end{equation}

para la configuración sin freno.

La selección del valor M4 debe permanecer documentada porque la inercia depende de la configuración mecánica considerada.

\subsection{Parámetros todavía pendientes}

La reconstrucción no incorpora como parámetros activos de M4 determinados valores históricos de fricción.

En particular:

\begin{equation}
B
\end{equation}

y:

\begin{equation}
T_{\mathrm{fric}}
\end{equation}

permanecen pendientes de verificación.

Valores históricos asociados a CloudE no se trasladan automáticamente a la planta actual.

Esta decisión permite mantener una separación clara entre:

\begin{equation}
\text{parámetro confirmado}
\end{equation}

y:

\begin{equation}
\text{valor histórico recuperado}.
\end{equation}

\subsection{Clase \texttt{PMSMParameters}}

La clase \texttt{PMSMParameters} concentra los parámetros necesarios para ejecutar la planta M4.

Su función es evitar que las ecuaciones dependan de valores dispersos en diferentes módulos.

Conceptualmente:

\begin{equation}
\Theta=
\{
R_s,
L_d,
L_q,
\lambda_m,
p,
J
\}.
\end{equation}

El objeto de parámetros representa este conjunto como una unidad.

La planta recibe así una configuración explícita y verificable.

Esta estructura facilita la creación de diferentes configuraciones futuras sin modificar las ecuaciones fundamentales.

\subsection{Validación de parámetros}

La implementación incorpora validaciones básicas antes de ejecutar la planta.

Entre las condiciones comprobadas se encuentran:

\begin{itemize}
    \item resistencia positiva;
    \item inductancias positivas;
    \item flujo permanente positivo;
    \item inercia positiva;
    \item número de pares de polos consistente;
    \item valores numéricos válidos.
\end{itemize}

Estas comprobaciones no constituyen una validación física completa de la máquina.

Su finalidad es impedir que errores elementales de parametrización pasen inadvertidos durante la ejecución.

\subsection{Estados de la planta}

La clase \texttt{PMSMBase} utiliza cuatro estados principales:

\begin{equation}
x=
\begin{bmatrix}
i_d&
i_q&
\omega_m&
\theta_m
\end{bmatrix}^{T}.
\end{equation}

Los dos primeros representan las corrientes en el marco $dq$.

El tercero representa la velocidad mecánica del rotor.

El cuarto representa la posición mecánica.

Estos estados corresponden directamente a la formulación matemática desarrollada anteriormente.

\subsection{Entradas de la planta}

La interfaz \texttt{PlantInput} representa las magnitudes que la implementación actual utiliza para avanzar la planta.

Incluye:

\begin{equation}
v_a,\quad v_b,\quad v_c,
\end{equation}

el par de carga:

\begin{equation}
T_{\mathrm{load}},
\end{equation}

y la posición mecánica utilizada por la implementación:

\begin{equation}
\theta_m.
\end{equation}

La presencia de $\theta_m$ como entrada debe entenderse como una característica de la implementación actual y no como una afirmación de que la posición del rotor sea una entrada física independiente de la máquina.

La consolidación futura de esta interfaz deberá distinguir de forma más estricta entre estado interno y señal externa de posición.

\subsection{Salidas de la planta}

La interfaz \texttt{PlantOutput} expone:

\begin{equation}
i_a,\quad i_b,\quad i_c,
\end{equation}

\begin{equation}
i_d,\quad i_q,
\end{equation}

\begin{equation}
\omega_m,\quad \theta_m,
\end{equation}

y:

\begin{equation}
T_e.
\end{equation}

Esta interfaz permite que la misma planta sea utilizada por diferentes capas de control y análisis.

La salida no contiene directamente parámetros del controlador.

La separación permite utilizar la planta como componente independiente.

\subsection{Transformación de variables trifásicas}

La implementación recibe variables trifásicas y las transforma al marco de referencia utilizado por las ecuaciones.

La secuencia conceptual es:

\begin{equation}
abc
\rightarrow
\alpha\beta
\rightarrow
dq.
\end{equation}

La implementación actual utiliza una transformación de Clarke basada en:

\begin{equation}
i_\alpha=i_a,
\end{equation}

y:

\begin{equation}
i_\beta=
\frac{i_b-i_c}{\sqrt{3}}.
\end{equation}

Posteriormente se aplica la transformación de Park utilizando la posición eléctrica.

La transformación inversa permite reconstruir las variables trifásicas de salida.

\subsection{Posición eléctrica}

La posición eléctrica se obtiene mediante:

\begin{equation}
\theta_e=p\theta_m.
\end{equation}

De forma equivalente, la velocidad eléctrica es:

\begin{equation}
\omega_e=p\omega_m.
\end{equation}

Esta relación constituye uno de los principales acoplamientos entre los dominios mecánico y eléctrico.

El número utilizado actualmente es:

\begin{equation}
p=4,
\end{equation}

correspondiente a ocho polos físicos.

\subsection{Ecuaciones eléctricas implementadas}

La dinámica eléctrica utilizada por \texttt{PMSMBase} es:

\begin{equation}
\frac{di_d}{dt}
=
\frac{
v_d-R_si_d+\omega_eL_qi_q
}
{L_d},
\end{equation}

y:

\begin{equation}
\frac{di_q}{dt}
=
\frac{
v_q-R_si_q-\omega_e(L_di_d+\lambda_m)
}
{L_q}.
\end{equation}

Estas ecuaciones corresponden directamente a la formulación matemática establecida para el modelo concentrado.

La implementación no introduce en esta etapa dependencias angulares de inductancia ni mapas electromagnéticos.

\subsection{Par electromagnético}

El par electromagnético utilizado como relación de referencia es:

\begin{equation}
T_e=
\frac{3}{2}p
\left[
\lambda_mi_q+
(L_d-L_q)i_di_q
\right].
\end{equation}

Dado que el modelo M4 utiliza:

\begin{equation}
L_d=L_q,
\end{equation}

la expresión se reduce a:

\begin{equation}
T_e=
\frac{3}{2}p\lambda_mi_q.
\end{equation}

Esta ecuación constituye la relación de par utilizada por la implementación M4. No se presenta como una reproducción literal de una ley interna de Simscape.

\subsection{Dinámica mecánica}

La velocidad mecánica se integra mediante:

\begin{equation}
\frac{d\omega_m}{dt}
=
\frac{T_e-T_{\mathrm{load}}}{J}.
\end{equation}

La posición se obtiene mediante:

\begin{equation}
\frac{d\theta_m}{dt}
=
\omega_m.
\end{equation}

El modelo actual no introduce explícitamente términos separados de fricción viscosa o par de fricción porque dichos parámetros permanecen pendientes de verificación.

Por tanto, $T_{\mathrm{load}}$ representa el par resistente externo utilizado en el experimento.

\subsection{Integración numérica}

La implementación actual utiliza integración explícita de Euler.

Para un estado genérico:

\begin{equation}
x_{k+1}
=
x_k+\Delta t\,\dot{x}_k.
\end{equation}

Aplicado a las corrientes:

\begin{equation}
i_{d,k+1}
=
i_{d,k}
+
\Delta t
\frac{di_d}{dt},
\end{equation}

\begin{equation}
i_{q,k+1}
=
i_{q,k}
+
\Delta t
\frac{di_q}{dt}.
\end{equation}

Para la velocidad:

\begin{equation}
\omega_{m,k+1}
=
\omega_{m,k}
+
\Delta t
\frac{T_e-T_{\mathrm{load}}}{J}.
\end{equation}

La posición se actualiza mediante:

\begin{equation}
\theta_{m,k+1}
=
\theta_{m,k}
+
\Delta t\,\omega_{m,k}.
\end{equation}

El método se mantiene explícito porque facilita la inspección del comportamiento numérico y la reproducibilidad del experimento.

\subsection{Paso temporal de referencia}

El experimento base utiliza:

\begin{equation}
\Delta t=50\,\mu\mathrm{s}.
\end{equation}

Esto corresponde a una frecuencia de integración de:

\begin{equation}
f_s=20\,\mathrm{kHz}.
\end{equation}

El paso temporal constituye una condición numérica del experimento y no un parámetro físico de la máquina.

Esta distinción permite cambiar el integrador o realizar estudios de convergencia sin alterar la parametrización física.

\subsection{Condiciones iniciales}

La implementación base inicializa las variables de estado en cero.

Conceptualmente:

\begin{equation}
i_d(0)=0,
\end{equation}

\begin{equation}
i_q(0)=0,
\end{equation}

\begin{equation}
\omega_m(0)=0,
\end{equation}

y:

\begin{equation}
\theta_m(0)=0.
\end{equation}

Estas condiciones corresponden al experimento de arranque utilizado como referencia.

Una condición inicial diferente constituye una configuración experimental diferente y debe quedar registrada junto con los resultados.

\subsection{Peculiaridad actual de \texorpdfstring{$\theta_m$}{theta m}}

La implementación actual presenta una característica que debe conservarse explícitamente en la trazabilidad.

\texttt{PlantInput} contiene $\theta_m$, mientras que \texttt{PMSMBase} también mantiene una variable interna para la posición.

Durante el paso de simulación, la posición recibida se utiliza para calcular la posición eléctrica y se asigna al estado interno antes de la actualización dinámica.

Esto implica que actualmente existe una mezcla entre:

\begin{equation}
\text{posición externa}
\end{equation}

y:

\begin{equation}
\text{posición integrada internamente}.
\end{equation}

La condición no se modifica silenciosamente porque forma parte del estado actual del código.

Debe considerarse una tarea de consolidación futura de la interfaz de la planta.

\subsection{Ausencia de un solver FEM}

La implementación \texttt{PMSMBase} no constituye un solver de elementos finitos.

Su nivel de fidelidad corresponde a un modelo concentrado basado en parámetros.

La estructura es:

\begin{equation}
\text{PMSMBase}
=
\text{modelo matemático parametrizado}.
\end{equation}

No debe afirmarse que la implementación actual reproduce por sí misma la distribución espacial del campo magnético.

La dependencia angular detallada, las saturaciones locales y los armónicos espaciales requieren una representación de mayor fidelidad.

\subsection{Ausencia de mapas electromagnéticos activos}

El repositorio contiene módulos destinados a una futura representación mediante mapas.

Sin embargo, la configuración M4 activa utiliza relaciones analíticas.

El flujo de referencia es:

\begin{equation}
\psi_d=L_di_d+\lambda_m,
\end{equation}

\begin{equation}
\psi_q=L_qi_q.
\end{equation}

Las tablas recuperadas de etapas anteriores permanecen separadas y marcadas como información recuperada pendiente de verificación semántica.

Por tanto, no se consideran automáticamente mapas validados de la máquina.

\subsection{Evolución hacia mapas}

Cuando se disponga de datos electromagnéticos suficientemente trazables, las relaciones analíticas podrán sustituirse o complementarse mediante funciones dependientes del punto de operación.

Una representación futura puede adoptar la forma:

\begin{equation}
\psi_d=
f(i_d,i_q,\theta_e),
\end{equation}

\begin{equation}
\psi_q=
g(i_d,i_q,\theta_e).
\end{equation}

De forma correspondiente, el par podrá evaluarse directamente desde el estado electromagnético.

La arquitectura actual permite introducir esta evolución sin redefinir la identidad de la planta.

\subsection{Evolución del nivel de fidelidad}

La reconstrucción Python se organiza alrededor de una evolución progresiva:

\begin{equation}
\mathrm{M4}
\rightarrow
\mathrm{mapas}
\rightarrow
\mathrm{modelo\ no\ lineal}
\rightarrow
\mathrm{FEM/ROM}
\rightarrow
\mathrm{Digital\ Twin}.
\end{equation}

Cada etapa debe incorporar nueva evidencia.

La etapa M4 proporciona una semilla reproducible.

Los mapas permitirán incorporar dependencias no representadas por el modelo concentrado.

El modelo no lineal podrá incorporar saturación y dependencia angular.

FEM y ROM podrán proporcionar información espacial o modelos reducidos derivados de representaciones de mayor fidelidad.

\subsection{Separación entre planta y controlador}

La reconstrucción Python mantiene separada la planta de las estrategias de control.

La planta responde a entradas eléctricas y mecánicas y produce variables físicas.

El FOC utiliza dichas variables para generar las referencias necesarias para el accionamiento.

La relación conceptual es:

\begin{equation}
\text{FOC}
\rightarrow
\text{accionamiento}
\rightarrow
\text{PMSMBase}.
\end{equation}

Y en realimentación:

\begin{equation}
\text{PMSMBase}
\rightarrow
\text{mediciones}
\rightarrow
\text{FOC}.
\end{equation}

Esta separación permite utilizar posteriormente la misma planta con otras estrategias.

\subsection{Separación de \texttt{foc} y \texttt{foc\_reconstructed}}

El repositorio contiene dos estructuras relacionadas con FOC.

El directorio:

\begin{equation}
\texttt{foc/}
\end{equation}

representa actualmente una estructura de interfaz o esqueleto pendiente de consolidación.

El directorio:

\begin{equation}
\texttt{foc\_reconstructed/}
\end{equation}

contiene el FOC funcional utilizado por los experimentos de integración actuales.

Esta distinción evita presentar una interfaz placeholder como si fuera el controlador funcional utilizado en los resultados.

\subsection{Parámetros del FOC reconstruido}

El FOC reconstruido utiliza parámetros que deben mantenerse separados de la planta.

Entre los valores históricos recuperados se encuentran:

\begin{equation}
K_{p,id}=K_{p,iq}=15.0796,
\end{equation}

\begin{equation}
K_{i,id}=K_{i,iq}=628.3185,
\end{equation}

\begin{equation}
K_{p,\omega}=5.2804,
\end{equation}

y:

\begin{equation}
K_{i,\omega}=1.2693.
\end{equation}

Estos valores describen una configuración particular del controlador.

No deben incorporarse a la clase de parámetros físicos de la PMSM.

\subsection{Desacoplamiento del FOC}

El FOC reconstruido incorpora términos de desacoplamiento basados en la formulación $dq$.

Los términos utilizados como referencia son:

\begin{equation}
v_{d,\mathrm{dec}}
=
-\omega_eL_qi_q,
\end{equation}

y:

\begin{equation}
v_{q,\mathrm{dec}}
=
\omega_e
(L_di_d+\lambda_m).
\end{equation}

Estos términos pertenecen al controlador.

Su existencia no modifica la ecuación física de la planta.

Esta separación es fundamental para que la misma planta pueda recibir diferentes estrategias de control.

\subsection{Límites del FOC}

El FOC reconstruido también incorpora límites de corriente y tensión.

El límite histórico de corriente $q$ utilizado en esta configuración es aproximadamente:

\begin{equation}
|i_q|\leq112.43\,\mathrm{A}.
\end{equation}

La tensión máxima se relaciona con el bus mediante:

\begin{equation}
V_{\max}
=
0.95
\frac{V_{DC}}{\sqrt{3}}.
\end{equation}

Estos límites pertenecen a la estrategia de accionamiento y control.

No deben interpretarse como sustitutos de los límites físicos documentados para la máquina.

\subsection{Arquitectura del experimento B0}

La primera referencia cerrada de la reconstrucción Python utiliza una cadena M3--M6--M4.

La arquitectura es:

\begin{equation}
\omega^\ast
\rightarrow
\mathrm{FOC}
\rightarrow
(v_d^\ast,v_q^\ast)
\rightarrow
\mathrm{M3}
\rightarrow
(v_a,v_b,v_c)
\rightarrow
\mathrm{M4}.
\end{equation}

La realimentación devuelve:

\begin{equation}
\mathrm{M4}
\rightarrow
(i_a,i_b,i_c,\omega_m,\theta_m)
\rightarrow
\mathrm{FOC}.
\end{equation}

Esta configuración constituye un experimento de referencia de la planta reconstruida.

\subsection{Condición de bus del experimento B0}

El experimento B0 utiliza:

\begin{equation}
V_{DC}=540\,\mathrm{V}.
\end{equation}

Este valor forma parte de la configuración experimental del accionamiento.

No se incorpora a \texttt{PMSMParameters} porque no constituye un parámetro intrínseco de la máquina.

La misma planta puede utilizarse posteriormente con otra condición de bus si se mantienen documentadas las nuevas condiciones.

\subsection{Referencia de velocidad de B0}

La referencia de velocidad del experimento B0 es:

\begin{equation}
\omega^\ast=157.079633\,\mathrm{rad/s}.
\end{equation}

La carga cambia de cero a:

\begin{equation}
T_{\mathrm{load}}=23\,\mathrm{N\,m}
\end{equation}

a los $0.05\,\mathrm{s}$.

El experimento tiene una duración total de:

\begin{equation}
t_{\mathrm{stop}}=20\,\mathrm{s}.
\end{equation}

Estas magnitudes constituyen la configuración del experimento y no deben confundirse con parámetros permanentes de la planta.

\subsection{Resultado de referencia B0}

La ejecución instrumentada del experimento B0 produjo, al final de la simulación, aproximadamente:

\begin{equation}
\omega_m=157.023729\,\mathrm{rad/s},
\end{equation}

con:

\begin{equation}
i_d=-1.122894\,\mathrm{A},
\end{equation}

\begin{equation}
i_q=35.743789\,\mathrm{A},
\end{equation}

y:

\begin{equation}
T_e=23.000056\,\mathrm{N\,m}.
\end{equation}

La referencia de corriente $q$ al final fue aproximadamente:

\begin{equation}
i_q^\ast=35.761422\,\mathrm{A}.
\end{equation}

Estos valores pertenecen a la configuración experimental B0 y no constituyen nuevos parámetros de la máquina.

\subsection{Métricas instrumentadas}

Entre las métricas obtenidas en B0 se encuentran:

\begin{equation}
e_{\omega,\mathrm{final}}
=
0.055903785\,\mathrm{rad/s},
\end{equation}

y un error relativo final aproximado de:

\begin{equation}
e_{\omega,\mathrm{rel}}
=
0.035589454\%.
\end{equation}

El tiempo de establecimiento al 1\% registrado fue aproximadamente:

\begin{equation}
t_s=6.12635\,\mathrm{s}.
\end{equation}

También se registraron:

\begin{equation}
|i_q|_{\max}=120.2163\,\mathrm{A},
\end{equation}

\begin{equation}
|i_d|_{\max}=3.39164\,\mathrm{A},
\end{equation}

y:

\begin{equation}
|V_{dq}|_{\max}=296.1843\,\mathrm{V}.
\end{equation}

La salida PWM permaneció dentro del intervalo:

\begin{equation}
0\leq D\leq1.
\end{equation}

Estas métricas constituyen resultados del experimento y no valores de parametrización física.

\subsection{Trazabilidad de los resultados B0}

La trayectoria experimental y las métricas fueron almacenadas con identificadores SHA-256.

La trayectoria B0 presenta el hash:

\begin{equation}
\texttt{68c89e3d88ed51d18dd9fe07061d1185169aa95c7dd6a43846d1ba173b914278}.
\end{equation}

El archivo de métricas presenta:

\begin{equation}
\texttt{c0f4fc2e56430b912089660173a18db8808bd6a42512e1691c977a68b1a002ea}.
\end{equation}

Estos identificadores permiten asociar un conjunto de resultados con una versión concreta de los datos generados.

El hash no demuestra por sí mismo la validez física del resultado. Su función es garantizar la identidad del artefacto almacenado.

\subsection{Reproducibilidad experimental}

La reproducibilidad requiere más que conservar un archivo de resultados.

Debe ser posible reconstruir la configuración utilizada mediante:

\begin{equation}
C=
\{
\Theta,
V_{DC},
\omega^\ast,
T_{\mathrm{load}},
x_0,
\Delta t,
t_{\mathrm{stop}},
\text{control}
\}.
\end{equation}

La ejecución de un experimento debe registrar estos elementos junto con sus resultados.

De esta forma, un resultado puede ser asociado a:

\begin{equation}
\text{código}
+
\text{parámetros}
+
\text{configuración}
+
\text{datos}.
\end{equation}

Esta asociación constituye una de las principales ventajas de la reconstrucción reproducible.

\subsection{Resultados y generación automática}

Los resultados experimentales no deben incorporarse manualmente al código fuente del reporte.

El flujo previsto es:

\begin{equation}
\text{experimento Python}
\rightarrow
\text{datos}
\rightarrow
\text{métricas}
\rightarrow
\text{generador}
\rightarrow
\text{LaTeX}.
\end{equation}

El reporte puede entonces incorporar valores y figuras generados automáticamente.

Esta metodología reduce el riesgo de copiar manualmente resultados y permite repetir el proceso cuando cambien las condiciones experimentales.

\subsection{Relación con el generador de datos LaTeX}

El directorio del reporte contiene un generador destinado a producir archivos consumibles por LaTeX.

El archivo:

\begin{equation}
\texttt{report/scripts/generate\_latex\_data.py}
\end{equation}

constituye la infraestructura para esta transferencia.

El archivo generado:

\begin{equation}
\texttt{report/datos/resultados\_pmsm.tex}
\end{equation}

debe contener macros derivadas de resultados experimentales.

La plantilla actual no debe interpretarse como una indicación de que todas las métricas del proyecto ya estén automatizadas.

La infraestructura está preparada para incorporar el experimento de referencia conforme se consoliden sus resultados.

\subsection{Reproducibilidad y control de versiones}

La implementación Python forma parte del repositorio Git del proyecto.

Esto permite asociar:

\begin{equation}
\text{código}
\leftrightarrow
\text{commit}
\leftrightarrow
\text{experimento}
\leftrightarrow
\text{resultado}.
\end{equation}

La identificación mediante Git complementa los hashes de los archivos experimentales.

Una ejecución reproducible debe conservar suficiente información para determinar qué versión del código produjo un resultado determinado.

\subsection{Separación entre implementación y evidencia histórica}

La reconstrucción no incorpora automáticamente todos los scripts encontrados durante la recuperación.

Los scripts diagnósticos antiguos, pruebas incompletas y artefactos históricos se consideran material de trazabilidad salvo que hayan sido seleccionados explícitamente como parte de la implementación activa.

Esta regla evita contaminar el modelo actual con decisiones pertenecientes a experimentos anteriores.

El mismo criterio se aplica al material ROM/FEM histórico.

\subsection{Relación con los artefactos recuperados}

El proceso de recuperación proporcionó diferentes fuentes de información.

Estas pueden clasificarse como:

\begin{itemize}
    \item documentación de fabricante;
    \item implementación MATLAB/Simscape;
    \item archivos históricos;
    \item parámetros M4;
    \item código Python reconstruido;
    \item resultados experimentales.
\end{itemize}

Cada fuente tiene una función diferente.

La reconstrucción Python no intenta resolver automáticamente las contradicciones entre fuentes.

Cuando existen discrepancias, estas deben conservarse y clasificarse hasta que exista evidencia suficiente para resolverlas.

\subsection{Planta reproducible como semilla}

El modelo M4 constituye una semilla computacional porque permite ejecutar experimentos sin depender de un modelo gráfico externo.

La semilla contiene:

\begin{equation}
\{
\text{parámetros},
\text{ecuaciones},
\text{estados},
\text{entradas},
\text{salidas},
\text{integrador}
\}.
\end{equation}

A partir de esta base pueden construirse experimentos más complejos sin modificar necesariamente la definición fundamental de la planta.

Esto permite utilizar la misma planta como entorno común de evaluación.

\subsection{Preparación para control avanzado}

Una planta explícita y reproducible facilita el desarrollo de diferentes estrategias de control.

La arquitectura prevista es:

\begin{equation}
\text{PMSM}
\rightarrow
\left\{
\begin{array}{c}
\mathrm{FOC}\\
\mathrm{MPC}\\
\mathrm{ANFIS}\\
\mathrm{TinyML}
\end{array}
\right.
\end{equation}

La planta permanece común mientras las estrategias de control cambian.

Esta condición permite que las diferencias observadas entre estrategias puedan asociarse a sus respectivas decisiones de control, siempre que se mantengan constantes las condiciones de la planta y del experimento.

\subsection{FOC como controlador establecido}

En la arquitectura actual, FOC constituye el controlador funcional utilizado para cerrar el primer experimento de referencia.

Esto no significa que sea la única estrategia futura.

Su función actual es proporcionar una primera cadena cerrada que permita verificar la interacción:

\begin{equation}
\text{referencia}
\rightarrow
\text{control}
\rightarrow
\text{planta}
\rightarrow
\text{realimentación}.
\end{equation}

Las estrategias MPC, ANFIS y TinyML se desarrollarán posteriormente sobre la misma base de planta y respetando las interfaces establecidas.

\subsection{Relación con MPC}

MPC requiere una representación de la planta que pueda utilizarse para predecir su evolución bajo diferentes acciones de control.

La reconstrucción Python proporciona una base para esa tarea.

Sin embargo, la existencia del modelo M4 no implica que el MPC ya esté implementado o validado.

La relación actual es preparatoria:

\begin{equation}
\text{M4 reproducible}
\rightarrow
\text{modelo predictivo}
\rightarrow
\text{MPC}.
\end{equation}

El modelo predictivo deberá conservar la trazabilidad de las simplificaciones adoptadas.

\subsection{Relación con ANFIS y neurodifuso}

La arquitectura prevista coloca ANFIS como capa supervisora o adaptativa y no como sustituto de la planta ni del controlador principal.

La separación conceptual es:

\begin{equation}
\text{ANFIS}
\rightarrow
\text{supervisión/adaptación}
\rightarrow
\text{FOC}.
\end{equation}

La planta permanece:

\begin{equation}
\text{PMSM}
\end{equation}

y el FOC continúa siendo el controlador funcional de la primera arquitectura cerrada.

La reconstrucción Python permite preparar los datos necesarios para futuras estrategias neurodifusas sin alterar la definición física de la planta.

\subsection{Relación con TinyML}

La misma filosofía se aplica a las futuras capas TinyML.

Los datos producidos por la planta pueden utilizarse para construir conjuntos de entrenamiento, validación y prueba.

La cadena futura puede representarse como:

\begin{equation}
\text{PMSM}
\rightarrow
\text{datos}
\rightarrow
\text{modelo TinyML}
\rightarrow
\text{supervisión}.
\end{equation}

La generación de estos datos debe conservar las condiciones físicas y experimentales utilizadas para producirlos.

De esta manera, la planta se convierte en una fuente controlada de escenarios para el desarrollo de inteligencia aplicada al control.

\subsection{Contratos de interfaz}

La reconstrucción Python mantiene la separación conceptual entre bloques mediante interfaces.

La planta no debe conocer la implementación interna del controlador.

El controlador tampoco debe depender de detalles internos de integración de la planta.

La comunicación conceptual es:

\begin{equation}
\text{controlador}
\rightarrow
[\text{contrato}]
\rightarrow
\text{planta},
\end{equation}

y:

\begin{equation}
\text{planta}
\rightarrow
[\text{contrato}]
\rightarrow
\text{controlador}.
\end{equation}

Esta separación será necesaria para trasladar posteriormente la arquitectura a Simulink/Simscape.

\subsection{Ventaja de la planta común}

Una planta común permite comparar diferentes estrategias sin reconstruir el sistema físico para cada controlador.

La condición puede expresarse como:

\begin{equation}
PMSM_{\mathrm{common}}
=
PMSM_{\mathrm{FOC}}
=
PMSM_{\mathrm{MPC}}
=
PMSM_{\mathrm{ANFIS}}
=
PMSM_{\mathrm{TinyML}}.
\end{equation}

La igualdad anterior representa la intención arquitectónica de utilizar una misma planta física de referencia, no la afirmación de que todos los modelos futuros tendrán exactamente el mismo nivel de fidelidad.

Conforme la planta evolucione, las diferentes estrategias deberán compararse sobre versiones identificadas de dicha planta.

\subsection{Control de fidelidad}

La evolución de la planta debe estar acompañada por una identificación clara del nivel de fidelidad.

El modelo M4 actual corresponde a una representación concentrada.

Las futuras versiones podrán incorporar:

\begin{itemize}
    \item mapas electromagnéticos;
    \item saturación;
    \item dependencia angular;
    \item armónicos;
    \item pérdidas adicionales;
    \item dependencia térmica;
    \item modelos FEM;
    \item modelos reducidos ROM.
\end{itemize}

Cada aumento de fidelidad debe conservar una relación explícita con la versión anterior.

Esto permite determinar si una mejora en una estrategia de control proviene realmente del controlador o de una modificación simultánea de la planta.

\subsection{Criterio de evolución}

La evolución del modelo no debe realizarse únicamente por complejidad.

Una representación de mayor fidelidad debe incorporar información adicional que pueda justificarse mediante datos o modelos de referencia.

El criterio conceptual es:

\begin{equation}
\text{mayor fidelidad}
\Rightarrow
\text{mayor información física verificable}.
\end{equation}

No:

\begin{equation}
\text{mayor complejidad}
\Rightarrow
\text{mayor validez}.
\end{equation}

Esta distinción es especialmente importante para evitar introducir parámetros o fenómenos no sustentados.

\subsection{Estado actual de la reconstrucción}

La implementación actual permite ejecutar una planta PMSM concentrada con parametrización M4 y conectarla a una cadena cerrada FOC.

El modelo actual no debe describirse como:

\begin{itemize}
    \item un modelo FEM;
    \item un modelo no lineal completo;
    \item un conjunto de mapas electromagnéticos validados;
    \item un modelo térmico completo;
    \item un Digital Twin terminado.
\end{itemize}

Su función actual es proporcionar una semilla reproducible y trazable sobre la cual puedan construirse esos niveles posteriores.

\subsection{Valor metodológico de la reconstrucción}

La reconstrucción Python permite transformar una recuperación de archivos en una metodología reproducible.

El conocimiento queda distribuido entre:

\begin{equation}
\text{documentación}
+
\text{ecuaciones}
+
\text{código}
+
\text{experimentos}
+
\text{resultados}.
\end{equation}

Esta estructura reduce la dependencia de un único archivo o herramienta.

También permite que una modificación futura pueda localizarse en términos de:

\begin{equation}
\text{fuente}
\rightarrow
\text{parámetro}
\rightarrow
\text{ecuación}
\rightarrow
\text{código}
\rightarrow
\text{resultado}.
\end{equation}

Esta trazabilidad constituye uno de los fundamentos del Digital Twin desarrollado en este trabajo.

\subsection{Cierre del capítulo}

La reconstrucción reproducible en Python convierte la formulación de la planta PMSM en una implementación explícita y versionable.

La estructura actual separa los parámetros M4, las ecuaciones de la planta, las interfaces de entrada y salida y las capas de control. Esta separación permite utilizar la misma planta como plataforma común para FOC y para las futuras estrategias MPC, ANFIS y TinyML.

La implementación M4 constituye actualmente una representación concentrada y reproducible. Los mapas electromagnéticos, el modelo no lineal, FEM y ROM representan etapas futuras de aumento de fidelidad y no se consideran capacidades ya implementadas.

La relación con MATLAB/Simscape permanece complementaria: MATLAB/Simscape aporta la referencia de implementación y evidencia histórica, mientras que Python proporciona la reconstrucción explícita, reproducible y versionable.

Con esta base queda establecida una semilla computacional sobre la cual puede desarrollarse la validación de la planta y, posteriormente, su evolución hacia el Digital Twin PMSM.

El siguiente capítulo documentará la validación de la planta, estableciendo qué propiedades pueden verificarse actualmente, qué métricas deben utilizarse y cuáles son todavía los límites de evidencia del modelo M4.


\section{Validación de la planta}

\subsection{Propósito de la validación}

La validación de la planta PMSM tiene como objetivo determinar qué propiedades de la representación M4 pueden verificarse con la evidencia disponible y cuáles requieren niveles posteriores de modelado o evidencia experimental.

La validación se plantea de manera progresiva. No se considera que una única simulación de lazo cerrado demuestre por sí misma la validez completa de la planta. En su lugar, se separan las verificaciones matemáticas, numéricas, funcionales y de correspondencia con la documentación de referencia.

El criterio general es:

\begin{equation}
\begin{aligned}
\text{modelo}
&\rightarrow \text{implementación}
\rightarrow \text{prueba reproducible}
\\
&\rightarrow \text{métrica}
\rightarrow \text{evidencia}
\rightarrow \text{alcance de validez}.
\end{aligned}
\end{equation}

Esta separación permite evitar que un resultado correcto bajo una condición particular sea interpretado como una validación general de todos los fenómenos físicos de la máquina.

\subsection{Niveles de validación}

La validación actual puede organizarse en varios niveles complementarios:

\begin{itemize}
    \item consistencia de parámetros;
    \item consistencia matemática;
    \item consistencia de implementación;
    \item comportamiento numérico;
    \item comportamiento de lazo cerrado;
    \item correspondencia con datos documentales;
    \item validación de fenómenos de mayor fidelidad.
\end{itemize}

Los primeros niveles pueden ejecutarse sobre la implementación M4 actual. Los últimos requieren información que todavía no se encuentra incorporada de manera activa en la planta.

\subsection{Consistencia de los parámetros}

La primera verificación consiste en comprobar que los parámetros utilizados por la planta coincidan con la fuente de verdad establecida para M4.

La implementación actual concentra los parámetros en \texttt{parameters.py}, donde cada magnitud posee información de trazabilidad.

Los valores utilizados por la planta M4 incluyen:

\begin{itemize}
    \item $R_s=0.05\,\Omega$;
    \item $L_d=1.2\times10^{-3}\,\mathrm{H}$;
    \item $L_q=1.2\times10^{-3}\,\mathrm{H}$;
    \item $\lambda_m=0.107245\,\mathrm{Wb}$;
    \item $p=4$ pares de polos;
    \item $J=0.00416\,\mathrm{kg\,m^2}$.
\end{itemize}

La presencia de estos valores en la configuración de la planta permite verificar que el experimento utiliza una parametrización explícita y no valores ocultos dentro del algoritmo de control.

\subsection{Consistencia entre documentación y modelo}

La validación documental compara los parámetros empleados por M4 con la caracterización de la máquina de referencia.

Esta comparación no significa que todos los valores deban coincidir numéricamente con una única fuente. Algunos parámetros corresponden directamente a datos del fabricante, mientras que otros son derivados o constituyen decisiones de modelado.

Por ejemplo, la resistencia de referencia publicada por Siemens es de $0.046\,\Omega$ a $20\,^{\circ}\mathrm{C}$, mientras que M4 utiliza $0.05\,\Omega$. La diferencia debe conservarse explícitamente en la trazabilidad y no eliminarse mediante una sustitución silenciosa.

De manera similar, la documentación proporciona una inductancia $L_D=1.2\,\mathrm{mH}$, mientras que M4 utiliza $L_d=L_q=1.2\,\mathrm{mH}$. La igualdad entre ambos ejes constituye una hipótesis del modelo concentrado actual y no una afirmación de que la máquina física carezca de anisotropía en cualquier condición de operación.

\subsection{Verificación del flujo magnético}

El flujo permanente utilizado por M4 puede obtenerse a partir del par específico documentado para la máquina.

La relación empleada es:

\begin{equation}
k_T=
\frac{3}{2}\sqrt{2}\,p\,\lambda_m.
\end{equation}

Por tanto:

\begin{equation}
\lambda_m=
\frac{k_T}
{\left(\frac{3}{2}\right)\sqrt{2}\,p}.
\end{equation}

Con $k_T=0.91\,\mathrm{N\,m/A}$ y $p=4$ se obtiene aproximadamente:

\begin{equation}
\lambda_m\approx0.107245\,\mathrm{Wb}.
\end{equation}

Esta magnitud es la utilizada actualmente por la planta M4.

La documentación proporciona además un valor de $k_E$, que permite construir una referencia independiente. La comparación entre ambas estimaciones no se utiliza para declarar automáticamente un error del modelo, debido a que las condiciones de temperatura, tolerancias y convenciones de magnitud pueden ser diferentes.

\subsection{Verificación de la estructura matemática}

La planta M4 implementa las ecuaciones eléctricas en coordenadas $dq$:

\begin{equation}
v_d=
R_si_d+
L_d\frac{di_d}{dt}
-
\omega_eL_qi_q.
\end{equation}

\begin{equation}
v_q=
R_si_q+
L_q\frac{di_q}{dt}
+
\omega_e
\left(
L_di_d+\lambda_m
\right).
\end{equation}

Las ecuaciones diferenciales implementadas son:

\begin{equation}
\frac{di_d}{dt}
=
\frac{
v_d-R_si_d+\omega_eL_qi_q
}{
L_d
}.
\end{equation}

\begin{equation}
\frac{di_q}{dt}
=
\frac{
v_q-R_si_q-
\omega_e(L_di_d+\lambda_m)
}{
L_q
}.
\end{equation}

La relación entre velocidad eléctrica y mecánica se establece mediante:

\begin{equation}
\omega_e=p\omega_m.
\end{equation}

La dinámica mecánica utiliza:

\begin{equation}
J\frac{d\omega_m}{dt}
=
T_e-T_{\mathrm{load}}.
\end{equation}

Y la posición mecánica:

\begin{equation}
\frac{d\theta_m}{dt}
=
\omega_m.
\end{equation}

La validación matemática consiste en verificar que las expresiones implementadas en \texttt{PMSMBase} sean consistentes con la formulación establecida en los capítulos anteriores.

\subsection{Verificación del par electromagnético}

La expresión de referencia utilizada por M4 es:

\begin{equation}
T_e=
\frac{3}{2}p
\left[
\lambda_mi_q+
(L_d-L_q)i_di_q
\right].
\end{equation}

Dado que el modelo actual utiliza:

\begin{equation}
L_d=L_q,
\end{equation}

la expresión se reduce a:

\begin{equation}
T_e=
\frac{3}{2}p\lambda_mi_q.
\end{equation}

Esta simplificación permite comprobar directamente la relación entre corriente $i_q$ y par electromagnético dentro del modelo concentrado.

No debe interpretarse esta expresión como una ley universal de la máquina física bajo cualquier condición. Cuando se incorporen saturación, anisotropía efectiva, dependencia angular, armónicos o mapas electromagnéticos, la relación puede requerir una representación diferente.

\subsection{Verificación de las condiciones iniciales}

La implementación M4 inicia sus estados principales en cero:

\begin{equation}
x(0)=
\begin{bmatrix}
0&
0&
0&
0
\end{bmatrix}^{T}.
\end{equation}

Por tanto:

\begin{equation}
i_d(0)=i_q(0)=\omega_m(0)=\theta_m(0)=0.
\end{equation}

Esta condición constituye una condición inicial del experimento y no una propiedad permanente de la máquina.

Mantenerla explícita permite reproducir los resultados y evita introducir estados iniciales implícitos.

\subsection{Integración numérica}

La implementación actual utiliza integración explícita de Euler.

Para un estado genérico:

\begin{equation}
x_{k+1}
=
x_k+
\Delta t\,f(x_k,u_k,\Theta).
\end{equation}

El paso utilizado en la reconstrucción B0 es:

\begin{equation}
\Delta t=50\,\mu\mathrm{s}.
\end{equation}

Por tanto, la frecuencia de actualización correspondiente es:

\begin{equation}
f_s=
\frac{1}{\Delta t}
=
20\,\mathrm{kHz}.
\end{equation}

Esta condición permite reproducir exactamente la discretización utilizada por el experimento de referencia B0.

\subsection{Alcance de la validación numérica}

La validación numérica comprueba que la implementación puede ejecutar una trayectoria completa sin producir estados o señales fuera de las condiciones establecidas por el experimento.

Entre las variables de interés se encuentran:

\begin{itemize}
    \item velocidad mecánica;
    \item corriente $i_d$;
    \item corriente $i_q$;
    \item par electromagnético;
    \item tensiones $v_d$ y $v_q$;
    \item magnitud del vector de tensión;
    \item señales PWM;
    \item posición mecánica.
\end{itemize}

La ausencia de errores numéricos durante una simulación no constituye por sí misma una validación física. Constituye una evidencia de consistencia computacional bajo las condiciones ensayadas.

\subsection{Experimento B0}

El experimento B0 constituye la línea base reproducible de la reconstrucción cerrada M3--M6--M4.

La arquitectura es:

\begin{equation}
\omega^\ast
\rightarrow
\mathrm{FOC}
\rightarrow
(v_d^\ast,v_q^\ast)
\rightarrow
\mathrm{M3}
\rightarrow
\mathrm{M4}
\rightarrow
(i_{abc},\omega_m,\theta_m,T_e).
\end{equation}

La realimentación conecta las variables observadas de la planta con el controlador mediante la interfaz definida para el experimento.

Las condiciones principales de B0 son:

\begin{itemize}
    \item $V_{DC}=540\,\mathrm{V}$;
    \item $\omega^\ast=157.079633\,\mathrm{rad/s}$;
    \item $\Delta t=50\,\mu\mathrm{s}$;
    \item duración total de $20\,\mathrm{s}$;
    \item aplicación de una carga de $23\,\mathrm{N\,m}$ a $t=0.05\,\mathrm{s}$.
\end{itemize}

Estas condiciones pertenecen al experimento y no deben confundirse con parámetros intrínsecos de la máquina PMSM.

\subsection{Resultado de velocidad}

Al finalizar B0 se obtiene:

\begin{equation}
\omega_m=
157.023729\,\mathrm{rad/s}.
\end{equation}

Frente a la referencia:

\begin{equation}
\omega^\ast=
157.079633\,\mathrm{rad/s},
\end{equation}

el error final es:

\begin{equation}
e_\omega=
\omega^\ast-\omega_m
=
0.055903785\,\mathrm{rad/s}.
\end{equation}

El error relativo respecto a la referencia es aproximadamente:

\begin{equation}
e_{\omega,\mathrm{rel}}
=
\SI{0.035589454}{\percent}.
\end{equation}

Este resultado caracteriza el comportamiento del experimento B0 bajo sus condiciones específicas. No se utiliza como una medida universal de exactitud física de la máquina.

\subsection{Tiempo de establecimiento}

El tiempo de establecimiento de referencia utilizado para B0, considerando la banda de \SI{1}{\percent}, es:

\begin{equation}
t_s=
6.12635\,\mathrm{s}.
\end{equation}

Esta métrica caracteriza la respuesta del conjunto planta--controlador y, por tanto, no debe atribuirse exclusivamente a la planta M4.

El resultado depende de la configuración del FOC, de los límites de corriente y tensión, de la referencia de velocidad, de la carga y del paso de integración.

\subsection{Corrientes de referencia}

Al finalizar B0 se obtiene aproximadamente:

\begin{equation}
i_d=-1.122894\,\mathrm{A},
\end{equation}

y:

\begin{equation}
i_q=35.743789\,\mathrm{A}.
\end{equation}

La referencia de corriente $i_q$ es:

\begin{equation}
i_q^\ast=
35.761422\,\mathrm{A}.
\end{equation}

La cercanía entre $i_q$ e $i_q^\ast$ constituye una evidencia del seguimiento de corriente dentro de la condición final del experimento.

No obstante, el comportamiento transitorio debe analizarse por separado debido a que el máximo absoluto de $i_q$ observado durante B0 alcanza aproximadamente:

\begin{equation}
|i_q|_{\max}
=
120.2163\,\mathrm{A}.
\end{equation}

Este valor debe interpretarse en el contexto de los límites utilizados por el FOC reconstruido y no como una nueva especificación nominal de la máquina.

\subsection{Par electromagnético}

El par final obtenido en B0 es:

\begin{equation}
T_e=
23.000056\,\mathrm{N\,m}.
\end{equation}

La proximidad entre el par electromagnético final y la carga aplicada de $23\,\mathrm{N\,m}$ es consistente con la condición de régimen estacionario del experimento.

El resultado verifica además la coherencia interna entre el par calculado por la planta y la dinámica mecánica implementada.

No implica que el modelo concentrado reproduzca todos los mecanismos de generación de par presentes en la máquina física.

\subsection{\texorpdfstring{Tensión en el plano $dq$}{Tensión en el plano dq}}

La magnitud máxima del vector de tensión $dq$ observada durante B0 es aproximadamente:

\begin{equation}
|V_{dq}|_{\max}
=
296.1843\,\mathrm{V}.
\end{equation}

La tensión utilizada por el controlador está condicionada por el límite definido en la implementación del FOC:

\begin{equation}
V_{\max}
=
0.95
\frac{V_{DC}}{\sqrt{3}}.
\end{equation}

Para $V_{DC}=540\,\mathrm{V}$, este límite constituye una condición del accionamiento utilizado en el experimento y no un parámetro intrínseco del motor.

\subsection{Validación de las señales PWM}

Durante B0 no se detectaron valores PWM fuera del intervalo:

\begin{equation}
0\leq D_{A,B,C}\leq1.
\end{equation}

El resultado registrado para el número de muestras fuera del intervalo permitido es cero.

Esta comprobación verifica la consistencia de la interfaz entre el controlador y el modelo eléctrico M3 bajo las condiciones de B0.

No demuestra por sí sola la fidelidad de un inversor físico real, debido a que el modelo utilizado es una representación computacional del accionamiento.

\subsection{Coherencia energética y de potencia}

Una etapa posterior de validación deberá incorporar balances energéticos y de potencia.

Para el modelo eléctrico pueden analizarse:

\begin{equation}
P_{\mathrm{el}}
=
v_ai_a+
v_bi_b+
v_ci_c.
\end{equation}

La potencia mecánica puede expresarse como:

\begin{equation}
P_{\mathrm{mec}}
=
T_e\omega_m.
\end{equation}

La diferencia entre potencia eléctrica y mecánica debe interpretarse considerando las pérdidas y los elementos que estén realmente incluidos en el modelo.

En la implementación M4 actual no se dispone todavía de un modelo completo y validado de pérdidas eléctricas, magnéticas, mecánicas y térmicas. Por ello, el balance energético completo se considera una línea de validación futura.

\subsection{Validación de la dinámica mecánica}

La ecuación:

\begin{equation}
J\dot{\omega}_m
=
T_e-T_{\mathrm{load}}
\end{equation}

permite verificar la consistencia entre la diferencia de pares y la aceleración obtenida numéricamente.

En régimen estacionario, cuando la aceleración se aproxima a cero:

\begin{equation}
T_e\approx T_{\mathrm{load}},
\end{equation}

siempre dentro de la resolución numérica y de los elementos mecánicos incluidos en el modelo.

La validación no incluye actualmente un modelo activo de $B$ ni de $T_{\mathrm{fric}}$. Ambos parámetros permanecen pendientes de verificación y no deben incorporarse como valores históricos sin trazabilidad.

\subsection{Correspondencia con la implementación MATLAB/Simscape}

La relación entre Python y MATLAB/Simscape se utiliza como correspondencia metodológica y de trazabilidad.

La reconstrucción Python implementa explícitamente las ecuaciones que han sido documentadas en los capítulos anteriores. MATLAB/Simscape constituye la referencia de implementación física y documental recuperada del proyecto.

La validación no se formula como una afirmación de reproducción interna exacta de Simscape.

En particular, la ecuación de par implementada en M4 se considera una ecuación de referencia del modelo concentrado y no una afirmación sobre la ley interna exacta utilizada por un bloque propietario de Simscape.

La comparación futura deberá realizarse mediante variables observables y condiciones de ensayo equivalentes, en lugar de asumir identidad de implementación.

\subsection{Correspondencia con datos del fabricante}

La validación documental futura deberá comparar el comportamiento calculado con los datos publicados para la máquina de referencia.

Entre los elementos de comparación se encuentran:

\begin{itemize}
    \item velocidad nominal;
    \item par nominal;
    \item corriente nominal;
    \item velocidad mecánica máxima;
    \item par máximo;
    \item corriente máxima;
    \item características $M-n$;
    \item constantes $k_T$ y $k_E$;
    \item condiciones térmicas asociadas a los datos.
\end{itemize}

La comparación debe conservar las condiciones de temperatura, tensión del enlace DC, tipo de alimentación y región de operación indicadas en la documentación.

No debe compararse una trayectoria de simulación bajo una condición determinada con un punto de catálogo obtenido bajo otra condición sin registrar previamente esa diferencia.

\subsection{\texorpdfstring{Validación de la curva $M-n$}{Validación de la curva M-n}}

Las curvas oficiales $M-n$ constituyen una fuente de evidencia especialmente relevante para las siguientes etapas.

La documentación Siemens presenta regiones de operación asociadas a diferentes condiciones del convertidor y del enlace DC.

La futura digitalización de estas curvas deberá conservar:

\begin{itemize}
    \item fuente documental;
    \item página;
    \item condición eléctrica;
    \item condición térmica;
    \item variante documental;
    \item configuración de la máquina;
    \item versión del modelo utilizado para la comparación.
\end{itemize}

La digitalización deberá realizarse a partir de la gráfica documental disponible. No se introducirán puntos inventados para completar regiones que no puedan extraerse con suficiente precisión.

\subsection{Validación de los mapas electromagnéticos}

Los archivos recuperados durante la reconstrucción contienen estructuras relacionadas con mapas de flujo, torque e inductancias.

Sin embargo, su presencia dentro de los artefactos recuperados no constituye automáticamente evidencia suficiente para utilizarlos como mapas validados de la máquina.

Antes de incorporarlos al modelo activo deberá verificarse:

\begin{itemize}
    \item significado físico de cada eje;
    \item unidades;
    \item convenciones de signo;
    \item coordenadas de corriente;
    \item referencia angular;
    \item relación entre las tablas;
    \item procedencia;
    \item correspondencia con la máquina de referencia.
\end{itemize}

Mientras esta verificación no se complete, las tablas permanecen como material recuperado y trazable, pero no como fuente activa del modelo M4.

\subsection{Validación de la dependencia angular}

El modelo M4 actual no incorpora una dependencia angular detallada de las inductancias ni del flujo.

Por tanto, todavía no se ha validado mediante el modelo actual:

\begin{itemize}
    \item variación angular del flujo;
    \item ondulación de torque;
    \item cogging torque;
    \item armónicos espaciales;
    \item efectos de ranurado;
    \item dependencia angular de las inductancias.
\end{itemize}

Estas propiedades pertenecen a niveles superiores de fidelidad y deberán validarse mediante mapas, FEM, ROM o evidencia experimental adecuada.

\subsection{Validación térmica}

La caracterización documental proporciona información térmica de la máquina, pero la planta M4 no incorpora todavía una dinámica térmica completa.

Por ello, los valores térmicos almacenados en la parametrización se consideran datos de caracterización y trazabilidad.

La validación térmica futura deberá incorporar al menos:

\begin{itemize}
    \item pérdidas eléctricas;
    \item pérdidas dependientes de la velocidad;
    \item pérdidas magnéticas cuando corresponda;
    \item temperatura de devanados;
    \item temperatura del imán;
    \item condiciones de refrigeración;
    \item evolución temporal.
\end{itemize}

No se presenta el modelo M4 actual como un modelo térmico validado.

\subsection{Validación de pérdidas}

La planta actual representa de manera explícita la resistencia estatórica, pero no contiene un modelo completo de pérdidas magnéticas, mecánicas y térmicas.

Por tanto, una validación energética completa queda fuera del alcance de la versión M4.

La incorporación posterior de pérdidas deberá realizarse de manera trazable y sin mezclar parámetros históricos de otras reconstrucciones con la parametrización actual.

\subsection{Resultado reproducible y trazabilidad}

Los resultados B0 cuentan con identificadores criptográficos para conservar la relación entre la trayectoria calculada y sus métricas.

La trayectoria registrada posee el hash:

\begin{equation}
H_{\mathrm{traj}}
=
\texttt{68c89e3d88ed51d18dd9fe07061d1185169aa95c7dd6a43846d1ba173b914278}.
\end{equation}

Las métricas asociadas poseen el hash:

\begin{equation}
H_{\mathrm{metrics}}
=
\texttt{c0f4fc2e56430b912089660173a18db8808bd6a42512e1691c977a68b1a002ea}.
\end{equation}

Estos identificadores permiten comprobar posteriormente que los archivos de resultados utilizados para generar las tablas y figuras del reporte corresponden al mismo conjunto de resultados.

El objetivo de esta trazabilidad es evitar que los números presentados en el documento se desvinculen del experimento que los produjo.

\subsection{Separación entre validación de planta y validación de control}

Es necesario distinguir entre una propiedad de la planta y una propiedad del conjunto planta--controlador.

Por ejemplo, el error final de velocidad, el tiempo de establecimiento y el máximo transitorio de corriente dependen de la interacción entre:

\begin{itemize}
    \item planta M4;
    \item modelo eléctrico M3;
    \item FOC;
    \item límites de tensión;
    \item límites de corriente;
    \item referencia;
    \item carga;
    \item paso de integración.
\end{itemize}

Por tanto, B0 valida principalmente la capacidad del conjunto implementado para ejecutar una trayectoria cerrada reproducible.

La validación aislada de la planta requiere además pruebas en las que puedan estudiarse sus ecuaciones y respuestas sin atribuir automáticamente al motor los resultados producidos por el controlador.

\subsection{Discrepancias conocidas}

La validación debe conservar las discrepancias conocidas en lugar de ocultarlas mediante ajustes posteriores.

Entre ellas se encuentran:

\begin{itemize}
    \item $\lambda_m=0.107245\,\mathrm{Wb}$ en la planta M4 frente a $\lambda_m=0.1100\,\mathrm{Wb}$ en la reconstrucción histórica del FOC;
    \item $225\,\mathrm{A}$ como corriente máxima documentada de la máquina frente a los límites históricos utilizados por el FOC reconstruido;
    \item $R_s=0.05\,\Omega$ en M4 frente a $0.046\,\Omega$ documentado a $20\,^{\circ}\mathrm{C}$;
    \item ausencia actual de $B$ y $T_{\mathrm{fric}}$ en la parametrización activa;
    \item igualdad $L_d=L_q$ como hipótesis del modelo concentrado.
\end{itemize}

Estas diferencias forman parte de la trazabilidad del proceso de reconstrucción y constituyen puntos de trabajo para las siguientes iteraciones.

\subsection{Criterio de aceptación del modelo M4}

El modelo M4 puede considerarse apto para su función actual cuando se cumplan simultáneamente las siguientes condiciones:

\begin{itemize}
    \item los parámetros utilizados sean identificables y trazables;
    \item las ecuaciones implementadas correspondan a la formulación documentada;
    \item la integración numérica sea reproducible;
    \item las interfaces de planta estén definidas;
    \item los experimentos puedan repetirse con la misma configuración;
    \item los resultados puedan asociarse a archivos y versiones concretas;
    \item las limitaciones del modelo estén explícitamente documentadas.
\end{itemize}

Este criterio no exige que M4 reproduzca todos los fenómenos de la máquina física. Exige que su alcance de representación esté claramente delimitado.

\subsection{Evolución de la validación}

La validación deberá evolucionar junto con el nivel de fidelidad de la planta.

El proceso previsto es:

\begin{equation}
\mathrm{M4}
\rightarrow
\mathrm{Mapas}
\rightarrow
\mathrm{NoLineal}
\rightarrow
\mathrm{FEM/ROM}
\rightarrow
\mathrm{DigitalTwin}.
\end{equation}

En cada etapa deberán incorporarse nuevas pruebas específicas.

Para los mapas se validarán coordenadas, unidades y correspondencia física.

Para el modelo no lineal se validarán saturación, anisotropía y dependencia angular.

Para FEM y ROM se validarán resultados espaciales, extracción de parámetros y reducción de orden.

Para el Digital Twin se validará la correspondencia entre diferentes niveles de representación y la capacidad de utilizar una misma planta para diferentes estrategias de control.

\subsection{Relación con las estrategias de control}

La validación de la planta constituye una condición previa para comparar de manera significativa FOC, MPC, ANFIS y TinyML.

Una estrategia de control no debe evaluarse simultáneamente sobre versiones diferentes de la planta sin registrar explícitamente ese cambio.

La arquitectura prevista mantiene:

\begin{equation}
\text{Planta común}
+
\mathrm{Contratos}
+
\mathrm{Controlador}
\rightarrow
\mathrm{Experimento}.
\end{equation}

FOC constituye actualmente el controlador funcional de referencia.

MPC representa una estrategia posterior que podrá utilizar la representación de la planta para predicción.

ANFIS se plantea como capa supervisora o adaptativa y no como sustituto del controlador principal.

TinyML corresponde a una etapa posterior de exploración de inteligencia embebida.

\subsection{Límites de la validación actual}

La evidencia disponible permite validar actualmente la consistencia matemática, la ejecución reproducible, la parametrización explícita y el comportamiento del experimento B0 bajo las condiciones definidas.

Todavía no permite afirmar una validación completa de:

\begin{itemize}
    \item saturación magnética;
    \item mapas electromagnéticos físicos validados;
    \item dependencia angular completa;
    \item cogging torque;
    \item armónicos espaciales;
    \item pérdidas completas;
    \item dinámica térmica;
    \item modelo FEM;
    \item ROM derivado y validado;
    \item correspondencia experimental completa en todo el rango de operación.
\end{itemize}

Esta distinción es fundamental para conservar la integridad científica de la reconstrucción.

\subsection{Síntesis de la validación}

La validación actual establece una base reproducible para la planta PMSM M4.

La consistencia de parámetros permite identificar qué valores utiliza realmente la implementación. La formulación matemática permite verificar las ecuaciones. La integración numérica permite reproducir la trayectoria. B0 demuestra que la cadena M3--M6--M4 puede ejecutarse de manera cerrada bajo una configuración definida.

Los resultados obtenidos deben interpretarse dentro del alcance del modelo y del experimento.

La evidencia actual no se extiende automáticamente a fenómenos que M4 todavía no representa. Por ello, la siguiente evolución debe aumentar la fidelidad física sin perder la trazabilidad establecida.

\subsection{Cierre del capítulo}

La validación de la planta PMSM se establece como un proceso progresivo y trazable, no como una única prueba final.

El modelo M4 actual permite validar la consistencia de la parametrización, las ecuaciones implementadas, la integración numérica y la ejecución reproducible del experimento B0.

La correspondencia con MATLAB/Simscape y con la documentación Siemens proporciona referencias adicionales, pero cada comparación debe conservar las condiciones bajo las cuales fue obtenida.

Las limitaciones actuales quedan explícitamente delimitadas. Los mapas electromagnéticos, el modelo no lineal, FEM, ROM, pérdidas completas y dinámica térmica requieren etapas posteriores de desarrollo y validación.

Con esta base, el siguiente capítulo puede abordar los mapas electromagnéticos como una nueva fuente de información física para aumentar progresivamente la fidelidad de la planta.


\section{Mapas electromagnéticos}

\subsection{Propósito del capítulo}

El nivel M4 desarrollado hasta este punto representa la máquina PMSM mediante una parametrización concentrada y relaciones analíticas de flujo y par. Esta representación establece una base computacional reproducible y físicamente trazable.

El siguiente nivel de enriquecimiento consiste en incorporar información electromagnética dependiente de las variables de operación. Para ello se consideran los mapas de flujo, inductancia y torque recuperados durante la reconstrucción del proyecto, así como la información que pueda generarse posteriormente mediante modelos de mayor fidelidad.

El objetivo de este capítulo no es presentar automáticamente estos datos como mapas físicos validados. Su propósito es establecer qué estructuras existen, qué significado físico pueden tener, qué variables deben verificarse y bajo qué criterios podrán incorporarse posteriormente a la planta.

Los mapas constituyen, por tanto, un mecanismo de enriquecimiento de la representación electromagnética. No sustituyen por sí mismos a la planta M4 ni constituyen todavía una validación independiente de la máquina.

La evolución conceptual puede expresarse como:

\begin{equation}
\begin{aligned}
\text{representación concentrada}
&\rightarrow \text{información electromagnética}
\\
&\rightarrow \text{mapas trazables}
\rightarrow \text{representación no lineal}.
\end{aligned}
\end{equation}

La información FEM podrá constituir, a su vez, una fuente independiente de datos de alta fidelidad para generar, contrastar o validar estos mapas.

\subsection{Motivación física}

En una representación concentrada, las relaciones utilizadas por M4 son:

\begin{equation}
\psi_d=L_di_d+\lambda_m
\end{equation}

\begin{equation}
\psi_q=L_qi_q.
\end{equation}

Estas expresiones constituyen una aproximación de referencia en la cual $L_d$, $L_q$ y $\lambda_m$ permanecen constantes dentro del modelo.

Una máquina PMSM real puede presentar, dependiendo de su geometría, material magnético, nivel de corriente y posición del rotor, fenómenos como saturación, anisotropía, dependencia angular del flujo, variación de las inductancias y componentes adicionales del par electromagnético.

Por tanto, un modelo basado exclusivamente en parámetros constantes no representa necesariamente todos los estados electromagnéticos posibles de la máquina. La incorporación de mapas permite ampliar esa representación cuando la información disponible proporciona las variables y condiciones necesarias.

\subsection{Origen de los datos recuperados}

Durante la reconstrucción del proyecto se localizaron artefactos que contienen estructuras relacionadas con magnitudes electromagnéticas. Entre ellas aparecen variables y matrices asociadas a corrientes, flujo, derivadas de flujo, inductancias y torque.

Estas estructuras constituyen evidencia de que existieron datos electromagnéticos dentro de los artefactos recuperados, pero su presencia no determina por sí misma:

\begin{itemize}
\item la procedencia exacta de cada muestra;
\item las unidades originales;
\item la convención de signos;
\item la orientación angular utilizada;
\item la definición exacta de los ejes de corriente;
\item la correspondencia entre dimensiones de las matrices y variables físicas;
\item el procedimiento mediante el cual fueron obtenidos los datos;
\item si los datos proceden de FEM, identificación, interpolación u otra fuente.
\end{itemize}

Por esta razón, los datos recuperados permanecen inicialmente en estado de trazabilidad pendiente.

\subsection{Estructuras electromagnéticas recuperadas}

Entre las estructuras identificadas durante la reconstrucción aparecen nombres asociados a:

\begin{itemize}
\item vectores de corriente;
\item vectores de coordenadas;
\item matrices de flujo;
\item derivadas de flujo respecto a corriente;
\item derivadas de flujo respecto a posición;
\item matrices de torque;
\item matrices asociadas a $L_d$;
\item matrices asociadas a $L_q$;
\item componentes de flujo en coordenadas $dq0$;
\item torque asociado a diferentes condiciones de corriente y posición;
\item coeficientes relacionados con pérdidas.
\end{itemize}

Entre los artefactos recuperados se encuentran estructuras como:

\begin{verbatim}
iVec
xVec
dfluxdiMatrix
dfluxdxMatrix
TorqueMatrix
LdMatrix
LqMatrix
PmMatrix
fluxD_dq0
fluxQ_dq0
fluxA_dq0
Torque_dq0
fluxD_IBX
fluxQ_IBX
fluxA_IBX
Torque_IBX
rotorLossCoeff_k_h_DQ
k_J_DQ
k_ex_DQ
\end{verbatim}

Estos nombres se conservan como trazabilidad de los artefactos recuperados. No se interpretan como una especificación definitiva de la interfaz física hasta completar la verificación semántica correspondiente.

En particular, nombres como \texttt{PmMatrix}, \texttt{xVec} o \texttt{fluxA\_dq0} se conservan inicialmente como identificadores de los artefactos y no se les asigna una interpretación física adicional que no haya sido demostrada.

\subsection{Variables de corriente}

Un mapa electromagnético requiere identificar de manera inequívoca las variables independientes utilizadas para construirlo.

En el caso de un modelo $dq$, las variables fundamentales esperadas son:

\begin{equation}
i_d,\qquad i_q.
\end{equation}

Si además existe dependencia angular, debe incorporarse una coordenada de posición:

\begin{equation}
\theta_e
\end{equation}

o una variable angular equivalente cuya convención pueda determinarse.

Un conjunto de datos que contenga únicamente valores de corriente no permite reconstruir un mapa físico completo si no se conoce el significado de cada eje.

Por ello, antes de utilizar las matrices recuperadas se deberán verificar:

\begin{itemize}
\item rango de $i_d$;
\item rango de $i_q$;
\item unidades;
\item resolución de cada eje;
\item orden de los índices;
\item correspondencia entre cada muestra y su condición angular;
\item convención de corriente utilizada.
\end{itemize}

\subsection{Flujo en el eje d}

En el modelo M4, el flujo en el eje $d$ se representa mediante:

\begin{equation}
\psi_d=L_di_d+\lambda_m.
\end{equation}

En un mapa electromagnético general, esta relación puede escribirse como:

\begin{equation}
\psi_d=f_d(i_d,i_q,\theta_e).
\end{equation}

La diferencia entre ambas representaciones es fundamental.

La primera corresponde a un modelo concentrado de parámetros constantes. La segunda permite representar una dependencia respecto de corriente y posición cuando los datos disponibles tienen esa resolución y trazabilidad.

La transición de una expresión analítica a una función tabulada no debe realizarse únicamente sustituyendo una matriz en el código. Es necesario verificar previamente que la matriz representa la magnitud indicada y que sus ejes son compatibles con la formulación utilizada por la planta.

\subsection{Flujo en el eje q}

De manera análoga, M4 utiliza:

\begin{equation}
\psi_q=L_qi_q.
\end{equation}

Un modelo basado en mapas puede representarse como:

\begin{equation}
\psi_q=f_q(i_d,i_q,\theta_e).
\end{equation}

La dependencia respecto de $i_d$ puede aparecer cuando existe acoplamiento electromagnético y saturación. La dependencia angular puede aparecer cuando la distribución espacial del campo y la geometría de la máquina generan variaciones con la posición del rotor.

Estas dependencias no se activan en la implementación M4 actual.

\subsection{Inductancias diferenciales}

Las estructuras recuperadas incluyen matrices identificadas como $L_d$ y $L_q$. Para utilizarlas correctamente es necesario determinar si representan inductancias constantes, secantes o diferenciales.

Una posibilidad física es definir inductancias diferenciales mediante derivadas de flujo:

\begin{equation}
L_{dd}=
\frac{\partial\psi_d}{\partial i_d},
\end{equation}

\begin{equation}
L_{dq}=
\frac{\partial\psi_d}{\partial i_q},
\end{equation}

\begin{equation}
L_{qd}=
\frac{\partial\psi_q}{\partial i_d},
\end{equation}

\begin{equation}
L_{qq}=
\frac{\partial\psi_q}{\partial i_q}.
\end{equation}

Sin embargo, no debe asumirse que las matrices recuperadas representan exactamente estas derivadas hasta verificar su definición original.

La existencia de nombres como \texttt{dfluxdiMatrix} constituye una pista relevante para la investigación, pero no sustituye la comprobación de unidades, dimensiones y significado.

\subsection{Torque electromagnético}

En M4, el torque de referencia es:

\begin{equation}
T_e=
\frac{3}{2}p
\left[
\lambda_m i_q+
(L_d-L_q)i_di_q
\right].
\end{equation}

Para el caso $L_d=L_q$, esta expresión se reduce a:

\begin{equation}
T_e=
\frac{3}{2}p\lambda_m i_q.
\end{equation}

Un mapa electromagnético puede representar el torque como:

\begin{equation}
T_e=f_T(i_d,i_q,\theta_e).
\end{equation}

Esta representación permite capturar componentes que no aparecen explícitamente en el modelo concentrado, siempre que los datos hayan sido obtenidos y validados bajo condiciones físicas conocidas.

La matriz \texttt{TorqueMatrix} recuperada se conserva como candidata a esta representación, pero permanece fuera del camino activo hasta completar su verificación.

\subsection{Dependencia angular}

La dependencia angular constituye una de las diferencias principales entre el nivel M4 y una representación electromagnética de mayor resolución.

En M4, las expresiones analíticas no contienen una dependencia detallada de $\theta_e$ en las inductancias o en el flujo.

Un mapa general puede expresarse como:

\begin{equation}
\psi_d=f_d(i_d,i_q,\theta_e),
\qquad
\psi_q=f_q(i_d,i_q,\theta_e).
\end{equation}

De igual manera:

\begin{equation}
T_e=f_T(i_d,i_q,\theta_e).
\end{equation}

La presencia de esta tercera coordenada permite representar potencialmente variaciones espaciales, pero requiere una definición inequívoca del cero angular, periodicidad, sentido de giro y sistema de referencia.

\subsection{Relación entre flujo y torque}

Los mapas de flujo y torque no deben tratarse como conjuntos independientes sin comprobar su consistencia.

Si se dispone de:

\begin{equation}
\psi_d=f_d(i_d,i_q,\theta_e)
\end{equation}

y:

\begin{equation}
\psi_q=f_q(i_d,i_q,\theta_e),
\end{equation}

el torque electromagnético deberá mantener una relación consistente con las magnitudes de flujo y corriente utilizadas por la convención $dq$.

La verificación futura deberá comprobar que las unidades, signos y convenciones empleadas por los mapas permiten recuperar un torque físicamente consistente.

\subsection{Ejes y dimensiones}

Un mapa multidimensional solamente puede utilizarse de forma reproducible si sus ejes están documentados.

Para cada matriz deberán identificarse como mínimo:

\begin{itemize}
\item nombre físico de cada eje;
\item unidad;
\item número de muestras;
\item mínimo;
\item máximo;
\item orden de almacenamiento;
\item periodicidad angular, cuando corresponda;
\item orientación positiva;
\item origen angular;
\item correspondencia con la matriz de salida.
\end{itemize}

La dimensión de una matriz no es suficiente para deducir automáticamente su significado físico.

Por ejemplo, una matriz tridimensional puede representar distintas combinaciones de corriente y posición. La interpretación correcta debe derivarse de los vectores de coordenadas y de la evidencia de generación de los datos.

\subsection{Unidades y convenciones}

Las unidades constituyen una parte de la trazabilidad del mapa.

Las magnitudes esperadas incluyen:

\begin{itemize}
\item corriente en A;
\item flujo en Wb;
\item inductancia en H;
\item torque en N,m;
\item posición angular en rad o grados;
\item coeficientes de pérdidas en las unidades correspondientes a su definición.
\end{itemize}

No se realizará ninguna conversión silenciosa de unidades. Toda transformación necesaria para integrar los mapas a la planta deberá quedar documentada y, cuando sea posible, automatizada.

\subsection{Estado actual de los mapas}

El estado actual del proyecto distingue entre tres niveles:

\begin{enumerate}
\item estructuras recuperadas;
\item mapas electromagnéticos semánticamente verificados;
\item mapas electromagnéticos validados físicamente.
\end{enumerate}

Las estructuras recuperadas se encuentran disponibles para investigación y trazabilidad.

Los mapas semánticamente verificados serán aquellos para los cuales se haya establecido el significado de sus ejes, unidades, dimensiones, convenciones y variables.

La validación física requerirá además una referencia independiente adecuada, por ejemplo datos documentales, resultados FEM reproducibles o evidencia experimental.

Por tanto:

\begin{equation}
\text{dato recuperado}
\neq
\text{mapa verificado}
\neq
\text{mapa validado}.
\end{equation}

\subsection{Estado de implementación en Python}

La implementación Python conserva una representación analítica de referencia como camino activo del nivel M4.

La representación de mapas de flujo mantiene el modelo:

\begin{equation}
\psi_d=L_di_d+\lambda_m
\end{equation}

\begin{equation}
\psi_q=L_qi_q.
\end{equation}

Los datos recuperados no se activan automáticamente. La interfaz de mapas conserva el principio de que una estructura recuperada debe pasar primero por una etapa de verificación antes de modificar el comportamiento reproducible de la planta M4.

Esta decisión evita que una matriz recuperada de procedencia incompleta altere silenciosamente los resultados de referencia.

\subsection{Criterios para activar un mapa}

Antes de activar un mapa electromagnético deberán verificarse como mínimo:

\begin{itemize}
\item procedencia;
\item unidades;
\item dimensiones;
\item ejes;
\item orientación;
\item convenciones $dq$;
\item cobertura del dominio de operación;
\item continuidad o comportamiento en los bordes;
\item coherencia entre flujo, inductancias y torque;
\item comportamiento fuera de la región tabulada;
\item método de interpolación;
\item reproducibilidad de la evaluación.
\end{itemize}

La activación deberá generar además una versión identificable del conjunto de datos para que los resultados puedan asociarse con una fuente concreta.

\subsection{Interpolación y extrapolación}

Una vez verificados los mapas, será necesario definir cómo se evalúan entre puntos tabulados.

Una interpolación puede expresarse de manera general como:

\begin{equation}
\hat{f}(x)=\mathcal{I}\left(f_{\mathrm{map}},x\right),
\end{equation}

donde $\mathcal{I}$ representa el operador de interpolación seleccionado.

La extrapolación fuera del dominio tabulado constituye un problema diferente. No debe suponerse que una interpolación válida dentro del dominio puede extenderse indefinidamente fuera de él.

Por ello, una implementación futura deberá definir explícitamente:

\begin{itemize}
\item límites de cada eje;
\item comportamiento en los bordes;
\item política de extrapolación;
\item saturación de consulta, si se utiliza;
\item advertencias o errores fuera de dominio.
\end{itemize}

\subsection{Consistencia con M4}

La transición hacia mapas debe conservar una referencia con respecto al modelo M4.

En regiones donde los mapas puedan aproximarse por relaciones lineales, sus resultados deberán poder compararse con:

\begin{equation}
\psi_d=L_di_d+\lambda_m
\end{equation}

y:

\begin{equation}
\psi_q=L_qi_q.
\end{equation}

De igual manera, el torque podrá compararse con:

\begin{equation}
T_e=
\frac{3}{2}p
\left[
\lambda_mi_q+
(L_d-L_q)i_di_q
\right].
\end{equation}

La comparación no busca demostrar que el mapa y M4 deban coincidir en todo el dominio. Su función es identificar dónde la representación concentrada deja de reproducir adecuadamente la información electromagnética disponible.

\subsection{Saturación y anisotropía}

Una de las motivaciones principales para incorporar mapas es representar fenómenos que no aparecen en la hipótesis de inductancias constantes.

La saturación puede producir una variación de las derivadas de flujo respecto de las corrientes. La anisotropía puede producir diferencias entre las respuestas de los ejes $d$ y $q$.

Estas propiedades deberán determinarse a partir de los datos, no suponerse solamente por la existencia de una matriz denominada $L_d$ o $L_q$.

La incorporación de mapas permitirá posteriormente construir una representación no lineal cuando las variables y relaciones físicas hayan sido verificadas.

\subsection{Pérdidas asociadas a los artefactos recuperados}

Algunas estructuras recuperadas contienen nombres relacionados con coeficientes de pérdidas, entre ellos:

\begin{verbatim}
rotorLossCoeff_k_h_DQ
k_J_DQ
k_ex_DQ
\end{verbatim}

Estos elementos no se incorporan actualmente al modelo M4.

La existencia de coeficientes recuperados no determina por sí misma:

\begin{itemize}
\item qué mecanismo de pérdida representan;
\item cómo fueron identificados;
\item cuáles son sus unidades;
\item bajo qué velocidad y temperatura fueron obtenidos;
\item si pertenecen a la misma variante de máquina;
\item si fueron derivados de FEM, medición o ajuste.
\end{itemize}

Por ello se mantienen como material de trazabilidad pendiente de interpretación.

\subsection{Relación con FEM}

Los mapas electromagnéticos pueden proceder de modelos FEM, pero no debe inferirse automáticamente que la existencia de una matriz recuperada demuestra que se ejecutó un solver FEM dentro de la reconstrucción actual.

La relación metodológica prevista puede representarse como:

\begin{equation}
\mathrm{FEM}
\rightarrow
\text{extracción de datos}
\rightarrow
\text{mapas}.
\end{equation}

Si los artefactos recuperados proceden efectivamente de una cadena FEM, esa procedencia deberá demostrarse mediante sus archivos de origen, documentación, metadatos o evidencia reproducible.

La planta Python actual no contiene un solver FEM activo.

\subsection{Relación con el modelo no lineal}

Los mapas constituyen una fuente potencial para evolucionar el modelo M4 hacia una representación no lineal.

El modelo futuro podrá reemplazar relaciones constantes por funciones tabuladas:

\begin{equation}
\psi_d=f_d(i_d,i_q,\theta_e)
\end{equation}

\begin{equation}
\psi_q=f_q(i_d,i_q,\theta_e).
\end{equation}

El objetivo será mantener la misma estructura de estados e interfaces de planta cuando sea posible, modificando la representación interna de los fenómenos electromagnéticos.

Esto permitirá conservar la separación entre planta y control y facilitará posteriormente la comparación entre diferentes niveles de fidelidad.

\subsection{Preparación para FEM y ROM}

Los mapas constituyen un puente entre la representación concentrada y las representaciones electromagnéticas de mayor fidelidad.

Una cadena posible es:

\begin{equation}
\mathrm{FEM}
\rightarrow
\text{datos electromagnéticos}
\rightarrow
\text{mapas}
\rightarrow
\text{modelo no lineal y/o ROM}.
\end{equation}

Esta relación no implica que todo mapa deba proceder de FEM ni que todo modelo no lineal deba construirse exclusivamente a partir de mapas FEM. La fuente de cada conjunto de datos deberá mantenerse explícita.

Cada transición deberá conservar la correspondencia entre variables físicas y variables computacionales.

\subsection{Validación futura de los mapas}

La validación de los mapas deberá realizarse progresivamente.

Una primera etapa será la validación semántica:

\begin{itemize}
\item identificación de ejes;
\item unidades;
\item dimensiones;
\item convenciones;
\item origen angular;
\item procedencia.
\end{itemize}

Una segunda etapa será la validación matemática:

\begin{itemize}
\item consistencia entre flujo e inductancias;
\item consistencia del torque;
\item continuidad;
\item simetrías esperadas;
\item periodicidad angular;
\item comportamiento en el dominio de operación.
\end{itemize}

Una tercera etapa será la validación física mediante una referencia independiente.

Esta separación evita que la correcta lectura de una matriz se confunda con la demostración de que representa fielmente la máquina.

\subsection{Relación con la planta común}

La incorporación futura de mapas debe preservar la planta como elemento común de la arquitectura de control.

El objetivo es que FOC, MPC, ANFIS y TinyML puedan utilizar una misma representación física de referencia, con el nivel de fidelidad registrado explícitamente.

La arquitectura conceptual permanece:

\begin{equation}
\text{control}
\rightarrow
\text{contrato}
\rightarrow
\text{accionamiento}
\rightarrow
\text{planta}
\rightarrow
\text{contrato}
\rightarrow
\text{realimentación}.
\end{equation}

La sustitución de M4 por una representación basada en mapas no debe modificar silenciosamente las interfaces del sistema.

\subsection{Trazabilidad de versiones}

Cada conjunto de mapas que llegue a utilizarse experimentalmente deberá identificarse mediante una versión o hash.

Como mínimo deberán conservarse:

\begin{itemize}
\item fuente;
\item fecha de extracción;
\item archivo de origen;
\item variables;
\item unidades;
\item ejes;
\item método de interpolación;
\item versión del modelo;
\item hash del conjunto de datos;
\item estado de validación.
\end{itemize}

De esta manera, un resultado experimental podrá asociarse inequívocamente con la versión de planta utilizada.

\subsection{Estado actual y límites}

El estado actual permite afirmar que existen estructuras recuperadas relacionadas con mapas electromagnéticos y que dichas estructuras pueden servir como material para una futura evolución de la planta.

No permite afirmar todavía que:

\begin{itemize}
\item los mapas estén físicamente validados;
\item todos sus ejes estén semánticamente verificados;
\item sus unidades hayan sido confirmadas;
\item procedan necesariamente de FEM;
\item representen exactamente la unidad Siemens de referencia;
\item puedan activarse sin una etapa adicional de verificación.
\end{itemize}

Por tanto, el camino activo de M4 permanece basado en la representación analítica conservadora, mientras la información electromagnética recuperada constituye una fuente de enriquecimiento en proceso de trazabilidad y verificación.

\subsection{Cierre del capítulo}

Los mapas electromagnéticos representan un mecanismo fundamental para aumentar el contenido físico de la representación PMSM más allá de los parámetros concentrados de M4. Su valor potencial reside en permitir que el modelo evolucione desde parámetros constantes hacia relaciones dependientes de corriente y, cuando exista evidencia suficiente, de posición angular.

La reconstrucción actual conserva estructuras relacionadas con flujo, torque, inductancias y pérdidas. Estas estructuras constituyen material de trazabilidad y una fuente para futuras investigaciones, pero no se presentan todavía como mapas electromagnéticos físicamente validados.

El siguiente paso metodológico consiste en verificar sus ejes, unidades, dimensiones, convenciones y procedencia. Una vez cerrada esa trazabilidad podrá evaluarse su incorporación a un modelo no lineal y su relación con información FEM y modelos reducidos.

Así, los mapas permiten aumentar la fidelidad física sin perder reproducibilidad, trazabilidad ni separación entre planta y control.


\section{Evolución hacia modelo no lineal, FEM y ROM}

\subsection{Propósito del capítulo}

La planta PMSM desarrollada hasta este punto se encuentra en un nivel de representación concentrado, denominado M4. Este nivel proporciona una representación matemática, ejecutable y reproducible de la máquina a partir de un conjunto reducido de parámetros físicamente trazables.

El objetivo de las etapas posteriores no es sustituir arbitrariamente esta representación, sino aumentar progresivamente la cantidad y resolución de los fenómenos físicos representados. La evolución debe conservar al mismo tiempo la reproducibilidad, la trazabilidad de los datos y la separación entre la planta y las capas de control.

La arquitectura de evolución no se interpreta como una cadena estrictamente secuencial. M4 constituye la base activa; los mapas y el modelo no lineal permiten enriquecer la representación de la planta; FEM puede proporcionar una fuente independiente de información electromagnética de mayor fidelidad; y ROM puede utilizar esa información para construir representaciones computacionalmente más eficientes.

De forma conceptual:

\begin{equation}
\mathrm{M4}
\rightarrow
\text{enriquecimiento electromagnético}
\rightarrow
\text{modelos de mayor fidelidad}
\end{equation}

mientras que:

\begin{equation}
\mathrm{FEM}
\rightarrow
\text{datos de alta fidelidad}
\rightarrow
\text{mapas y/o ROM}.
\end{equation}

Estos niveles pueden posteriormente integrarse dentro de la arquitectura del Digital Twin.

\subsection{Principio de aumento progresivo de fidelidad}

Un modelo de mayor complejidad no es automáticamente un modelo de mayor fidelidad.

La fidelidad debe relacionarse con la capacidad de representar fenómenos físicos relevantes mediante información cuya procedencia pueda ser demostrada.

Por ello, cada incremento de fidelidad deberá responder al menos a tres preguntas:

\begin{itemize}
\item qué fenómeno físico adicional se pretende representar;
\item qué información permite representarlo;
\item cómo se verificará que la nueva representación es consistente con una referencia.
\end{itemize}

Este criterio evita incorporar complejidad únicamente por razones de implementación y mantiene una relación explícita entre física, datos y modelo computacional.

\subsection{Nivel M4 como referencia}

El nivel M4 constituye actualmente la referencia activa de la planta.

Su formulación utiliza:

\begin{equation}
\psi_d=L_di_d+\lambda_m
\end{equation}

\begin{equation}
\psi_q=L_qi_q
\end{equation}

y el torque electromagnético de referencia:

\begin{equation}
T_e=
\frac{3}{2}p
\left[
\lambda_mi_q+
(L_d-L_q)i_di_q
\right].
\end{equation}

La dinámica eléctrica se mantiene en forma concentrada:

\begin{equation}
\frac{di_d}{dt} =
\frac{v_d-R_si_d+\omega_eL_qi_q}{L_d}.
\end{equation}

\begin{equation}
\frac{di_q}{dt} =
\frac{v_q-R_si_q-\omega_e(L_di_d+\lambda_m)}{L_q}.
\end{equation}

La dinámica mecánica se representa mediante:

\begin{equation}
J\frac{d\omega_m}{dt} = T_e-T_{\mathrm{load}}.
\end{equation}

Este modelo constituye la referencia computacional frente a la cual podrán compararse las representaciones posteriores.

\subsection{Qué significa aumentar la fidelidad}

El aumento de fidelidad puede incorporar progresivamente:

\begin{itemize}
\item dependencia del flujo respecto de las corrientes;
\item saturación magnética;
\item anisotropía;
\item dependencia angular;
\item variación de las inductancias;
\item componentes adicionales del torque;
\item pérdidas electromagnéticas;
\item comportamiento térmico;
\item efectos geométricos;
\item información procedente de FEM;
\item reducción computacional mediante ROM.
\end{itemize}

No todos estos fenómenos deben incorporarse simultáneamente.

La arquitectura se diseña para permitir que cada fenómeno se incorpore cuando exista información suficiente para justificarlo.

\subsection{Transición desde parámetros constantes}

En M4, las relaciones de flujo pueden escribirse como funciones de las corrientes:

\begin{equation}
\psi_d=f_d(i_d,i_q)
=
L_di_d+\lambda_m
\end{equation}

\begin{equation}
\psi_q=f_q(i_d,i_q)
=
L_qi_q.
\end{equation}

Una representación de mayor fidelidad puede reemplazarlas por:

\begin{equation}
\psi_d=f_d(i_d,i_q,\theta_e)
\end{equation}

\begin{equation}
\psi_q=f_q(i_d,i_q,\theta_e).
\end{equation}

El torque podrá representarse entonces mediante una relación general:

\begin{equation}
T_e=f_T(i_d,i_q,\theta_e).
\end{equation}

La estructura de estados puede mantenerse:

\begin{equation}
x=
\begin{bmatrix}
i_d &
i_q &
\omega_m &
\theta_m
\end{bmatrix}^{T},
\end{equation}

aunque la función dinámica interna se vuelve más rica:

\begin{equation}
\dot{x}=f(x,u,\Theta,\mathcal{M}),
\end{equation}

donde $\mathcal{M}$ representa la información electromagnética adicional.

\subsection{Modelo no lineal}

El modelo no lineal constituye una extensión de la representación M4 en la que las relaciones electromagnéticas pueden depender de las condiciones de operación en lugar de estar restringidas a parámetros constantes.

Su finalidad es representar dependencias que el modelo M4 no captura directamente.

Entre ellas pueden encontrarse:

\begin{itemize}
\item saturación;
\item dependencia cruzada entre $i_d$ e $i_q$;
\item variación de inductancias;
\item dependencia angular;
\item comportamiento no lineal del flujo;
\item torque dependiente del estado electromagnético.
\end{itemize}

La formulación general puede expresarse mediante:

\begin{equation}
\boldsymbol{\psi}
=
\boldsymbol{\psi}(i_d,i_q,\theta_e).
\end{equation}

La matriz de derivadas del flujo puede representarse conceptualmente como:

\begin{equation}
\mathbf{L}_{\mathrm{dif}}
=
\frac{\partial\boldsymbol{\psi}}
{\partial\mathbf{i}},
\end{equation}

con:

\begin{equation}
\mathbf{i}
=
\begin{bmatrix}
i_d\
i_q
\end{bmatrix}.
\end{equation}

Estas relaciones definen la estructura matemática del siguiente nivel de representación. Su incorporación efectiva dependerá de la disponibilidad y verificación de los datos correspondientes.

\subsection{Estado actual del modelo no lineal en Python}

El repositorio contiene una estructura denominada \texttt{pmsm\_nonlinear.py} asociada a la siguiente etapa de evolución.

Actualmente esta estructura funciona como una frontera explícita para el desarrollo futuro y no como un modelo no lineal operativo.

M4 permanece como camino activo de la planta mientras las representaciones posteriores se verifican. Esta separación permite desarrollar la nueva arquitectura sin modificar silenciosamente la referencia reproducible existente.

\subsection{Mapas como entrada del modelo no lineal}

Los mapas descritos en el capítulo anterior constituyen una posible fuente de información para construir el modelo no lineal.

La relación conceptual es:

\begin{equation}
\text{datos electromagnéticos}
\rightarrow
\text{mapas}
\rightarrow
\text{funciones interpoladas}
\rightarrow
\text{modelo no lineal}.
\end{equation}

La incorporación deberá realizarse únicamente después de verificar:

\begin{itemize}
\item ejes;
\item unidades;
\item dimensiones;
\item convenciones;
\item procedencia;
\item dominio de validez;
\item coherencia entre magnitudes.
\end{itemize}

De esta forma, el modelo no lineal no se convierte en un contenedor de datos de procedencia desconocida.

\subsection{Interpolación como parte del modelo}

Cuando una magnitud física se encuentra tabulada, la interpolación pasa a formar parte de la representación computacional.

Puede escribirse:

\begin{equation}
\hat{\psi}_d
=
\mathcal{I}_d
\left(
i_d,i_q,\theta_e
\right)
\end{equation}

y:

\begin{equation}
\hat{\psi}_q
=
\mathcal{I}_q
\left(
i_d,i_q,\theta_e
\right).
\end{equation}

La función de interpolación debe considerarse parte de la configuración del modelo.

Cambiar el método de interpolación puede modificar el resultado de la simulación incluso cuando el conjunto de datos original permanezca idéntico.

Por ello, los experimentos futuros deberán registrar el método utilizado.

\subsection{Dominio de validez}

Todo modelo basado en mapas posee un dominio de validez definido por sus ejes.

Si:

\begin{equation}
i_d\in[i_{d,\min},i_{d,\max}]
\end{equation}

\begin{equation}
i_q\in[i_{q,\min},i_{q,\max}]
\end{equation}

y:

\begin{equation}
\theta_e\in[\theta_{\min},\theta_{\max}],
\end{equation}

la evaluación del modelo debe respetar esos límites.

La extrapolación fuera del dominio deberá considerarse explícitamente.

No se asumirá que una función interpolada reproduce el comportamiento físico fuera de la región para la cual fueron generados sus datos.

\subsection{Validación del modelo no lineal}

La validación del modelo no lineal deberá realizarse contra una referencia independiente.

Las comparaciones podrán incluir:

\begin{itemize}
\item flujo;
\item torque;
\item inductancias;
\item corriente;
\item respuesta frente a posición angular;
\item respuesta frente a velocidad;
\item pérdidas;
\item energía.
\end{itemize}

El criterio de aceptación deberá definirse para cada magnitud y dominio de operación.

Una coincidencia puntual no será suficiente para establecer la validez del modelo en todo el dominio.

\subsection{Introducción del método FEM}

El método de elementos finitos puede utilizarse para obtener una representación electromagnética de mayor resolución espacial.

En una cadena de modelado, FEM puede proporcionar información sobre:

\begin{itemize}
\item distribución de campo magnético;
\item densidad de flujo;
\item enlaces de flujo;
\item torque;
\item inductancias;
\item dependencia angular;
\item pérdidas específicas del modelo utilizado.
\end{itemize}

La salida de FEM puede convertirse posteriormente en datos tabulados para construir mapas o alimentar otras representaciones computacionales.

La relación conceptual es:

\begin{equation}
\text{geometría}
\rightarrow
\text{mallado}
\rightarrow
\mathrm{FEM}
\rightarrow
\text{solución electromagnética}
\rightarrow
\text{datos}.
\end{equation}

La cadena anterior representa una posible fuente de información de alta fidelidad y no implica que todos los mapas recuperados durante la reconstrucción tengan necesariamente ese origen.

\subsection{FEM y la geometría de la máquina}

Una simulación FEM requiere una representación geométrica y material suficientemente definida.

Para la máquina de referencia, la genealogía de geometría debe mantenerse separada de la parametrización M4.

La existencia histórica de archivos de geometría, STEP, materiales o modelos FEM no implica automáticamente que todos ellos correspondan exactamente a la unidad Siemens de referencia.

La correspondencia deberá establecerse mediante trazabilidad documental y geométrica.

Por ello, cualquier modelo FEM futuro deberá registrar:

\begin{itemize}
\item geometría utilizada;
\item materiales;
\item condiciones de frontera;
\item malla;
\item posición angular;
\item condiciones de corriente;
\item condiciones térmicas;
\item solver;
\item versión;
\item archivos de entrada;
\item archivos de salida.
\end{itemize}

\subsection{FEM como referencia de alta fidelidad}

En la arquitectura propuesta, FEM puede convertirse en una referencia de alta fidelidad para evaluar representaciones reducidas o modelos basados en mapas.

Esto no significa que FEM deba ejecutarse durante cada experimento de control.

Por el contrario, el objetivo puede ser utilizar FEM para generar un conjunto de datos representativo y posteriormente construir una representación computacionalmente más eficiente.

La relación puede escribirse como:

\begin{equation}
\mathrm{FEM}
\rightarrow
\text{dataset de alta fidelidad}
\rightarrow
\text{mapas y/o modelo reducido}
\rightarrow
\text{simulación de control}.
\end{equation}

La utilidad de esta cadena dependerá de que el modelo FEM y los datos generados sean suficientemente representativos y verificables.

\subsection{Estado actual de FEM}

La implementación Python actual no contiene un solver FEM activo.

Por tanto, FEM se considera actualmente una etapa de evolución y una posible fuente de datos de alta fidelidad, no una capacidad ya integrada en la planta.

Los artefactos electromagnéticos recuperados descritos anteriormente tampoco demuestran por sí mismos que se haya reconstruido el procedimiento FEM que los originó.

La procedencia deberá establecerse antes de utilizar cualquier dato como referencia FEM.

\subsection{Modelo reducido}

Un modelo reducido, o ROM, busca conservar las características relevantes de una representación de mayor fidelidad reduciendo el costo computacional.

Conceptualmente:

\begin{equation}
\mathcal{M}*{\mathrm{alta\ fidelidad}}
\longrightarrow
\mathcal{M}*{\mathrm{ROM}}.
\end{equation}

Cuando la fuente sea un modelo FEM, la reducción podrá expresarse como:

\begin{equation}
\mathcal{M}*{\mathrm{FEM}}
\longrightarrow
\mathcal{M}*{\mathrm{ROM}}.
\end{equation}

El objetivo no es conservar toda la información disponible, sino identificar una representación suficiente para el propósito de simulación definido.

La reducción deberá evaluarse respecto de una referencia y no solamente respecto de su estabilidad numérica.

\subsection{ROM y control en tiempo de ejecución}

El interés de un ROM dentro del Digital Twin está relacionado con la posibilidad de utilizar una representación de mayor riqueza física durante experimentos de control sin asumir necesariamente el costo computacional de un solver FEM completo.

Esto puede ser especialmente relevante para:

\begin{itemize}
\item simulaciones de larga duración;
\item barridos paramétricos;
\item MPC;
\item optimización;
\item generación de datasets;
\item entrenamiento de modelos inteligentes;
\item evaluación de estrategias híbridas.
\end{itemize}

La reducción debe preservar las variables necesarias para las interfaces de planta.

\subsection{Criterios de un ROM}

Un ROM futuro deberá evaluarse al menos mediante:

\begin{itemize}
\item error respecto a la referencia;
\item dominio de operación;
\item costo computacional;
\item estabilidad numérica;
\item continuidad;
\item conservación de magnitudes físicas relevantes;
\item sensibilidad a las condiciones de operación.
\end{itemize}

El criterio de reducción deberá relacionarse con el objetivo de uso.

Un ROM diseñado para una simulación de control puede tener requisitos diferentes de uno destinado a análisis electromagnético detallado.

\subsection{Conservación de energía y consistencia física}

El aumento de fidelidad no debe evaluarse únicamente mediante diferencias punto a punto.

También será necesario verificar la consistencia física de magnitudes como energía, potencia y torque.

Una relación fundamental es:

\begin{equation}
P_m=T_e\omega_m.
\end{equation}

La potencia eléctrica y la potencia mecánica deberán poder relacionarse mediante las pérdidas representadas por el nivel de modelo utilizado.

A medida que se incorporen pérdidas y fenómenos adicionales, la contabilidad energética se convertirá en una herramienta importante para detectar inconsistencias.

\subsection{Pérdidas}

El modelo M4 actual no incorpora una descripción detallada de pérdidas electromagnéticas.

Las estructuras recuperadas que contienen nombres asociados a coeficientes de pérdidas permanecen pendientes de interpretación.

Una futura extensión podrá distinguir, según la información disponible:

\begin{itemize}
\item pérdidas de cobre;
\item pérdidas en hierro;
\item pérdidas dependientes de frecuencia;
\item pérdidas mecánicas;
\item pérdidas adicionales.
\end{itemize}

Cada mecanismo deberá incorporarse únicamente cuando exista una definición y parametrización trazables.

\subsection{Efectos térmicos}

La temperatura modifica parámetros eléctricos y electromagnéticos.

En el proyecto ya existen datos térmicos almacenados como información de trazabilidad, pero M4 no constituye todavía un modelo térmico completo.

Una evolución futura puede representar:

\begin{equation}
R_s=R_s(T)
\end{equation}

y, cuando exista evidencia suficiente:

\begin{equation}
\lambda_m=\lambda_m(T).
\end{equation}

La extensión térmica deberá conservar las condiciones térmicas bajo las cuales fueron definidos los datos del fabricante.

\subsection{Relación con la máquina Siemens}

La evolución de fidelidad debe permanecer vinculada a la máquina de referencia.

Los datos documentales del fabricante constituyen una referencia externa para parámetros, límites y curvas de operación.

La validación futura podrá comparar el modelo contra regiones documentadas de velocidad, torque y corriente.

La comparación deberá conservar las condiciones correspondientes al convertidor, temperatura y régimen de operación.

\subsection{Curvas par--velocidad como referencia}

Las curvas par--velocidad documentadas por Siemens constituyen una referencia relevante para evaluar el comportamiento global de la máquina.

Estas curvas permiten comparar regiones como:

\begin{itemize}
\item operación continua;
\item operación de sobrecarga;
\item regiones de velocidad elevada;
\item límites de torque;
\item condiciones asociadas al enlace de continua.
\end{itemize}

Antes de digitalizar una curva deberá conservarse la condición correspondiente y la identificación documental.

No se deben generar puntos de una curva únicamente por estimación visual cuando el objetivo sea utilizarla como dato de validación cuantitativa.

Esta referencia podrá integrarse posteriormente como evidencia de trazabilidad física del Digital Twin, complementando las comparaciones internas entre niveles de fidelidad.

\subsection{Relación con el Digital Twin}

El Digital Twin se plantea como una arquitectura evolutiva en la que diferentes representaciones de la planta pueden coexistir o utilizarse según el propósito de simulación, manteniendo una relación trazable con la máquina física.

Por ello, el Digital Twin no debe entenderse como un modelo posterior único al que se llega después de una secuencia obligatoria de representaciones.

La relación conceptual es:

\begin{equation}
\boxed{
\begin{array}{c}
\text{Máquina física}
\\[2mm]
\downarrow
\\
\text{Planta PMSM trazable}
\\[2mm]
\begin{array}{ccc}
\swarrow & & \searrow \\
\text{M4} & & \text{FEM / alta fidelidad} \\
\downarrow & & \downarrow \\
\text{mapas / no lineal} & & \text{mapas / ROM}
\end{array}
\\[2mm]
\downarrow
\\
\text{Arquitectura Digital Twin}
\end{array}
}
\end{equation}

Esta organización permite que M4 continúe siendo la referencia computacional reproducible mientras se desarrollan representaciones de mayor fidelidad.

El término Digital Twin representa, por tanto, una arquitectura y una relación trazable con la planta física, no simplemente la ejecución de un archivo de simulación.

\subsection{Interfaces de planta}

La evolución interna de la planta deberá preservar, cuando sea posible, las interfaces ya establecidas.

Actualmente se dispone de:

\begin{equation}
u_{\mathrm{planta}}
=
[v_a,v_b,v_c,T_{\mathrm{load}},\theta_m]
\end{equation}

y:

\begin{equation}
y_{\mathrm{planta}}
=
[i_a,i_b,i_c,i_d,i_q,\omega_m,\theta_m,T_e].
\end{equation}

El objetivo es que la sustitución de M4 por una representación de mayor fidelidad no obligue a reconstruir toda la arquitectura de control.

Si un nivel posterior requiere nuevas variables, estas deberán introducirse mediante contratos definidos y trazables.

\subsection{Contratos entre niveles de fidelidad}

La arquitectura de contratos también aplica a la evolución de la planta.

Un nivel de fidelidad no deberá acceder de manera implícita a datos internos de otro nivel.

La estructura conceptual es:

\begin{equation}
\mathrm{control}
\rightarrow
\mathrm{contrato}
\rightarrow
\mathrm{planta}
\rightarrow
\mathrm{contrato}
\rightarrow
\mathrm{realimentación}.
\end{equation}

De esta forma, M4, una representación basada en mapas o un ROM pueden ocupar una frontera de planta compatible siempre que implementen el contrato correspondiente.

\subsection{FOC, MPC y capas inteligentes}

La evolución de la planta debe ser independiente de la elección del controlador.

FOC constituye actualmente el controlador funcional de referencia del proyecto.

MPC representa una estrategia de control posterior que podrá utilizar la planta para optimización y predicción.

ANFIS se mantiene como capa supervisora, adaptativa o predictiva y no como sustituto de FOC.

TinyML representa una posible capa inteligente futura.

Una planta con mayor fidelidad deberá poder utilizarse para estudiar estas estrategias sin modificar sus parámetros físicos de forma silenciosa.

\subsection{Validación cruzada entre niveles}

Una estrategia de validación futura consiste en ejecutar condiciones equivalentes sobre diferentes niveles de fidelidad.

Por ejemplo:

\begin{equation}
\mathrm{M4}
\leftrightarrow
\mathrm{mapas}
\leftrightarrow
\mathrm{ROM}
\leftrightarrow
\mathrm{referencia\ FEM}.
\end{equation}

Las diferencias podrán analizarse en función de:

\begin{itemize}
\item corriente;
\item torque;
\item velocidad;
\item flujo;
\item potencia;
\item energía;
\item tiempo de cálculo.
\end{itemize}

Esta comparación permitirá identificar qué fenómenos son relevantes para cada objetivo de simulación.

\subsection{Reproducibilidad}

Cada nivel de fidelidad deberá conservar una configuración reproducible.

Como mínimo deberán registrarse:

\begin{itemize}
\item versión del código;
\item parámetros;
\item dataset;
\item mapas;
\item geometría, cuando corresponda;
\item solver;
\item condiciones iniciales;
\item entradas;
\item paso de integración;
\item método de interpolación;
\item métricas de validación;
\item hash de los artefactos.
\end{itemize}

Esta información permitirá distinguir entre cambios de planta, cambios de controlador y cambios exclusivamente experimentales.

\subsection{Estado actual de la evolución}

El estado actual puede resumirse de la siguiente manera:

\begin{itemize}
\item M4 es la planta activa y reproducible;
\item existen estructuras electromagnéticas recuperadas;
\item los mapas recuperados permanecen pendientes de verificación;
\item el modelo no lineal está preparado como etapa futura, pero no está implementado como camino activo;
\item no existe actualmente un solver FEM activo en Python;
\item ROM constituye una etapa futura de reducción;
\item la arquitectura del Digital Twin se diseña para admitir progresivamente estos diferentes niveles de representación.
\end{itemize}

Esta separación permite continuar el desarrollo sin presentar como terminado aquello que todavía se encuentra en investigación.

\subsection{Criterio de transición entre niveles}

La transición de un nivel al siguiente deberá realizarse cuando el nuevo nivel pueda demostrar una ventaja física o metodológica verificable.

En términos generales:

\begin{equation}
\mathrm{nuevo\ nivel}
\Rightarrow
\mathrm{nueva\ informaci\acute{o}n}
+
\mathrm{verificaci\acute{o}n}
+
\mathrm{reproducibilidad}.
\end{equation}

Agregar código sin agregar información física o capacidad de validación no constituye por sí mismo un incremento de fidelidad.

\subsection{Cierre del capítulo}

La evolución hacia modelos no lineales, FEM y ROM constituye la ruta mediante la cual la planta PMSM puede incorporar progresivamente mayor contenido físico y diferentes niveles de resolución computacional.

M4 permanece como referencia activa porque proporciona una realización reproducible y trazable de la planta sobre la cual pueden compararse las representaciones posteriores.

Los mapas electromagnéticos recuperados representan una fuente potencial para enriquecer la representación de la planta. El modelo no lineal podrá utilizar esta información una vez verificados sus ejes, unidades, procedencia y coherencia física.

FEM podrá utilizarse posteriormente como referencia de mayor fidelidad y como fuente de datos para construir o contrastar mapas y representaciones reducidas. El ROM permitirá estudiar modelos de mayor riqueza física con un costo computacional compatible con experimentos de control, siempre que su error y dominio de validez sean conocidos.

La arquitectura resultante no depende de una única cadena de modelos. Se basa en una planta PMSM físicamente trazada que puede incorporar progresivamente diferentes niveles de representación y que, finalmente, puede integrarse dentro del Digital Twin.

La regla fundamental permanece:

\begin{equation}
\boxed{
\text{mayor fidelidad física}
+
\text{trazabilidad}
+
\text{reproducibilidad}
+
\text{interfaces estables}
}
\end{equation}

De esta manera, la evolución de la planta aumenta su contenido físico sin perder la separación entre planta, accionamiento y control.


\section{Digital Twin PMSM}

Contenido pendiente de desarrollo técnico y validación.


\section{Transformaciones de coordenadas}

Contenido pendiente de desarrollo técnico y validación.


\section{Control FOC}

Contenido pendiente de desarrollo técnico y validación.


\section{Control predictivo MPC}

Contenido pendiente de desarrollo técnico y validación.


\section{ANFIS y control neurodifuso}

Contenido pendiente de desarrollo técnico y validación.


\section{TinyML}

Contenido pendiente de desarrollo técnico y validación.


\section{Arquitectura híbrida y contratos}

Contenido pendiente de desarrollo técnico y validación.


\section{Resultados}

Contenido pendiente de desarrollo técnico y validación.


\section{Discusión}

Contenido pendiente de desarrollo técnico y validación.


\section{Conclusiones}

Contenido pendiente de desarrollo técnico y validación.


\appendix

\section{Trazabilidad documental}

Este apéndice queda deliberadamente como espacio de trazabilidad. Su permanencia en el reporte final se decidirá al cierre del desarrollo.

\begin{longtable}{p{0.29\textwidth}p{0.53\textwidth}}
\toprule
Fuente & Uso \\
\midrule
\texttt{01\_M4\_HARD\_DATA.md} & Parámetros M4 confirmados \\
\texttt{02\_CONTRACTS.md} & Contratos y fronteras de señal \\
\texttt{BASELINE\_MIGRACION\_001.md} & Evidencia y autopsia del modelo histórico M1--M8 \\
\texttt{foc\_reconstructed/} & Reconstrucción funcional de M6 \\
\texttt{m3\_reconstructed/} & Reconstrucción funcional de M3 \\
\texttt{plant/} & Modelo y parametrización PMSM \\
\texttt{integration/} & Experimentos y trazas \\
\texttt{tests/} & Pruebas de consistencia de software \\
\bottomrule
\end{longtable}



\end{document}
